\documentclass[11pt]{article}
\ifdefined\pdfinfoomitdate\pdfinfoomitdate=1\fi
\ifdefined\pdftrailerid\pdftrailerid{}\fi
\usepackage[margin=1.05in]{geometry}
\usepackage{amsmath,amssymb,amsthm,mathtools}
\usepackage{microtype}
\usepackage[T1]{fontenc}
\usepackage{lmodern}
\usepackage{enumitem}
\usepackage[hidelinks,hypertexnames=false]{hyperref}
\usepackage{xcolor}
\usepackage{listings}

\newtheorem{theorem}{Theorem}[section]
\newtheorem{proposition}[theorem]{Proposition}
\newtheorem{lemma}[theorem]{Lemma}
\newtheorem{corollary}[theorem]{Corollary}
\newtheorem{conjecture}[theorem]{Conjecture}
\theoremstyle{definition}
\newtheorem{definition}[theorem]{Definition}
\newtheorem{remark}[theorem]{Remark}
\newtheorem{question}[theorem]{Question}

\newcommand{\Q}{\mathbb Q}
\newcommand{\Z}{\mathbb Z}
\newcommand{\C}{\mathbb C}
\newcommand{\Res}{\operatorname{Res}}
\newcommand{\Disc}{\operatorname{Disc}}
\newcommand{\ord}{\operatorname{ord}}
\newcommand{\vp}{v_p}
\newcommand{\Package}{\mathcal P}
\newcommand{\T}{\mathcal T}
\newcommand{\dd}{\,\mathrm d}

\title{Reciprocal Package Coprimality for the Polynomials \(Q_{a,b}\)}
\author{Johannes Schmitt}
\date{26 June 2026}

\begin{document}
\maketitle


\begin{abstract}
For positive integers \(a,b\), set
\[
 Q_{a,b}(x)=
 \frac{a x^{a+b}-(a+b)x^a+b}{(x^{\gcd(a,b)}-1)^2}.
\]
The two polynomials \(Q_{a,b}\) and \(Q_{b,a}\) are reciprocal partners, so an
unordered pair \(\{a,b\}\) naturally carries the package
\(P_{a,b}=Q_{a,b}Q_{b,a}\).  We prove that these packages are pairwise coprime
over \(\Q\): if \(\{a,b\}\neq\{c,d\}\) and both pairs are unequal, then
\(\gcd(P_{a,b},P_{c,d})=1\) in \(\Q[x]\).

The proof converts a hypothetical common factor into an off-unit-circle
collision among the divided differences \(L_k(x)=(1-x^k)/k\).  Analytic
rigidity, local power sieves, height bounds, and endpoint arithmetic then
reduce every collision to a finite list of exact package states.  The final
step is computer-assisted but deterministic: the bundled verifiers replay the
upper-range elimination, the complete unit-coefficient branch, the
both-nonunit branch, the large one-nonunit branch, and the corrected residual
one-nonunit branch \(D\leq16\,583\), \(2\leq U\leq D\), \(d\leq2601\).  The
only terminal residual pairs are eliminated by modular gcd certificates.
\end{abstract}

\tableofcontents


\section{Statement, context, and proof map}

The central object of the paper is the family
\[
 Q_{a,b}(x)=
 \frac{a x^{a+b}-(a+b)x^a+b}{(x^{\gcd(a,b)}-1)^2}
 \qquad(a,b\in\Z_{>0}).
\]
When \(a\neq b\), write
\[
 P_{a,b}(x)=Q_{a,b}(x)Q_{b,a}(x).
\]
The product depends only on the unordered pair \(\{a,b\}\); it is the
\emph{reciprocal package} attached to that pair.  The denominator in
\(Q_{a,b}\) removes the forced double cyclotomic factor, and
Lemma~\ref{lem:unitcircle} below shows that no root of any \(Q_{a,b}\)
lies on the unit circle.

\begin{theorem}[Main package-coprimality theorem]\label{thm:main}
Let \(a,b,c,d\) be positive integers with \(a\neq b\), \(c\neq d\), and
\(\{a,b\}\neq\{c,d\}\).  Then
\[
 \gcd\bigl(P_{a,b},P_{c,d}\bigr)=1
 \qquad\text{in }\Q[x].
\]
Equivalently, no irreducible factor over \(\Q\) occurs in the reciprocal
packages of two distinct unordered unequal pairs.
\end{theorem}

This statement sits between two nearby problems.  The strongest natural target
would be irreducibility of every nonconstant \(Q_{a,b}\); in the primitive
case this includes Borisov's irreducibility problem for the \(abc\)-polynomials
\cite{Borisov1998}.  The package-coprimality theorem is weaker than universal
irreducibility, but it is uniform in the whole imprimitive family because
\[
 Q_{rA,rB}(x)=rQ_{A,B}(x^r).
\]
It is also strong enough to give each package a factor-theoretic identity that
cannot be shared with any other package.

The proof has five layers.
\begin{enumerate}[label=(\arabic*)]
\item Sections~\ref{sec:dictionary}--\ref{sec:finite-support} translate
package sharing into a graph of equal divided differences and prove the first
orientation facts: no unit-circle roots, no three equal values, disjoint
additive triples, and finite support for any fixed factor orbit.
\item Sections~\ref{sec:tail}--\ref{sec:rational} add analytic and
algebraic rigidity.  The tail-conjugacy identity, sign law, phase bounds,
Schur deformation, and rational-root classification are used later as
structural sieves rather than as a direct contradiction.
\item Sections~12--17 build the arithmetic reduction: prime separation,
\(S\)-unit finiteness, proper-divisibility rigidity, self-power rigidity, local
power sieves, resultant information, sparse Wronskians, and explicit
height--degree bounds.
\item Section~\ref{sec:computation} eliminates the upper range
\(6\,816\,242\leq D\leq175\,394\,637\).  Section~\ref{sec:lowerstates}
then records the lower-range state variables and the exact endpoint conditions
used by the finite searches.
\item Sections~\ref{sec:unitlowercomputation} and
\ref{sec:nonunitcomputation} finish the proof by deterministic
computer-assisted eliminations of, respectively, the unit-coefficient branch
and all nonunit branches.
\end{enumerate}

The computational part is used only after the preceding sections have reduced
a counterexample to explicit integer data.  Each search is conservative: rows
are kept unless a stated mathematical filter removes them, and the verification
scripts recompute the arithmetic conditions and terminal finite-field
certificates from the recorded data.  Thus the trust boundary is the usual one
for a reproducible computer-assisted proof package: exact source code,
deterministic data files, checksums, and independent replayers, not a formal
proof-assistant certificate.

For reference, the following summary lists the main mathematical and
computational outputs in the order in which they are used.  It is an index to
the proof rather than a substitute for the later arguments.

\begin{theorem}[Summary]\label{thm:summary}
\leavevmode
\begin{enumerate}[label=(\roman*)]
\item Every noncyclotomic irreducible reciprocal factor orbit occurs in only
finitely many packages.  Hence the package-sharing graph is locally finite.

\item For every primitive unequal pair \((a,b)\), including boundary pairs,
\[
 Q_{a,b}\mid Q_{c,d}\quad\Longrightarrow\quad(c,d)=(a,b).
\]

\item Let \(A,B\) be coprime and unequal, and suppose \(Q_{A,B}\) is
irreducible.  If \(\beta\) is a root, then
\[
 \beta\notin\Q(\beta)^p
 \qquad\text{for every prime }p.
\]
Every lift \(\{rA,rB\}\) is therefore globally isolated.

\item If \(I(X)\) counts unordered unequal pairs of total at most \(X\)
whose packages are not globally isolated, then
\[
 I(X)=O\!\left(X^{20/11}\log X\right).
\]
Thus a density-one set of package vertices is globally isolated from the
entire infinite family.

\item If an irreducible factor of degree \(e\) divides a primitive
\(Q_{A,B}\), then
\[
 \operatorname{rad}(A+B)\mid e(e+2).
\]
For every prime \(p\mid A+B\), including \(p=2\) with arbitrary
\(v_2(A+B)\), a factor root is not a \(p\)-th power in its own field.

\item If \(p^\kappa\Vert B\) and a degree-\(e\) factor root is a
\(p\)-th power in its own field, then
\[
 e\leq A,\qquad pA\mid e\kappa,
\]
for every prime \(p\), with the reciprocal statement when
\(p^\kappa\Vert A\).  Thus the only local power room at a bad endpoint is
inside its nonunit block, and it requires \(\kappa\geq p\).

\item If \(N=qk\) with \(q\) odd prime and
\(Q_{A,B}=\prod_i h_i\), then the factor degrees satisfy the exact
partition identities of Corollary~\ref{cor:partitionsignature}.  In
particular, if \(q>k+1\), exactly one factor has degree congruent to
\(-2\pmod q\) and there are at most \(k\) factors.  As a consequence,
primitive total \(q\) or \(2q\), with \(q\) an odd prime, always gives
an irreducible \(Q_{A,B}\).

\item Every normalized four-index counterexample satisfies
\[
 \max\{m,n,p,q\}\leq175\,394\,637.
\]
If the maximum is at least \(6\,816\,242\), the primitive minimal polynomial
of a common root is monic with constant coefficient \(\pm1\).

\item An exact integer computation eliminates that entire unit-normalized
range.  Consequently every normalized counterexample satisfies
\[
 \boxed{\max\{m,n,p,q\}\leq6\,816\,241.}
\]
The computation is reduced to \(761\) same-shape rows, all excluded by the
incompatible bounds \(d\leq2rs\leq38\,364\) and \(d\geq48\,688\).

\item A second exact computation eliminates every counterexample in the
remaining box whose primitive minimal polynomial has unit leading and
constant coefficients.  Therefore every counterexample is genuinely
nonunit.

\item The new nonunit computations eliminate the both-nonunit branch, the
one-nonunit branch with \(D\geq16\,584\), and the corrected residual
one-nonunit branch
\[
 D\leq16\,583,\qquad |V|=1,\qquad 2\leq U\leq D,\qquad d\leq2601.
\]
The residual branch produces \(233\,278\) conservative package states, four
terminal state pairs, two package pairs, and eight modular gcd certificates,
all replayed by the bundled verifier.  Consequently no normalized
counterexample remains.

\item If a primitive irreducible polynomial \(h\), of degree \(e\) and
Mahler measure \(M(h)\), divides selected orientations of two primitive
packages, then the radical of the product of their six additive-triple
entries is at most
\[
 2eM(h)^{4-1/e},
\]
and at most \(2eM(h)^{2-1/e}\) in the monic, unit-constant case.  The
leading and constant coefficients also satisfy the four endpoint
divisibilities \((17.24)\), and their prime divisors force the exact slope
matching \((17.25)\)--\((17.26)\).
\end{enumerate}
\end{theorem}

Items (viii)--(xi) are the completion path for Theorem~\ref{thm:main}.  The
height--degree argument and the computation in Section~\ref{sec:computation}
eliminate the upper unit-normalized range; Section~\ref{sec:unitlowercomputation}
eliminates all unit factors in the remaining box; and
Section~\ref{sec:nonunitcomputation} eliminates the both-nonunit branch, the
one-nonunit branch with \(D\geq16\,584\), and the corrected residual
one-nonunit range \(D\leq16\,583\).  Thus no residual branch remains after the
finite arithmetic eliminations.

\section{Reciprocal normalization and the collision dictionary}
\label{sec:dictionary}

The first task is to replace the factor question by an object that can be followed combinatorially.  This section introduces the divided differences \(L_k\), removes the cyclotomic double roots, and records the primitive normalization used throughout the rest of the proof.

For \(0<m<n\), define
\[
 H_{m,n}(x)=mx^n-nx^m+(n-m),
 \qquad
 L_k(x)=\frac{1-x^k}{k}.
\]
Then
\begin{equation}
 H_{m,n}(z)=0
 \quad\Longleftrightarrow\quad
 L_m(z)=L_n(z),
 \tag{2.1}
\end{equation}
and
\begin{equation}
 Q_{a,b}(x)=
 \frac{H_{a,a+b}(x)}{(x^{\gcd(a,b)}-1)^2}.
 \tag{2.2}
\end{equation}

\subsection{Reciprocal packages}

Let \(g=\gcd(a,b)\) and \(N=a+b\).  Direct calculation gives
\begin{equation}
 x^{N-2g}Q_{a,b}(1/x)=Q_{b,a}(x).
 \tag{2.3}
\end{equation}
The monomial in \((2.3)\) is harmless because \(Q_{a,b}(0)=b\neq0\).

\begin{definition}
The reciprocal package of the unordered pair \(\{a,b\}\) is
\[
 \Package_{\{a,b\}}=\{Q_{a,b},Q_{b,a}\},
\]
or, equivalently for factor questions, the product \(Q_{a,b}Q_{b,a}\).
\end{definition}

An edge \(\{m,n\}\), \(0<m<n\), in the collision graph corresponds to the
unordered package \(\{m,n-m\}\).  The two orientations of that package are the
complementary edges
\[
 \{m,n\},\qquad \{n-m,n\}.
 \tag{2.4}
\]

\subsection{The unit circle is completely removed}

The following elementary lemma closes the first logical gap identified in the
audit.

\begin{lemma}[Unit-circle roots]\label{lem:unitcircle}
Let \(e=\gcd(m,n)\).  If \(H_{m,n}(z)=0\) and \(|z|=1\), then
\[
 z^m=z^n=1,
 \qquad\text{hence}\qquad z^e=1.
\]
Conversely, every \(e\)-th root of unity is an exactly double root of
\(H_{m,n}\).  Consequently,
\[
 \frac{H_{m,n}(x)}{(x^e-1)^2}
\]
has no root on the unit circle.
\end{lemma}

\begin{proof}
The collision equation can be written
\[
 z^m=\left(1-\frac mn\right)+\frac mn z^n.
\]
The right side is a strict convex combination of two points of the closed unit
disk.  Its modulus is one, so equality in the triangle inequality forces
\(z^n=1\), and then \(z^m=1\).

If \(\zeta^e=1\), then \(H_{m,n}(\zeta)=H'_{m,n}(\zeta)=0\), while
\[
 H''_{m,n}(\zeta)
 =mn(n-m)\zeta^{m-2}\neq0.
\]
Thus the multiplicity is exactly two.  These are all the unit-circle roots by
the first paragraph.
\end{proof}

These are also all multiple roots.  Indeed,
\[
 H'_{m,n}(x)=mnx^{m-1}(x^{n-m}-1).
\]
If \(H_{m,n}(x)=H'_{m,n}(x)=0\), then \(x\neq0\),
\(x^{n-m}=1\), and therefore
\[
 0=H_{m,n}(x)=(n-m)(1-x^m).
\]
Thus \(x^m=x^n=1\), so \(x^{\gcd(m,n)}=1\).  Consequently every
\(Q_{a,b}\) is square-free and all of its roots lie off the unit circle.
Also
\[
 H_{m,2m}(x)=m(x^m-1)^2,\qquad Q_{m,m}(x)=m,
\]
so an off-unit-circle collision always corresponds to an unequal pair.

\subsection{Primitive reduction}

If \(a=gA\), \(b=gB\), and \(\gcd(A,B)=1\), then
\begin{equation}
 Q_{a,b}(x)=gQ_{A,B}(x^g).
 \tag{2.5}
\end{equation}
Likewise, if a common-root configuration contains the edges \(\{m,n\}\) and
\(\{p,q\}\), and
\[
 s=\gcd(m,n,p,q),\qquad w=z^s,
\]
then
\begin{equation}
 L_{sk}(z)=\frac1sL_k(w).
 \tag{2.6}
\end{equation}
Thus a minimal counterexample may be normalized by
\[
 \gcd(m,n,p,q)=1.
 \tag{2.7}
\]
In that normalization,
\(\gcd(\gcd(m,n),\gcd(p,q))=1\), exactly the coprime-gcd regime occurring in
Schinzel's bivariate theorem.

\section{No three equal values and additive-triple separation}

The dictionary from Section~\ref{sec:dictionary} is useful only if a single root cannot create many unrelated collisions.  The no-three theorem supplies that control: it turns the collision graph at an off-unit-circle point into a matching and forces the additive triples of two hypothetical packages to be disjoint.

For fixed \(z\), let \(\mathcal G_z\) be the graph on the positive integers in
which \(\{m,n\}\) is an edge precisely when \(L_m(z)=L_n(z)\).

\begin{theorem}[No-three theorem]\label{thm:nothree}
If \(|z|\neq1\), there are no three distinct positive integers \(i,j,k\) with
\[
 L_i(z)=L_j(z)=L_k(z).
\]
Consequently, \(\mathcal G_z\) is a matching.
\end{theorem}

\begin{proof}
Let the common value be \(A\), and put \(R=|z|^2\neq1\).  Then
\[
 z^r=1-rA\qquad(r=i,j,k).
\]
Consider the real smooth function
\[
 F(t)=|1-tA|^2-R^t.
\]
It vanishes at the four distinct real points \(0,i,j,k\), whereas
\[
 F'''(t)=-(\log R)^3R^t
\]
never vanishes.  Three applications of Rolle's theorem give a contradiction.
\end{proof}

Because the two complementary edges in \((2.4)\) share their upper endpoint,
Theorem~\ref{thm:nothree} rules out their simultaneous occurrence at an
off-unit-circle root.  We may therefore state the exact equivalence that was too early
in the first attack.

\begin{proposition}[Exact collision reformulation]\label{prop:equivalence}
The reciprocal packages of distinct unordered pairs are pairwise coprime if
and only if
\[
 |E(\mathcal G_z)|\leq1
 \qquad\text{for every }z\in\C^\times,\ \lvert z\rvert\neq1.
 \tag{3.1}
\]
\end{proposition}

\begin{proof}
If two package products share an irreducible factor, choose a complex root
\(z\) of that factor.  Lemma~\ref{lem:unitcircle} gives \(|z|\neq1\), and
an orientation from each package gives two collision edges at \(z\).
Conversely, an edge equation has rational coefficients, so a common root of
two edges is algebraic; because it is off the unit circle, its minimal
polynomial divides the corresponding two \(Q\)-polynomials.  Finally, two
edges representing the same unordered package would be the complementary
edges in \((2.4)\), which share their upper endpoint and are forbidden by
Theorem~\ref{thm:nothree}.  Thus two distinct graph edges are exactly a common
factor between two distinct packages.
\end{proof}

The addition law
\begin{equation}
 (u+v)L_{u+v}(z)=uL_u(z)+vz^uL_v(z)
 \tag{3.2}
\end{equation}
is fundamental.

\begin{lemma}[Propagation]\label{lem:propagation}
If
\[
 L_u=L_{u+v}
 \qquad\text{and}\qquad
 L_v=L_w,
\]
then
\[
 L_u=L_{u+w}.
\]
If \(w\neq v\), this contradicts Theorem~\ref{thm:nothree}.
\end{lemma}

\begin{proof}
The first equality and \((3.2)\) give \(L_u=z^uL_v\).  Hence
\[
 (u+w)L_{u+w}=uL_u+wz^uL_w
 =uL_u+wz^uL_v=(u+w)L_u.
\]
\end{proof}

For an unordered pair define its additive triple
\[
 \T(a,b)=\{a,b,a+b\}.
\]

\begin{corollary}[Disjoint additive triples]\label{cor:triples}
If two distinct packages occur in a hypothetical common-root collision, then
\begin{equation}
 \T(a,b)\cap\T(c,d)=\varnothing.
 \tag{3.3}
\end{equation}
Thus all six integers
\[
 a,b,a+b,c,d,c+d
\]
are pairwise distinct.
\end{corollary}

\begin{proof}
Choose one collision edge from each package.  Every member of an additive
triple is either an endpoint or the gap of its selected edge.  A shared
endpoint contradicts the matching conclusion of Theorem~\ref{thm:nothree}.
A shared gap becomes a shared endpoint after replacing \(z\) by \(1/z\).
For the endpoint--gap case, write the second edge as
\(\{r,r+s\}\) and suppose, for example, that \(s\) is an endpoint of
the first edge \(\{s,w\}\).  Then
\[
 L_r=L_{r+s},\qquad L_s=L_w,
\]
and Lemma~\ref{lem:propagation} gives \(L_r=L_{r+w}\), producing three
distinct equal values.  The other cross cases are symmetric.  These exhaust
the possibilities.
\end{proof}

This simultaneously eliminates equal totals, equal gaps, shared entries, and
all endpoint--gap cross equalities.

\section{Finite support of every factor orbit and local finiteness}
\label{sec:finite-support}

Before looking for a uniform contradiction, we first prove a qualitative finiteness statement.  A fixed noncyclotomic factor orbit can occur in only finitely many packages; this explains why package privacy is, in principle, a finite exact question for each individual package.

For a primitive irreducible polynomial \(h\in\Z[x]\) with \(h(0)\neq0\),
define its normalized reciprocal by
\[
 h^\vee(x)=\operatorname{prim}\!\left(x^{\deg h}h(1/x)\right),
\]
where the sign is chosen so that the leading coefficient is positive.  We
write \([h]=\{h,h^\vee\}\) for its reciprocal orbit.  Since a package
product is reciprocal, occurrence of either member implies occurrence of both
in the product.

The following elementary estimate is independent of the \(p\)-adic theory.
It is the strongest global finiteness statement obtained in this round.

\begin{theorem}[Finite collision support of a fixed point]
\label{thm:pointfinite}
Fix \(z\in\C^\times\) with \(0<r=|z|<1\).  If
\(L_m(z)=L_n(z)\), \(0<m<n\), and \(k=n-m\), then
\begin{equation}
 \frac{k}{m}
 \leq \frac{2r^m}{1-r^m},
 \tag{4.1}
\end{equation}
and
\begin{equation}
 n\leq\frac{2m}{1-r^m}.
 \tag{4.2}
\end{equation}
In particular, \(\mathcal G_z\) has only finitely many edges, and those edges
are effectively bounded in terms of \(r\).
\end{theorem}

\begin{proof}
The addition law and the collision \(L_m=L_{m+k}\) give
\(L_m=z^mL_k\).  Hence
\[
 k(1-z^m)=mz^m(1-z^k).
\]
Taking absolute values yields
\[
 \frac{k}{m}
 =r^m\frac{|1-z^k|}{|1-z^m|}
 \leq r^m\frac{1+r^k}{1-r^m}
 \leq\frac{2r^m}{1-r^m},
\]
which is \((4.1)\).  Since \(k\geq1\), this forces
\(1\leq2mr^m/(1-r^m)\); the right side tends to zero, so only finitely
many \(m\) occur.  The original equality also gives
\[
 n|1-z^m|=m|1-z^n|\leq2m,
\]
and \((4.2)\) follows from \(|1-z^m|\geq1-r^m\).
\end{proof}

For example, one may take
\[
 M(r)=\max\left\{m\geq1:\frac{2mr^m}{1-r^m}\geq1\right\},
\]
with the convention \(M(r)=0\) when the set is empty.  Then every edge has
\(m\leq M(r)\), followed by the explicit bound \((4.2)\).

For a fixed package \(\{c,d\}\), define
\[
 E_{c,d}(X)=\#\bigl\{\{a,b\}:a+b\leq X,
 \ \gcd(P_{a,b},P_{c,d})\neq1\bigr\}.
\]

\begin{theorem}[Finite support of factors; local finiteness]
\label{thm:localfinite}
Every noncyclotomic reciprocal irreducible orbit \([h]\) occurs in only
finitely many package products \(P_{a,b}\).  Consequently every package has
only finitely many sharing partners.  In particular, for each fixed
\(\{c,d\}\),
\[
 E_{c,d}(X)=O_{c,d}(1).
\]
\end{theorem}

\begin{proof}
Choose a member of \([h]\) having a root \(\alpha\) with
\(0<|\alpha|<1\); this is possible because there are no unit-circle roots,
and replacing \(h\) by \(h^\vee\) inverts all roots.  If \([h]\) occurs
in \(P_{a,b}\), reciprocity of the product ensures that \(\alpha\) is a
root of \(P_{a,b}\).  It therefore gives one of the collision edges
\(\{a,a+b\}\) or \(\{b,a+b\}\).  Theorem~\ref{thm:pointfinite} leaves
only finitely many such edges and hence finitely many unordered packages.

A fixed package product has only finitely many irreducible reciprocal orbits.
The union of their finite supports is exactly its set of sharing partners, so
that set is finite.
\end{proof}

\begin{corollary}[Finite decidability of privacy]
\label{cor:privacydecidable}
For every fixed unordered pair \(\{c,d\}\), whether it has a globally
private factor is decidable by a finite exact calculation.  More precisely,
after factoring \(P_{c,d}\), choose for each reciprocal factor orbit an
algebraic root in the open unit disk and a certified rational upper bound
\(\rho<1\) for its modulus.  Replacing \(r\) by \(\rho\) in
Theorem~\ref{thm:pointfinite} gives an explicit finite list of packages, and
exact polynomial gcds decide all remaining comparisons.
\end{corollary}

\begin{remark}
The bounds obtained from a root modulus need not be practical or uniform in
\(c,d\).  The point is logical: target (a), package by package, is no longer
an intrinsically infinite comparison.  A uniform proof that at least one
orbit survives the finite comparison is still missing.
\end{remark}

\section{The divided-difference sign law}

The next several sections derive analytic restrictions on any two-loop configuration.  This one begins with a one-dimensional sign law: once a collision \(L_m=L_n\) is fixed, the graph of \(|1-tA|^2-|z|^{2t}\) crosses zero with a completely controlled sign pattern.

After replacing \(z\) by \(1/z\) and replacing every edge by its complementary
edge, we may assume \(0<|z|<1\).  Put \(R=|z|^2\).

\begin{theorem}[Sign law]\label{thm:sign}
Suppose \(L_m(z)=L_n(z)=A\), where \(0<m<n\) and \(0<|z|<1\).  Then for every
real \(t\notin\{0,m,n\}\),
\[
 \operatorname{sgn}\bigl(|1-tA|^2-R^t\bigr)
 =\operatorname{sgn}\bigl(t(t-m)(t-n)\bigr).
 \tag{5.1}
\]
\end{theorem}

\begin{proof}
Let \(F_A(t)=|1-tA|^2-R^t\).  We have
\(F_A(0)=F_A(m)=F_A(n)=0\), and
\[
 F_A'''(t)=-(\log R)^3R^t>0.
\]
The third divided-difference mean-value theorem gives
\[
 \frac{F_A(t)}{t(t-m)(t-n)}
 =F_A[t,0,m,n]=\frac{F_A'''(\xi)}{3!}>0
\]
for some \(\xi\) in the convex hull of the four nodes.
\end{proof}

In particular, the affine interpolation \(1-tA\) lies strictly inside the
circle \(|w|=|z|^t\) for \(m<t<n\), and strictly outside it for
\(0<t<m\) and \(t>n\).

\section{Exact tail conjugacy and the strict chord obstruction}
\label{sec:tail}

Here the sign law becomes a geometric obstruction.  A first collision conjugates the later tail of the sequence \(L_k(z)\) to a rescaled copy centered at the gap; a second collision must therefore satisfy a strict chord condition.  This is the main analytic mechanism used in the early reductions.

This section contains the main analytic advance of the second attack.

\begin{theorem}[Exact tail-conjugacy identity]\label{thm:tail}
Assume
\[
 L_m(z)=L_{m+a}(z)=A,
 \qquad m,a\geq1.
\]
Then
\begin{equation}
 A=z^mL_a(z),
 \tag{6.1}
\end{equation}
and, for every integer \(j\geq1\),
\begin{equation}
 L_{m+j}(z)-A
 =\frac{j}{m+j}z^m\bigl(L_j(z)-L_a(z)\bigr).
 \tag{6.2}
\end{equation}
The identity extends to every positive real \(j\) after choosing a logarithm
of \(z\) and defining \(z^j\) continuously.
\end{theorem}

\begin{proof}
Apply \((3.2)\) with \((u,v)=(m,a)\):
\[
 (m+a)L_{m+a}=mL_m+az^mL_a.
\]
Since \(L_{m+a}=L_m\), this yields \((6.1)\).  Applying \((3.2)\) with
\((u,v)=(m,j)\), subtracting \((m+j)A\), and using \((6.1)\), gives
\[
 (m+j)(L_{m+j}-A)=jz^m(L_j-L_a),
\]
which is \((6.2)\).
\end{proof}

Thus a loop on the interval \([m,m+a]\) conjugates the entire tail, after an
affine translation and a simple scalar weight, to the original sequence
centered at its gap \(a\).

\begin{corollary}[Shifted one-crossing law]\label{cor:shiftedsign}
Under the hypotheses of Theorem~\ref{thm:tail}, for every real \(j>0\),
\begin{equation}
 \operatorname{sgn}
 \left(|1-jL_a(z)|^2-|z|^{2j}\right)
 =\operatorname{sgn}(j-a).
 \tag{6.3}
\end{equation}
\end{corollary}

\begin{proof}
From \((6.1)\),
\[
 1-(m+j)A=z^m(1-jL_a).
\]
Therefore
\[
 F_A(m+j)=|z|^{2m}
 \left(|1-jL_a|^2-|z|^{2j}\right).
\]
For \(j\neq a\), Theorem~\ref{thm:sign}, with \(n=m+a\), gives
\((6.3)\).  At \(j=a\), both sides are zero directly; here
\(\operatorname{sgn}(0)=0\).
\end{proof}

Now suppose a second loop occurs later.  Since the four endpoints are distinct,
we may choose the smallest endpoint and write the configuration as
\begin{equation}
 L_m=L_{m+a},
 \qquad
 L_{m+s}=L_{m+t},
 \qquad 0<s<t.
 \tag{6.4}
\end{equation}

\begin{corollary}[Strict chord constraint]\label{cor:chord}
Every configuration \((6.4)\) satisfies
\begin{equation}
 \frac{s}{m+s}(L_s-L_a)
 =\frac{t}{m+t}(L_t-L_a),
 \tag{6.5}
\end{equation}
and therefore
\begin{equation}
 L_t=(1-\lambda)L_a+\lambda L_s,
 \qquad
 \lambda=\frac{s(m+t)}{t(m+s)}\in(0,1).
 \tag{6.6}
\end{equation}
Moreover, \(L_s\neq L_a\), so \(L_t\) lies in the open line segment joining
\(L_a\) and \(L_s\).
\end{corollary}

\begin{proof}
Apply \((6.2)\) with \(j=s,t\), then use
\(L_{m+s}=L_{m+t}\).  This gives \((6.5)\), and solving it gives
\((6.6)\).  The inequalities \(0<\lambda<1\) are equivalent to \(0<s<t\).
The four edge endpoints are distinct, so \(s\neq a\).  If
\(L_s=L_a\), Lemma~\ref{lem:propagation} produces
\(L_m=L_{m+s}\); together with \(L_m=L_{m+a}\) this gives three
distinct equal values, contradicting Theorem~\ref{thm:nothree}.
\end{proof}

There is an equivalent block-centroid form.  Put \(b=t-s\).  The two
collisions imply
\begin{equation}
 z^sL_b
 =\frac{m}{m+s}L_a+\frac{s}{m+s}L_s.
 \tag{6.7}
\end{equation}
Thus the average of the geometric block of length \(b\) beginning at \(s\)
lies strictly between the two prefix averages \(L_a\) and \(L_s\), after the
common factor \(1-z\) is removed.

Clearing denominators in \((6.5)\) gives a useful sparse algebraic residue of
the two-loop problem:
\begin{equation}
 a(m+s)z^t-a(m+t)z^s-m(t-s)z^a+(a+m)(t-s)=0.
 \tag{6.8}
\end{equation}
Any surviving counterexample must satisfy both the original trinomial
\(H_{m,m+a}(z)=0\) and the quadrinomial \((6.8)\), as well as the strict
segment condition \((6.6)\) and the shifted sign law \((6.3)\).

\section{Phase localization, a positive cone, and the real-root subproblem}
\label{sec:phase}

The chord condition still leaves complex phases to control.  This section extracts explicit phase bounds from the trinomial equation, rewrites the quadrinomial obstruction as a positive-cone identity, and isolates the special order constraints that occur for negative real roots.

Write a collision as \(L_m=L_{m+a}\), and let
\(z=re^{i\theta}\), \(0<r<1\), with \(\theta\in(-\pi,\pi]\).
The trinomial equation has the exact form
\begin{equation}
 z^m\bigl(m+a-mz^a\bigr)=a.
 \tag{7.1}
\end{equation}
The denominator on the right of
\(z^m=a/(m+a-mz^a)\) lies in a disk centered at the positive real number
\(m+a\), of radius \(mr^a\).

\begin{theorem}[Exact phase bounds]\label{thm:phase}
For every collision \(L_m=L_{m+a}\) with \(0<|z|=r<1\),
\begin{equation}
 \operatorname{dist}(m\theta,2\pi\Z)
 \leq
 \arcsin\!\left(\frac{mr^a}{m+a}\right).
 \tag{7.2}
\end{equation}
If a second collision is written as
\[
 L_{m+s}=L_{m+t},\qquad 0<s<t,
\]
then, with \(b=t-s\),
\begin{align}
 \operatorname{dist}((m+s)\theta,2\pi\Z)
 &\leq
 \arcsin\!\left(\frac{(m+s)r^b}{m+t}\right),
 \tag{7.3}\\
 \operatorname{dist}(s\theta,2\pi\Z)
 &\leq
 \arcsin\!\left(\frac{mr^a}{m+a}\right)
 +\arcsin\!\left(\frac{(m+s)r^b}{m+t}\right).
 \tag{7.4}
\end{align}
If in addition \(m<a\), the first bound is strictly smaller than \(\pi/6\).
\end{theorem}

\begin{proof}
The largest possible argument of a point in a disk centered at a positive real
number \(R\) with radius \(\rho<R\) is
\(\arcsin(\rho/R)\).  Apply this to \((7.1)\).  The second collision gives
\[
 z^{m+s}\bigl(m+t-(m+s)z^{t-s}\bigr)=t-s,
\]
which proves \((7.3)\).  Circular distance to \(2\pi\Z\) is subadditive,
so subtracting the two multiples of \(\theta\) gives \((7.4)\).  If
\(m<a\), then \(mr^a/(m+a)<1/2\).
\end{proof}

\begin{remark}[Two useful normalizations cannot be conflated]
Replacing \(z\) by \(1/z\) swaps the lower endpoint and the gap of every
edge.  Thus one may orient a two-loop configuration so that its smallest
entry among the two lowers and two gaps is a lower endpoint.  This operation,
however, also replaces \(|z|\) by \(|z|^{-1}\).  Without an additional
argument one cannot simultaneously impose that global-minimum normalization
and \(|z|<1\).  The stronger \(\pi/6\) estimate is therefore conditional on
\(m<a\) in the inside orientation, not an automatic global normalization.
\end{remark}

The quadrinomial in \((6.8)\) has a useful decomposition.  Continue to write
\(b=t-s\), and let its left side be \(R_{m,a;s,t}(x)\).

\begin{proposition}[Positive-cone decomposition]\label{prop:positivecone}
One has the exact identity
\begin{equation}
 R_{m,a;s,t}(x)=
 \frac{a(m+t)}{t}H_{s,t}(x)
 +\frac{m(t-s)}{t}H_{a,t}(x).
 \tag{7.5}
\end{equation}
If \(t>a\), then
\[
 \frac{R_{m,a;s,t}(x)}{(x-1)^2}
\]
is a polynomial all of whose coefficients are strictly positive.  At a
surviving common root,
\begin{equation}
 \frac{H_{s,t}(z)}{H_{a,t}(z)}
 =-\frac{m(t-s)}{a(m+t)}\in\mathbb R_{<0}.
 \tag{7.6}
\end{equation}
\end{proposition}

\begin{proof}
Expanding the right side of \((7.5)\) gives \((6.8)\).  For \(0<u<v\),
\begin{equation}
 \frac{H_{u,v}(x)}{(x-1)^2}
 =\sum_{j=0}^{u-1}(v-u)(j+1)x^j
  +\sum_{j=u}^{v-2}u(v-j-1)x^j,
 \tag{7.7}
\end{equation}
so all coefficients are positive.  When \(t>a\), both kernels in
\((7.5)\) are of this form and both scalar coefficients are positive.
At a root of \(R_{m,a;s,t}\), neither kernel value can vanish:
otherwise both would vanish and \(L_a=L_s=L_t\), contrary to the
no-three theorem (the indices are distinct in a genuine configuration).
Dividing \((7.5)\) therefore gives \((7.6)\).
\end{proof}

Thus the disjoint and interlacing interval patterns reduce to a concrete
sector problem: conditional on \(H_{m,m+a}(z)=0\), the two kernel values in
\((7.6)\) would have to lie on opposite rays.  The nested case \(t<a\)
remains algebraically different.

There is a further ordering theorem when the common root is real.

\begin{proposition}[Ordering of two hypothetical negative-real loops]
\label{prop:negativeorder}
Let \(z=-r\), \(0<r<1\).  Every collision edge \(\{m,n\}\) has
\(m\) even and \(n\) odd.  If two distinct collision edges satisfy
\(m<p\), then
\begin{equation}
 m<p,\qquad n<q,\qquad n-m>q-p.
 \tag{7.8}
\end{equation}
In particular, the later edge has larger total and smaller gap; the nested
interval pattern is impossible.
\end{proposition}

\begin{proof}
Put
\[
 E(X)=\frac{1-e^{-X}}{X},\qquad
 O(Y)=\frac{1+e^{-Y}}{Y}.
\]
Both functions are strictly decreasing on \((0,\infty)\).  Values of
\(L_k(-r)\) on even and odd indices are respectively
\(uE(uk)\) and \(uO(uk)\), where \(u=-\log r\).  Equal values cannot have
the same parity; an odd lower and an even upper are also impossible because
\(L_m(-r)>1/m>1/n>L_n(-r)\).  Hence the lower index is even and the upper
one odd.

For each \(X>0\), the equation \(E(X)=O(Y)\) has a unique solution
\(Y>X\).  Write it as \(Y=\phi(X)\), and let the common value be \(c\).
Using
\[
 E'(X)=\frac{1-c(X+1)}{X},\qquad
 O'(Y)=\frac{1-c(Y+1)}{Y},
\]
implicit differentiation gives
\[
 0<\phi'(X)
 =\frac{Y}{X}\frac{c(X+1)-1}{c(Y+1)-1}<1;
\]
the last inequality is equivalent to
\((1-c)Y>(1-c)X\).  Thus \(\phi\) is increasing while
\(\phi(X)-X\) is strictly decreasing.  Applying this at
\(X=um,up\) proves \(n<q\) and \(n-m>q-p\).
\end{proof}

This does not yet exclude two negative-real loops, but it reduces them to a
strictly ordered tradeoff that is invisible in the general complex case.

\section{Schur, Casoratian, and contraction formulations}
\label{sec:schur}

The same collision equation has several equivalent forms.  The Schur, Casoratian, and contraction viewpoints below are included because each exposes a different obstruction; together they clarify why the tail-conjugacy identity is the usable output of this approach.

Define the bivariate deformation
\begin{align}
 \mathcal F_{m,n}(x,t)
 &=(1-t^m)x^n+(t^n-1)x^m+t^m-t^n  \notag\\
 &=\det\!\begin{pmatrix}
 1&t^n&x^n\\
 1&t^m&x^m\\
 1&1&1
 \end{pmatrix}.
 \tag{8.1}
\end{align}
It is divisible by \((t-1)(x-1)(x-t)\).  Put
\[
 \mathcal R_{m,n}(x,t)=
 \frac{\mathcal F_{m,n}(x,t)}{(t-1)(x-1)(x-t)}.
\]
The alternant formula for Schur polynomials gives
\begin{equation}
 \mathcal R_{m,n}(x,t)
 =-s_{(n-2,m-1)}(1,t,x).
 \tag{8.2}
\end{equation}
At \(t=1\),
\begin{equation}
 \mathcal R_{m,n}(x,1)
 =-\frac{H_{m,n}(x)}{(x-1)^2}.
 \tag{8.3}
\end{equation}
If \(e=\gcd(m,n)\) and \([e](x)=(x^e-1)/(x-1)\), the exact relation to our
polynomial is
\begin{equation}
 \mathcal R_{m,n}(x,1)
 =-Q_{m,n-m}(x)[e](x)^2.
 \tag{8.4}
\end{equation}
Thus the specialization is literally \(-Q\) in the primitive case and differs
from it only by an explicit cyclotomic square in general.

For the one-variable specialization put
\begin{equation}
 u_k(x)=h_k(1,1,x)
 =\sum_{j=0}^k(k-j+1)x^j
 =\frac{x^{k+2}-(k+2)x+k+1}{(x-1)^2},
 \tag{8.5}
\end{equation}
with \(u_{-1}=0\).  Jacobi--Trudi gives
\begin{equation}
 \frac{H_{m,n}(x)}{(x-1)^2}
 =u_{n-2}u_{m-1}-u_{n-1}u_{m-2}.
 \tag{8.6}
\end{equation}
The sequence satisfies
\begin{equation}
 u_{k+1}=xu_k+k+2\quad(k\geq-1),
 \qquad
 u_k-2u_{k-1}+u_{k-2}=x^k\quad(k\geq1).
 \tag{8.7}
\end{equation}
If \(v_k=(u_k,u_{k+1})\), then
\begin{equation}
 \det(v_i,v_j)=(j+2)u_i-(i+2)u_j.
 \tag{8.8}
\end{equation}
Hence a collision is a repeated projective value in the sequence \([v_k]\).

Equivalently, let
\[
 r_k(x)=\frac{u_{k-2}(x)}k,
 \qquad r_1=0.
\]
Then
\begin{equation}
 r_{k+1}=\frac{k}{k+1}\bigl(1+xr_k\bigr),
 \tag{8.9}
\end{equation}
and
\begin{equation}
 x\neq1:\qquad
 H_{m,n}(x)=0
 \quad\Longleftrightarrow\quad
 r_m(x)=r_n(x).
 \tag{8.10}
\end{equation}
Indeed,
\begin{equation}
 r_k(x)=-\frac{L_k(x)}{(x-1)^2}-\frac1{x-1}.
 \tag{8.11}
\end{equation}
For \(|x|<1\), every map
\[
 T_k(w)=\frac{k}{k+1}(1+xw)
\]
is a strict affine contraction.  A collision says that the orbit beginning at
\(r_1=0\) hits the unique fixed point of an interval composition.  Formula
\((6.2)\) is the exact affine conjugacy of the tail induced by such a hit.

Pl\"ucker relations by themselves only turn two vanishing minors into a product
identity.  The recurrences \((8.7)\) do not yet yield the desired descent to a
third vanishing minor.  The strict chord identity \((6.6)\) is the strongest
closed consequence obtained from this route in the present round.

\section{The correct second jet of the Schur deformation}
\label{sec:secondjet}

A tempting first-order transversality argument is false, because the relevant Schur curves have the same tangent at the specialization.  The replacement is a second-order computation in symmetric coordinates; the explicit second jet is what survives for later use.

A first-order transversality argument cannot work: all reduced Schur curves
have the same tangent at \(t=1\).  The correct coordinates are
\[
 t=e^\tau,
 \qquad
 x=e^{\tau/2}y.
\]
If $D=|\lambda|$ denotes the size of the partition, homogeneity and symmetry
in the first two variables give
\[
 e^{-D\tau/2}s_\lambda(1,e^\tau,e^{\tau/2}y)
 =s_\lambda(e^{-\tau/2},e^{\tau/2},y),
\]
and the right side is even in \(\tau\).

For our two-row partition define
\[
 S_{m,n}(y,\tau)
 =s_{(n-2,m-1)}(e^{-\tau/2},e^{\tau/2},y).
\]
The bialternant formula yields
\begin{equation}
 S_{m,n}(y,\tau)=
 \frac{
 \sinh(m\tau/2)y^n-
 \sinh(n\tau/2)y^m+
 \sinh((n-m)\tau/2)
 }{
 \sinh(\tau/2)
 \bigl(y^2-2\cosh(\tau/2)y+1\bigr)
 }.
 \tag{9.1}
\end{equation}
At \(\tau=0\), this is \(H_{m,n}(y)/(y-1)^2\).

\begin{proposition}[Second-jet root branch]\label{prop:secondjet}
Suppose $z$ is a root of $H_{m,n}$ with $|z|\neq1$, and put
\[
 A=L_m(z)=L_n(z).
\]
Then the local root branch of \(S_{m,n}(y,\tau)=0\) through \((z,0)\) has the
form
\begin{equation}
 y_{m,n}(\tau)=z+
 \frac z{24}\left(m+n-\frac3A\right)\tau^2+O(\tau^4).
 \tag{9.2}
\end{equation}
\end{proposition}

\begin{proof}
Expanding \((9.1)\) at \(\tau=0\) and using \(H_{m,n}(z)=0\) gives
\begin{equation}
 S_{m,n}(z,\tau)
 =\frac{\tau^2}{24(z-1)^2}
 \bigl(m^3z^n-n^3z^m+(n-m)^3\bigr)
 +O(\tau^4).
 \tag{9.3}
\end{equation}
Since \(z^m=1-mA\) and \(z^n=1-nA\), the bracket equals
\begin{equation}
 mn(n-m)\bigl((m+n)A-3\bigr).
 \tag{9.4}
\end{equation}
On the other hand,
\begin{align}
 \left.\frac{\partial}{\partial y}
 \frac{H_{m,n}(y)}{(y-1)^2}\right|_{y=z}
 &=\frac{mnz^{m-1}(z^{n-m}-1)}{(z-1)^2} \notag\\
 &=-\frac{mn(n-m)A}{z(z-1)^2}.
 \tag{9.5}
\end{align}
Here \(A\neq0\), because \(A=0\) would give \(z^m=1\), contrary to
\(|z|\neq1\).  Hence the derivative is nonzero.  The implicit-function
theorem gives a unique local root branch.  Since
\(S_{m,n}(y,\tau)=S_{m,n}(y,-\tau)\), uniqueness forces
\(y_{m,n}(\tau)=y_{m,n}(-\tau)\); the branch is even.  Combining this with
\((9.3)--(9.5)\) gives \((9.2)\) with remainder \(O(\tau^4)\).
\end{proof}

If two edges \(\{m,n\}\), \(\{p,q\}\) share a root \(z\), with collision
values \(A\), \(B\), their comoving branches therefore separate as
\begin{equation}
 y_{m,n}(\tau)-y_{p,q}(\tau)
 =\frac z{24}
 \left(m+n-p-q-3\left(\frac1A-\frac1B\right)\right)\tau^2
 +O(\tau^4).
 \tag{9.6}
\end{equation}
In the comoving local parameter \(\tau\), the two root branches therefore
differ by \(O(\tau^2)\), and the order is exactly two unless the explicit
bracket in \((9.6)\) vanishes.
For example, if \(m+n=p+q\), higher contact would force \(A=B\), contradicting
the no-three theorem for disjoint edges.  This is the correct second-order
replacement for the incorrect first-order ``transversality'' formulation.

\section{A complete classification of rational roots}
\label{sec:rational}

Rational roots are a boundary case that should not be left implicit.  The positivity of the coefficient expansion makes the classification short, and the result is later used to rule out degree-one complementary factors in finite searches.

We first record the coefficient expansion, valid when \(\gcd(a,b)=1\):
\begin{equation}
 Q_{a,b}(x)=
 b\sum_{j=0}^{a-1}(j+1)x^j
 +a\sum_{j=a}^{a+b-2}(a+b-1-j)x^j.
 \tag{10.1}
\end{equation}
In particular, the constant term is \(b\), the leading coefficient is \(a\),
and all coefficients are positive.

\begin{theorem}[Rational-root classification]\label{thm:rational}
For arbitrary positive \(a,b\), the polynomial \(Q_{a,b}\) has a rational root
if and only if
\[
 (a,b)=(2,1)\quad\text{or}\quad(a,b)=(1,2).
\]
The roots are respectively \(-1/2\) and \(-2\).
\end{theorem}

\begin{proof}
First suppose \(\gcd(a,b)=1\).  By positivity of the coefficients there is no
positive root.  By reciprocity it suffices to classify a rational root
\(x=-r/s\in(-1,0)\), with \(r,s\) coprime positive integers and \(r<s\).
Write $y=r/s\in(0,1)$ and $N=a+b$.  If \(a\) is odd and \(N\) is even, then
\[
 H_{a,N}(-y)=ay^N+Ny^a+b>0.
\]
If \(a\) and \(N\) are both odd, then
\[
 H_{a,N}(-y)=-ay^N+Ny^a+b>0,
\]
because \(y^a>y^N\) and \(N>a\).  If \(a,N\) are both even, then
\[
 H_{a,N}(-y)=b-Ny^a+ay^N>0;
\]
its derivative is \(aN y^{a-1}(y^{N-a}-1)<0\), and it decreases from
\(b\) to \(0\).  Hence only \(a\) even and \(N\) odd can remain; in
particular \(b\) is odd.

The rational-root theorem applied to \((10.1)\) gives
\[
 r\mid b,\qquad s\mid a.
\]
Writing \(N=a+b\), the equation \(H_{a,N}(-r/s)=0\) becomes
\begin{equation}
 b s^{a+b}=(a+b)r^a s^b+a r^{a+b}.
 \tag{10.2}
\end{equation}
Let \(\ell\mid s\), and put \(t=v_\ell(s)\).  Since
\(\ell\nmid b(a+b)r\), the two terms on the right of \((10.2)\) have
valuations \(bt\) and \(v_\ell(a)\), whereas the left side has valuation
\((a+b)t>bt\).  If the two right-hand valuations were unequal, the
valuation of their sum would be their smaller value, strictly below that of
the left side.  They must therefore coincide, giving
\begin{equation}
 v_\ell(a)=b\,v_\ell(s).
 \tag{10.3}
\end{equation}
Likewise, if \(\ell\mid r\) and \(t=v_\ell(r)\), then the two right-hand
valuations are \(at\) and \((a+b)t\), while the left side has valuation
\(v_\ell(b)\).  Hence
\begin{equation}
 v_\ell(b)=a\,v_\ell(r).
 \tag{10.4}
\end{equation}
Consequently,
\begin{equation}
 s^b\mid a,
 \qquad
 r^a\mid b.
 \tag{10.5}
\end{equation}
If \(r\geq2\), then \(b\geq2^a>a\), while \(s\geq2\) and \((10.5)\) give
\(a\geq2^b>b\), a contradiction.  Hence \(r=1\).

Write \(a=s^bA\) with \(\gcd(A,s)=1\).  Dividing \((10.2)\) by \(s^b\)
gives
\begin{equation}
 b(s^a-1)=A(s^b+1)=a+A\leq\frac{3a}{2}.
 \tag{10.6}
\end{equation}
For \(a\geq3\), the left side is at least \(2^a-1>3a/2\).  Therefore
\(a=2\), and \((10.5)\) gives \(s=2\), \(b=1\), \(A=1\).  Thus the only
primitive root in the open unit disk is \(-1/2\) for \(Q_{2,1}\).  Reciprocity
gives \(-2\) for \(Q_{1,2}\).

Finally, for a nonprimitive pair write \(a=gA\), \(b=gB\).  By \((2.5)\), a
rational root \(x\) would make \(x^g\) a rational root of a primitive
\(Q_{A,B}\).  The only possibilities are \(-1/2\) and \(-2\).  If \(g>1\), then \(g\)
must be odd, and the equality \(g\,v_2(x)=\pm1\) is impossible for rational
\(x\).  Hence \(g=1\), and the two cases above are the only cases.
\end{proof}

This theorem closes the degree-one exception in the \(p\)-adic argument below.

\section{Full additive-triple prime separation}

We now move from analytic restrictions to arithmetic support.  The goal is to show that primes appearing in the six entries of two additive triples cannot overlap except in precisely controlled ways.
\label{sec:padic}

We first isolate the two standard local-field facts used below; see, for
example, \cite[Chapters~I--II]{SerreLocalFields}.

\begin{lemma}[Two local-field criteria]\label{lem:localcriteria}
Let \(p\) be a prime.
\begin{enumerate}[label=(\alph*)]
\item A root of unity of order prime to \(p\) belongs to
\(\Q_p^{\mathrm{nr}}\).  If \(\zeta\) is such a root of unity and
\(v_p(\beta-\zeta)\notin\Z\), then \(\Q_p(\beta)/\Q_p\) is ramified.
\item A monic polynomial in \(\Z_p[x]\) with unit constant term and
square-free reduction has only unit roots, and every root lies in a finite
unramified extension of \(\Q_p\).
\end{enumerate}
\end{lemma}

\begin{proof}
For (a), prime-to-\(p\) roots of unity are Teichmuller lifts and lie in the
maximal unramified extension.  If \(\beta\) were also contained in an
unramified extension, then the compositum containing \(\beta\) and \(\zeta\)
would be unramified, whose normalized value group is \(\Z\), contradicting
the nonintegral valuation of their difference.  Part (b) is the usual Hensel
lifting description of finite etale \(\Z_p\)-algebras.  Monicity makes every
root integral, and the unit constant term forces all root valuations to be
zero.
\end{proof}

\begin{lemma}[Good package away from its additive triple]
\label{lem:goodpackage}
Let \(p\) be odd.  If \(p\nmid cd(c+d)\), then every root of both
orientations in \(\Package_{\{c,d\}}\) is a \(p\)-adic unit lying in a
finite unramified extension of \(\Q_p\).
\end{lemma}

\begin{proof}
Write \(c=hC,d=hD\), with \(\gcd(C,D)=1\).  Then
\(p\nmid hCD(C+D)\).  Borisov's discriminant formula
\cite[Lemma~2.1]{Borisov1998} shows that the primitive polynomial
\(Q_{C,D}\), after division by its unit leading coefficient, is monic over
\(\Z_p\), has unit constant term, and has square-free reduction.
Furthermore,
\[
 \frac{d}{dx}Q_{C,D}(x^h)=hx^{h-1}Q'_{C,D}(x^h).
\]
At a root modulo \(p\), one has \(x\neq0\), \(p\nmid h\), and the last
factor is nonzero because the primitive reduction is square-free.  Thus
\(Q_{C,D}(x^h)\) also has square-free reduction.  Equation \((2.5)\) differs
from this composition only by a scalar.  Lemma~\ref{lem:localcriteria}(b)
applies, and the reciprocal orientation is identical.
\end{proof}

For reference, Borisov's Lemmas~2.8--2.11 give the following consequence for
the primitive polynomial \(Q_{A,B}=f_{A+B,A,B}\).

\begin{lemma}[Ramification supplied by a primitive additive-triple prime]
\label{lem:borisovramification}
Let \(\gcd(A,B)=1\), let \(N=A+B\), and suppose that \(p\mid ABN\).
\begin{enumerate}[label=(\alph*)]
\item If \(p\geq5\), every root of either orientation is either a nonunit or
lies at nonintegral-valued distance from a root of unity of order prime to
\(p\).  In the latter case its local field is ramified.
\item If \(p=3\), each fixed orientation has at most one unit root in
\(\Q_3^{\mathrm{nr}}\); every other root is either a nonunit or ramified.
\end{enumerate}
\end{lemma}

\begin{proof}
The three numbers \(A,B,N\) are pairwise coprime, so \(p\) divides exactly
one of them.  If \(p\mid N\), Lemmas~2.8--2.9 of \cite{Borisov1998} place
all roots in clusters around prime-to-\(p\) roots of unity.  At nontrivial
centers the successive distances have valuations
\[
 \frac1p,\ \frac1{p^2-p},\ldots,
\]
and at the center \(1\) the first layer has valuation \(1/(p-2)\), followed
by the same higher layers.  For \(p\geq5\) all are nonintegral.  For
\(p=3\), the first layer at \(1\) consists of exactly one root at integral
valuation; all other layers are fractional.

If \(p\mid B\), Borisov's Lemma~2.10 gives \(A\) nonunit roots and places
the remaining roots in the same root-of-unity clusters.  If \(p\mid A\),
Lemma~2.11 gives the reciprocal statement.  Lemma~\ref{lem:localcriteria}(a)
converts every fractional cluster distance into ramification.  The same
analysis applies after interchanging the two orientations.
\end{proof}

\begin{theorem}[Full additive-triple prime separation]
\label{thm:fullseparation}
Write \(a=gA,b=gB\), with \(\gcd(A,B)=1\).  Suppose
\begin{equation}
 p\mid AB(A+B),\qquad p\nmid cd(c+d).
 \tag{11.1}
\end{equation}
If \(p\geq5\), then
\begin{equation}
 \gcd_{\Q[x]}(P_{a,b},P_{c,d})=1.
 \tag{11.2}
\end{equation}
If \(g=1\), the same conclusion holds for \(p=3\).
\end{theorem}

\begin{proof}
Suppose an irreducible \(h\in\Q[x]\) divides both package products.  Fix an
orientation of each package divisible by \(h\), and choose a root
\(\alpha\in\overline{\Q}_p\) of \(h\).  Lemma~\ref{lem:goodpackage}
shows that \(\alpha\) is a unit in an unramified extension.  Put
\(\beta=\alpha^g\).  Then \(\beta\) is a root of the corresponding primitive
orientation \(Q_{A,B}\) or \(Q_{B,A}\), and it is again a unit in an
unramified extension.  For \(p\geq5\), this contradicts
Lemma~\ref{lem:borisovramification}(a).

Now let \(p=3\) and \(g=1\).  Every root of \(h\) is a root of the good
second orientation and therefore lies in \(\Q_3^{\mathrm{nr}}\).  The fixed
first orientation has at most one such unit root by
Lemma~\ref{lem:borisovramification}(b).  Since \(h\) is separable,
\(\deg h\leq1\).  The rational-root classification
(Theorem~\ref{thm:rational}) forces the package \(\{1,2\}\).  The same
theorem would force the second package to be \(\{1,2\}\) as well, contrary
to \(3\nmid cd(c+d)\).
\end{proof}

\begin{corollary}[Support compatibility]
\label{cor:supportcompatibility}
Fix \(\{c,d\}\) and let
\[
 S=\{p:p\mid6cd(c+d)\}.
\]
If \(P_{a,b}\) shares a factor with \(P_{c,d}\), and
\(a=gA,b=gB\) is its primitive reduction, then
\begin{equation}
 \operatorname{Supp}(AB(A+B))\subseteq S.
 \tag{11.3}
\end{equation}
If both packages are primitive, their primitive additive-triple products have
the same prime support outside \(\{2,3\}\).
\end{corollary}

\begin{proof}
Every prime \(p\geq5\) in the left side but outside \(S\) is excluded by
Theorem~\ref{thm:fullseparation}; the primes \(2,3\) were inserted into
\(S\).  For two primitive packages, apply the theorem in both directions.
\end{proof}

\section{The \texorpdfstring{\(S\)}{S}-unit reduction and finite primitive shapes}

Once the relevant prime support is controlled, the possible primitive shapes become finite.  This section packages that finiteness in the language of \(S\)-unit equations so that later counting arguments have a fixed arithmetic input.
\label{sec:sunit}

The support statement is much stronger than smoothness of the primitive total:
all three pairwise coprime entries \(A,B,A+B\) are supported on one fixed
finite set.

\begin{theorem}[Finite primitive-shape theorem]
\label{thm:finiteshapes}
Fix \(\{c,d\}\) and let \(S\) be as in
Corollary~\ref{cor:supportcompatibility}.  There are only finitely many
ordered primitive pairs \((A,B)\) that can occur as primitive reductions of
sharing partners.  If \(s=|S|\), their number is at most
\begin{equation}
 2^{16s+8}.
 \tag{12.1}
\end{equation}
\end{theorem}

\begin{proof}
Let \(N=A+B\), and put
\[
 u=\frac{A}{N},\qquad v=\frac{B}{N}.
\]
By \((11.3)\), both \(u\) and \(v\) belong to the rational \(S\)-unit
group \(U_S\), and \(u+v=1\).  The group
\(U_S\times U_S\) has rank \(2s\).  The theorem of
Beukers--Schlickewei \cite{BeukersSchlickewei1996} bounds the number of
solutions of \(x+y=1\) in a multiplicative subgroup of
\((\C^\times)^2\) of rank \(r\) by \(2^{8r+8}\), giving \((12.1)\).
The map \((A,B)\mapsto(A/N,B/N)\) is injective on primitive ordered pairs,
because each fraction is already in lowest terms.
\end{proof}

There is an elementary final step from shapes to actual scaled packages.

\begin{proposition}[Finite scaling for a fixed shape]
\label{prop:finitescaling}
Fix a package \(P_{c,d}\) and a primitive shape \((A,B)\).  Only finitely
many positive integers \(g\) can make \(P_{gA,gB}\) share a factor with
\(P_{c,d}\).
\end{proposition}

\begin{proof}
A shared factor supplies a root \(\alpha\) of \(P_{c,d}\) such that
\(\beta=\alpha^g\) is a root of \(P_{A,B}\).  There are finitely many
choices of \((\alpha,\beta)\), and \(|\alpha|\neq1\).  For a fixed pair,
\(|\alpha|^g=|\beta|\) determines at most one positive integer \(g\).
\end{proof}

Theorem~\ref{thm:finiteshapes} and Proposition~\ref{prop:finitescaling}
give a second proof of local finiteness.  Unlike the direct Archimedean proof,
they also describe the arithmetic shape of the finite neighborhood.  This is
the main reason target (a) deserves renewed attention: failure of privacy for
\(P_{c,d}\) is now a finite divisibility assertion among explicitly
constrained \(S\)-unit shapes and finitely many scalings.

\section{Global proper-divisibility rigidity}

A special danger is that one package might divide another after a composition.  The results in this section rule out that kind of proper divisibility for primitive unequal pairs, preparing the self-power arguments that follow.
\label{sec:nodivision}

The next theorem now includes the boundary primitive shapes.  Its principal
algebraic input is an exact denominator calculation.

\begin{lemma}[Exact diagonal denominators]\label{lem:diagonaldenominator}
Let \(a,b\) be coprime positive integers, put \(N=a+b\), and set
\[
 F(x)=1-\frac Nb x^a+\frac ab x^N.
\]
For every \(k\geq0\),
\begin{equation}
 [x^{ka}]F(x)^{-1}=\frac{C_k}{b^k},
 \tag{13.1}
\end{equation}
where
\begin{equation}
 C_k=N^k+
 \sum_{1\leq r\leq\lfloor k/N\rfloor}
 \binom{k-br}{ar}N^{k-Nr}(-a)^{ar}b^{br}.
 \tag{13.2}
\end{equation}
In particular, \(\gcd(C_k,b)=1\), so the coefficient in \((13.1)\) has
denominator exactly \(b^k\).
\end{lemma}

\begin{proof}
Expand geometrically:
\[
 F(x)^{-1}=\sum_{j\geq0}
 \left(\frac Nb x^a-\frac ab x^N\right)^j.
\]
If exactly \(t\) of the \(j\) selected monomials have exponent \(N\), the
resulting exponent is \((j-t)a+tN=ja+tb\).  For this to equal \(ka\),
coprimality forces
\[
 t=ar,\qquad j=k-br
\]
for an integer \(r\geq0\), and \(t\leq j\) is equivalent to \(Nr\leq k\).
The corresponding contribution is
\[
 \binom{k-br}{ar}\frac{N^{k-Nr}(-a)^{ar}}{b^{k-br}}.
\]
Putting all terms over \(b^k\) gives \((13.2)\).  Every term with \(r\geq1\)
is divisible by every prime dividing \(b\), whereas
\(N^k\equiv a^k\not\equiv0\pmod p\) for \(p\mid b\).  Thus \(C_k\) is
coprime to \(b\).
\end{proof}

\begin{theorem}[Proper-divisibility rigidity]\label{thm:nodivision}
Let \(a,b\) be coprime positive integers with \(a\neq b\).  For arbitrary
positive \(c,d\),
\begin{equation}
 Q_{a,b}(x)\mid Q_{c,d}(x)\quad\text{in }\Q[x]
 \qquad\Longrightarrow\qquad(c,d)=(a,b).
 \tag{13.3}
\end{equation}
\end{theorem}

\begin{proof}
If \(c=d\), then \(Q_{c,c}=c\) is constant and cannot be divisible by the
nonconstant polynomial \(Q_{a,b}\).  If the unordered packages agree, the
no-three theorem applied to the two
complementary edges shows that the two orientations are coprime; hence the
orientations in \((13.3)\) must agree.  We therefore assume that the packages
are distinct, so Corollary~\ref{cor:triples} gives disjoint additive triples.

First suppose \(a,b>1\).  The polynomial \(Q_{a,b}\) is primitive in
\(\Z[x]\), because its leading and constant coefficients are the coprime
integers \(a,b\).  Gauss's lemma gives
\[
 Q_{c,d}=Q_{a,b}R,\qquad R\in\Z[x].
\]
Comparison of leading and constant coefficients gives
\begin{equation}
 c=au,\qquad d=bv,\qquad u,v\in\Z_{>0}.
 \tag{13.4}
\end{equation}
Disjointness of the additive triples gives \(u,v\geq2\).

Let \(h=\gcd(c,d)\), let \([h](x)=(x^h-1)/(x-1)\), and put
\(S(x)=[h](x)^2R(x)\in\Z[x]\).  Writing \(N=a+b\), the identity between
the trinomial numerators is
\[
 H_{c,c+d}(x)=H_{a,N}(x)S(x).
\]
Reduction modulo the ideal \((x^{c})=(x^{au})\) gives
\[
 b v\equiv
 b\left(1-\frac Nb x^a+\frac ab x^N\right)S(x)
 \pmod{x^{au}},
\]
and hence
\begin{equation}
 S(x)\equiv
 v\left(1-\frac Nb x^a+\frac ab x^N\right)^{-1}
 \pmod{x^{au}}.
 \tag{13.5}
\end{equation}
For \(0\leq k<u\), the coefficient of \(x^{ka}\) on the left is integral.
Lemma~\ref{lem:diagonaldenominator} therefore gives \(b^k\mid v\), and at
\(k=u-1\),
\begin{equation}
 b^{u-1}\mid v.
 \tag{13.6}
\end{equation}
Reciprocating the divisibility and repeating the argument gives
\begin{equation}
 a^{v-1}\mid u.
 \tag{13.7}
\end{equation}
Thus \(v\geq2^{u-1}\) and \(u\geq2^{v-1}\).  If, say,
\(u\leq v\), then \(u\geq2^{v-1}\geq2^{u-1}\); for \(u\geq3\) this is
impossible.  Hence \(u=2\), and then \(2\geq2^{v-1}\) gives \(v=2\).
Thus \(u=v=2\).  The
assumed divisibility then gives
\(Q_{a,b}(x)\mid Q_{a,b}(x^2)\).  Squaring maps the finite root set into
itself, so every root is eventually periodic under squaring and hence a root
of unity, contradicting Lemma~\ref{lem:unitcircle}.  This eliminates the
interior case.

It remains to treat a boundary shape.  By reciprocity it is enough to assume
\(b=1\) and \(a>1\).  Comparison of leading coefficients gives
\begin{equation}
 c=au\qquad(u\in\Z_{>0}).
 \tag{13.8}
\end{equation}
Since the packages are distinct, additive-triple separation gives \(d\geq2\).
Reciprocating yields
\[
 Q_{1,a}\mid Q_{d,c}.
\]
Apply Lemma~\ref{lem:diagonaldenominator} to the source parameters \((1,a)\)
and target parameters \((d,c)\).  The target leading multiplier is \(d\),
and its constant multiplier is \(u=c/a\); taking the diagonal coefficient
with \(k=d-1\) gives
\begin{equation}
 a^{d-1}\mid u,
 \qquad\text{hence}\qquad c=au\geq a^d.
 \tag{13.9}
\end{equation}

Now
\[
 Q_{a,1}(x)=1+2x+\cdots+a x^{a-1}.
\]
By the Enestr\"om--Kakeya coefficient-ratio bound
\cite{AndersonSaffVarga1979}, every root \(\alpha\) satisfies
\begin{equation}
 r:=|\alpha|\leq\frac{a-1}{a}\leq e^{-1/a}<1.
 \tag{13.10}
\end{equation}
The same \(\alpha\) is a root of \(Q_{c,d}\), so the collision estimate
\((4.1)\), applied to the edge \(\{c,c+d\}\), gives
\[
 \frac dc\leq\frac{2r^c}{1-r^c},
 \qquad
 d\leq\frac{2c}{e^{c/a}-1}.
 \tag{13.11}
\]
The function \(x/(e^{x/a}-1)\) is decreasing for \(x>0\): with
\(y=x/a\), the numerator of its derivative is
\(e^y(1-y)-1<0\).  Using \((13.9)\) in \((13.11)\), and then
\(e^y-1>y^2\) for \(y\geq2\), yields
\[
 d\leq\frac{2a^d}{e^{a^{d-1}}-1}<2a^{2-d}.
\]
For \(d=2\) the right side is strictly smaller than \(2\), and for
\(d\geq3\) it is at most \(1\), a contradiction.  Thus no distinct target
exists.  Reciprocity treats the case \(a=1\).
\end{proof}

\begin{corollary}[Primitive irreducibility gives global isolation]
\label{cor:primitiveisolation}
Let \(a,b\) be coprime and unequal.  If \(Q_{a,b}\) is irreducible over
\(\Q\), then the package \(\{a,b\}\) is globally isolated from the entire
primitive and nonprimitive family.
\end{corollary}

\begin{proof}
The reciprocal \(Q_{b,a}\) is irreducible as well.  A common package factor
would make one of these two polynomials divide one orientation of the other
package, contradicting Theorem~\ref{thm:nodivision} unless the packages agree.
\end{proof}

\section{The Schinzel fibre bound, self-power rigidity, and isolation}

This section proves the isolation theorem for irreducible primitive bases.  The key point is that a shared factor in a lift would force an internal power relation for a root; the fibre bound and Kummer analysis exclude that relation.
\label{sec:kummer}

The proper-divisibility theorem controls a whole primitive polynomial.  We
next make Schinzel's trinomial-fibre argument explicit, then combine it with a
local Newton-edge obstruction and a Frobenius-quotient argument.

\begin{lemma}[Elementary binomial-fibre lemma]\label{lem:schinzelfibre}
Let \(0<m<n\), let \(d\geq1\), and assume
\(\gcd(\gcd(m,n),d)=1\).  For every \(c\in\C^\times\),
\[
 \deg\gcd\bigl(H_{m,n}(x),x^d-c\bigr)\leq1.
\]
\end{lemma}

\begin{proof}
Suppose that the gcd has degree at least two.  Since \(x^d-c\) is
separable, there are two distinct common roots \(\alpha\) and
\(\zeta\alpha\), where \(\zeta^d=1\) and \(\zeta\neq1\).  Put
\(g=\gcd(m,n)\).  The hypothesis gives \(\gcd(g,d)=1\), so
\(\xi\mapsto\xi^g\) is an automorphism of \(\mu_d\).  Because
\[
 H_{m,n}(x)=gH_{m/g,n/g}(x^g),
\]
the two transformed roots remain distinct, and it is enough to treat
\(\gcd(m,n)=1\).  Set \(r=n-m\); then \(m,r,n\) are pairwise coprime.

Write
\[
 u=\zeta^m,\qquad v=\zeta^n,\qquad
 X=\alpha^m,\qquad Y=\alpha^n.
\]
The two root equations are
\[
 mY-nX=-r,\qquad mvY-nuX=-r.
\]
If \(u=v\), these equations give \(u=v=1\), and the coprimality of
\(m,n\) gives \(\zeta=1\), a contradiction.  Hence
\[
 X=\frac{r(v-1)}{n(v-u)},\qquad
 Y=\frac{r(u-1)}{m(v-u)}.
\]
Since \(X^n=Y^m\), one obtains the master identity
\begin{equation}
 m^m r^r(\zeta^n-1)^n
 =n^n\zeta^{mr}(\zeta^m-1)^m(\zeta^r-1)^r.
 \tag{14.F1}\label{eq:fibre-master}
\end{equation}
None of \(\zeta^m,\zeta^n,\zeta^r\) is \(1\), by the displayed formulas
and \(\alpha\neq0\).

Let \(D=\ord(\zeta)=\ell^a s\), where \((\ell,s)=1\).  For each
rational prime \(\ell\), choose a place of \(\Q(\zeta)\) above \(\ell\)
and normalize its valuation by \(v_\ell(\ell)=1\).  The standard
cyclotomic valuation formula is
\[
 v_\ell(1-\zeta^k)=
 \begin{cases}
 0,&s\nmid k,\\[2mm]
 \displaystyle
 \frac1{\varphi(\ell^{a-\min(a,v_\ell(k))})},
 &s\mid k,\ \zeta^k\neq1.
 \end{cases}
\]
It follows by separating the prime-to-\(\ell\) Teichmuller part from the
\(\ell\)-power part and using \(\Phi_{\ell^j}(1)=\ell\).

Assume first that \(\ell\mid m\), and put \(t=v_\ell(m)\).  Taking
valuations in \eqref{eq:fibre-master} gives
\begin{equation}
 mt+n v_\ell(1-\zeta^n)
 =m v_\ell(1-\zeta^m)+r v_\ell(1-\zeta^r).
 \tag{14.F2}\label{eq:fibre-valuation}
\end{equation}
If \(s=1\), then \(a>t\), because \(\zeta^m\neq1\), and the valuation
identity becomes
\[
 t+\frac1{\varphi(\ell^a)}
 =\frac1{\varphi(\ell^{a-t})},
\]
or
\[
 t\varphi(\ell^a)+1=\ell^t.
\]
This is impossible: since \(a>t\),
\(t\ell^{a-1}(\ell-1)+1>\ell^t\).

Suppose \(s>1\).  Since \(m,r,n\) are pairwise coprime, \(s\) divides at
most one of them.  If it divides none, the valuation identity gives
\(mt=0\).  If \(s\mid n\), its left side has an additional positive term
and its right side has none.  If \(s\mid r\), then
\(mt=r/\varphi(\ell^a)\), so \(m\mid r\), contrary to \((m,r)=1\).
Finally, if \(s\mid m\), then
\[
 t=\frac1{\varphi(\ell^{a-t})},
\]
so necessarily
\[
 \ell=2,\qquad t=1,\qquad a=2.
\]
If instead \(\ell\mid n\), with \(t=v_\ell(n)\), valuation of
\eqref{eq:fibre-master} gives
\[
 n v_\ell(1-\zeta^n)
 =nt+m v_\ell(1-\zeta^m)+r v_\ell(1-\zeta^r).
\]
When \(s=1\), division by \(n=m+r\) gives exactly the same impossible
identity
\[
 t+\frac1{\varphi(\ell^a)}
 =\frac1{\varphi(\ell^{a-t})}.
\]
When \(s>1\), the prime-to-\(\ell\) part \(s\) divides at most one of
\(m,r,n\).  If \(s\nmid n\), the right side is strictly positive while the
left side is zero; if \(s\mid n\), the identity reduces to
\(t=1/\varphi(\ell^{a-t})\).  Thus again
\(\ell=2\), \(t=1\), and \(a=2\).

Finally, if \(\ell\mid r\), apply the already proved \(\ell\mid m\) case to
the reciprocal edge \((r,n)\): the identity
\[
 x^nH_{r,n}(x^{-1})=H_{m,n}(x)
\]
sends the two roots \(\alpha,\zeta\alpha\) to
\(\alpha^{-1},\zeta^{-1}\alpha^{-1}\), and \(\zeta^{-1}\) has the same
order as \(\zeta\).  Hence every prime divisor of \(mrn\) would have to be
\(2\), occurring to the first power.  Pairwise coprimality and \(n=m+r\)
force \((m,r,n)=(1,1,2)\).  But
\(H_{1,2}(x)=(x-1)^2\), which cannot share two distinct roots with the
separable polynomial \(x^d-c\).  This contradiction proves the lemma.
\end{proof}

\begin{theorem}[Confluent Schinzel degree bound]
\label{thm:schinzelbound}
Let \(0<m<n\), \(0<p<q\), and put
\[
 d_1=\gcd(m,n),\qquad d_2=\gcd(p,q).
\]
If \(\gcd(d_1,d_2)=1\), then
\begin{equation}
 \deg\gcd\bigl(H_{m,n},H_{p,q}\bigr)
 \leq
 \min\left\{\frac n{d_1},\frac q{d_2}\right\}.
 \tag{14.1}
\end{equation}
Consequently, if a noncyclotomic irreducible polynomial \(h\) of degree
\(e\) divides both trinomials, then
\begin{equation}
 e\leq
 \min\left\{\frac n{d_1},\frac q{d_2}\right\}-2.
 \tag{14.2}
\end{equation}
\end{theorem}

\begin{proof}
Write \(p=d_2p'\), \(q=d_2q'\).  As a polynomial in
\(y=x^{d_2}\), the second trinomial has degree \(q'\), so over \(\C\)
\[
 H_{p,q}(x)=\lambda
 \prod_{j=1}^{q'}(x^{d_2}-c_j),
\]
where the \(c_j\neq0\) are counted with multiplicity.  Since
\(\gcd(d_1,d_2)=1\), Lemma~\ref{lem:schinzelfibre} applies to every fibre,
and hence
\[
 \deg\gcd(H_{m,n},H_{p,q})
 \leq\sum_{j=1}^{q'}
 \deg\gcd(H_{m,n},x^{d_2}-c_j)
 \leq q'=\frac q{d_2}.
\]
Interchanging the two trinomials gives the other half of \((14.1)\).  Their
only common roots of unity are at \(1\), with multiplicity two in each.
Thus a common noncyclotomic factor of degree \(e\), together with
\((x-1)^2\), lies in the gcd, proving \((14.2)\).  The only external input in this proof has therefore been removed: the fibre estimate is the elementary Lemma~\ref{lem:schinzelfibre}.
\end{proof}

\subsection{Two internal-power obstructions}

The first lemma isolates the exact local mechanism used at primes dividing an
additive triple.

\begin{lemma}[Newton-edge power obstruction]\label{lem:edgepower}
Let \(L/\Q_p\) be unramified, let \(\zeta\in L\) be a root of unity of
order prime to \(p\), and let \(\beta\) be algebraic over \(L\).  Suppose
that for some integer \(e\geq1\),
\[
 v_p(\beta-\zeta)=\frac1e,
 \qquad [L(\beta):L]\leq e.
\]
Then \(\beta\notin L(\beta)^p\).
\end{lemma}

\begin{proof}
The value \(1/e\) in the value group forces the ramification index of
\(L(\beta)/L\) to be divisible by \(e\).  The degree hypothesis therefore
forces equality throughout: the extension has degree and ramification index
exactly \(e\), and its value group is \(e^{-1}\Z\).

Suppose \(\beta=w^p\) with \(w\in L(\beta)\).  Frobenius is an automorphism
of the residue field of \(L\), so there is a Teichmuller root
\(\eta\in L\) with \(\eta^p=\zeta\).  Reduction modulo the maximal ideal
gives \(w\equiv\eta\).  Write \(w=\eta+\delta\); then
\(v_p(\delta)\geq1/e\).  In
\[
 \beta-\zeta=(\eta+\delta)^p-\eta^p
 =\sum_{j=1}^{p-1}\binom pj\eta^{p-j}\delta^j+\delta^p,
\]
every intermediate term has valuation at least \(1+j/e>1/e\), and the last
has valuation at least \(p/e>1/e\).  This contradicts
\(v_p(\beta-\zeta)=1/e\).
\end{proof}

The odd-prime arguments below need only the following root packet.  It is
best taken directly from Borisov's complete Newton-polygon description rather
than inferred from a partial coefficient calculation.

\begin{lemma}[The center-\(1\) packet in every odd-prime role]
\label{lem:centeroneedge}
Let \(A,B\) be coprime, put \(N=A+B\), and let \(p\) be odd.  If
\(p^\kappa\) exactly divides one of \(A,B,N\), then exactly \(p-2\) roots
\(\gamma\) of \(Q_{A,B}\) satisfy
\[
 v_p(\gamma-1)=\frac1{p-2}.
\]
Consequently the Newton polygon of \(Q_{A,B}(1+Y)\) has a side of
horizontal length \(p-2\) and slope \(-1/(p-2)\).
\end{lemma}

\begin{proof}
Under the dictionary
\[
 Q_{A,B}=f_{N,A,B}
\]
for Borisov's notation \((a,b,c)=(N,A,B)\).  If \(p\mid N\), the asserted
packet is exactly Lemma~2.9 of \cite{Borisov1998}.  If \(p\mid B\),
Lemma~2.10 first separates the \(A\) nonunit roots and states that the
remaining roots occur in root-of-unity clusters exactly as in the total role;
the cluster centered at \(1\) therefore contains the same \(p-2\) roots.
If \(p\mid A\), Lemma~2.11 gives the reciprocal endpoint statement and the
same center-\(1\) packet.  The final assertion follows from the standard
Newton-polygon root--slope correspondence: roots with
\(v_p(\gamma-1)=1/(p-2)\) contribute a side of slope \(-1/(p-2)\), and
their number is its horizontal length.
\end{proof}

For an integer \(t\) and a prime \(p\), put
\[
 q_p(t)=\frac{t^p-t}{p},
 \qquad
 C_p(X,Y)=\frac{(X+Y)^p-X^p-Y^p}{p}.
\]
Both are integral.  The following factor-level theorem is stronger than what
is needed for an irreducible whole polynomial.

\begin{theorem}[Good-prime Frobenius bound for a proper factor]
\label{thm:frobeniusfactor}
Let \(A,B\) be coprime positive integers, put \(N=A+B\), and let \(h\) be
an irreducible factor of \(F=Q_{A,B}\) of degree \(e\).  Let \(\beta\) be
a root of \(h\).  If \(p\) is an odd prime with
\[
 p\nmid ABN
\]
and \(\beta\in\Q(\beta)^p\), then
\begin{equation}
 e\leq p-2.
 \tag{14.3}
\end{equation}
\end{theorem}

\begin{proof}
Choose \(h\in\Z[x]\) primitive with positive leading coefficient and write
\(F=hk\) in \(\Z[x]\).  The leading and constant coefficients of \(h\)
are coprime, because they divide the coprime integers \(A\) and \(B\).
Suppose \(\beta=w^p\) in \(K=\Q(\beta)\).  Then \(K=\Q(w)\), and the
norm identity gives
\[
 \operatorname{lc}(h)=u^p,\qquad h(0)=v_0^p
\]
for coprime integers \(u>0\) and \(v_0\), where the sign is absorbed into
\(v_0\) because \(p\) is odd.

Let \(G\in\Z[x]\) be the primitive minimal polynomial of \(w\), with
positive leading coefficient \(t\), and put \(s=G(0)\).  Since
\(G\mid h(x^p)\), one has \(t\mid u^p\) and \(s\mid v_0^p\), hence
\((t,s)=1\).  The norm identity gives
\[
 \left(\frac{s}{t}\right)^p
 =\left(\frac{v_0}{u}\right)^p.
\]
The odd \(p\)-th power map is injective on \(\Q\), and both fractions are
reduced.  Therefore
\[
 t=u,\qquad s=v_0.
\]
In particular, \(p\nmid u\), because \(u^p\mid A\).

Let \(w_1,\ldots,w_e\) be the conjugates of \(w\).  Since
\(\Q(w)=\Q(\beta)\), the multiset \(\{w_i^p\}\) is exactly the multiset
of conjugates of \(\beta\).  Since \(G/u\) is monic over
\(\Z_{(p)}\), the \(w_i\) are integral over \(\Z_{(p)}\).  In the
integral closure \(\mathcal O\) of \(\Z_{(p)}\),
\[
 \begin{aligned}
 G(x)^p
 &=u^p\prod_i(x-w_i)^p\\
 &\equiv u^p\prod_i(x^p-w_i^p)
 =h(x^p)\pmod {p\mathcal O[x]}.
 \end{aligned}
\]
Both sides have rational coefficients, and
\(p\mathcal O\cap\Q=p\Z_{(p)}\), so this is a coefficientwise congruence
in \(\Z_{(p)}[x]\).  Coefficientwise freshman's dream also gives
\(G(x)^p\equiv G(x^p)\) and
\(h(x)^p\equiv h(x^p)\pmod p\).  Hence
\(\overline{G(x^p)}=\overline{h(x^p)}\); injectivity of
\(f(x)\mapsto f(x^p)\) on \(\mathbb F_p[x]\) yields
\[
 \overline G=\overline h.
\]

Define
\[
 \Phi_p(f)(x)=\frac{f(x)^p-f(x^p)}p.
\]
Because \(u\) is a \(p\)-adic unit, the quotient
\[
 R=\Z_{(p)}[x]/(G)=\Z_{(p)}[x]/(G/u)
\]
is a free, hence torsion-free, \(\Z_{(p)}\)-module of rank \(e\).  Since \(\overline G=\overline h\), write
\(G=h+pE\) with \(E\in\Z_{(p)}[x]\).  In \(R\), one has
\[
 h=-pE,\qquad h(x^p)=0,
\]
the second equality following from \(G\mid h(x^p)\).  Therefore
\[
 \Phi_p(h)=(-1)^p p^{p-1}E^p\in pR.
\]
Since \(R/pR\simeq\mathbb F_p[x]/(\overline h)\), this proves
\[
 \overline h\mid\overline{\Phi_p(h)}.
\]

The product rule
\[
 \Phi_p(fg)=f^p\Phi_p(g)+g(x^p)\Phi_p(f)
\]
transfers this vanishing from \(h\) to \(F\).  Applying the
same rule to
\[
 H(x)=(x-1)^2F(x)=Ax^N-Nx^A+B
\]
transfers it from \(F\) to \(H\): at a root \(\alpha\) of
\(\overline h\), the second product-rule term contains
\(F(\alpha^p)=F(\alpha)^p=0\) in characteristic \(p\).  Thus
\[
 \overline{\Phi_p(H)}(\alpha)=0
 \quad\text{for every root }\alpha\text{ of }\overline h.
\]

Set
\[
 U=A\alpha^B,\qquad V=B\alpha^{-A}.
\]
The equation \(H(\alpha)=0\) gives \(U+V=N\).  Put
\[
 q_p(t)=\frac{t^p-t}{p},\qquad
 C_p(X,Y)=\frac{(X+Y)^p-X^p-Y^p}{p}.
\]
To make the coefficient calculation explicit, write
\[
 H=Ax^N-Nx^A+B.
\]
In \(\Phi_p(H)\), the mixed terms from the three summands contribute
\(x^{pA}C_p(U,V)\), while the coefficient corrections contribute
\[
 x^{pA}\left(\frac{q_p(A)}A U^p
 +\frac{q_p(B)}B V^p-q_p(N)\right).
\]
Using \(V=N-U\) gives
\begin{equation}
 \alpha^{-pA}\overline{\Phi_p(H)}(\alpha)
 =T_{p,A,B}(U),
 \tag{14.6}
\end{equation}
where
\begin{equation}
 T_{p,A,B}(U)=
 C_p(U,N-U)
 +\frac{q_p(A)}A U^p
 +\frac{q_p(B)}B(N-U)^p
 -q_p(N)
 \in\mathbb F_p[U].
 \tag{14.7}
\end{equation}
The hypothesis \(p\nmid ABN\) is essential here.  The polynomial has degree
at most \(p\), and the identity
\[
 q_p(N)=q_p(A)+q_p(B)+C_p(A,B)
\]
gives
\[
 T_{p,A,B}(A)=T'_{p,A,B}(A)=0,
 \qquad
 T''_{p,A,B}(A)=A^{-1}+B^{-1}=\frac N{AB}\neq0.
\]
Thus \(T_{p,A,B}\) is nonzero and \(U=A\) is an exactly double root.

It remains to count the other roots.  The reduction of \(F\) is square-free.
Indeed, a common root of \(F\) and \(F'\) would be a common root of \(H\)
and
\[
 H'(x)=ANx^{A-1}(x^B-1).
\]
It would satisfy \(x^A=x^B=1\), hence \(x=1\), whereas
\[
 F(1)=\frac{ABN}{2}\not\equiv0\pmod p.
\]
Therefore \(\overline h\) has \(e\) distinct roots.  The map
\[
 \alpha\longmapsto U=A\alpha^B
\]
is injective on roots of \(F\): equality of \(U\) also gives equality of
\(V=N-U\), so the quotient of two roots has both its \(A\)-th and
\(B\)-th powers equal to \(1\).  Moreover \(U=A\) would force
\(\alpha^A=\alpha^B=1\), hence \(\alpha=1\), which is not a root.
Consequently the \(e\) roots of \(\overline h\) give \(e\) distinct roots
of \(T_{p,A,B}\) outside the exactly double root \(A\).  The remaining
root-multiplicity budget is at most \(p-2\), proving \((14.3)\).
\end{proof}

\begin{theorem}[Self-power rigidity]\label{thm:selfpower}
Let \(A,B\) be coprime positive integers with \(A\neq B\), and suppose
\(F=Q_{A,B}\) is irreducible.  If \(\beta\) is a root and
\(K=\Q(\beta)\), then
\begin{equation}
 \beta\notin K^p
 \qquad\text{for every prime }p.
 \tag{14.8}
\end{equation}
\end{theorem}

\begin{proof}
Suppose \(\beta=w^p\) with \(w\in K\).  Put \(N=A+B\) and
\(D=N-2\).  Since the monic minimal polynomial of \(\beta\) is \(F/A\),
\begin{equation}
 \operatorname N_{K/\Q}(w)^p
 =\operatorname N_{K/\Q}(\beta)
 =(-1)^D\frac BA.
 \tag{14.9}
\end{equation}
For odd \(p\), coprimality implies
\[
 A=u^p,\qquad B=v^p
\]
for coprime positive integers \(u,v\).  For \(p=2\), equation \((14.9)\)
forces \(D\) even and \(A,B\) to be coprime squares; hence both are odd and
\(v_2(N)=1\).

First suppose \(p\mid ABN\).  For odd \(p\),
Lemma~\ref{lem:centeroneedge} provides in each of the three possible roles a
cluster at \(1\) containing exactly \(e=p-2\) roots
\(\gamma\) with
\[
 v_p(\gamma-1)=\frac1e.
\]
For \(p=2\), the preceding parity conclusion gives \(v_2(N)=1\).  Since
\(A,B\) are distinct coprime odd squares, \(N\geq10\), so \(N/2>1\);
Borisov's Lemma~2.8 provides a nontrivial prime-to-\(2\) center \(\zeta\)
with exactly \(e=2\) roots satisfying
\(v_2(\gamma-\zeta)=1/2\).

Because \(F\) is irreducible, a \(\Q\)-embedding of \(K\) into
\(\overline{\Q}_p\) may send \(\beta\) to one of these roots.  Let \(L\)
be the unramified field generated by the cluster center.  Every \(L\)-embedding fixes the cluster center and preserves the unique
extension of the \(p\)-adic valuation, so
\(v_p(\sigma(\beta)-\zeta)=v_p(\beta-\zeta)\).  Thus every
\(L\)-conjugate remains on the same edge, which contains only \(e\) roots;
hence \([L(\beta):L]\leq e\).  Lemma~\ref{lem:edgepower} contradicts
\(\beta=w^p\).

It remains that \(p\nmid ABN\).  Then \(p\) is odd, and
Theorem~\ref{thm:frobeniusfactor}, applied to the whole irreducible polynomial
of degree \(D\), gives \(D\leq p-2\).  But \(A=u^p\), \(B=v^p\), and
\(u\neq v\), so
\[
 D=A+B-2\geq1+2^p-2=2^p-1>p-2,
\]
a contradiction.  This proves \((14.8)\).
\end{proof}

\subsection{Kummer reduction and its complete elimination}

\begin{lemma}[Degree under extraction of a root]
\label{lem:kummerdegree}
Let \(f\in\Q[x]\) be irreducible of degree \(D\), let \(r\geq1\), and let
\(w\) be algebraic with \(f(w^r)=0\).  Put \(\beta=w^r\).  Then
\[
 [\Q(w):\Q]=[\Q(w):\Q(\beta)]D.
\]
In particular, equality of the two degrees is equivalent to
\(\Q(w)=\Q(\beta)\).
\end{lemma}

\begin{proof}
This is the tower law for \(\Q(\beta)\subseteq\Q(w)\), together with the
irreducibility of \(f\).
\end{proof}

\begin{theorem}[Exact Kummer reduction]
\label{thm:kummerreduction}
Let \(A,B\) be coprime and unequal, put \(N=A+B\), and suppose
\(Q_{A,B}\) is irreducible.  Assume that \(\{rA,rB\}\) shares a factor
with \(\{sC,sD\}\), where \(\gcd(C,D)=1\).  Independently interchange \(A,B\) and \(C,D\), as necessary, so that the
selected common root lies in \(Q_{rA,rB}\) and \(Q_{sC,sD}\).  Put
\[
 g=\gcd(r,s),\qquad r_0=r/g,\qquad s_0=s/g.
\]
Then either the two unordered packages agree, or \(r_0>1\) and a root
\(\beta\) of \(Q_{A,B}\) is an \(r_0\)-th power inside \(\Q(\beta)\).
\end{theorem}

\begin{proof}
Let \(\alpha\) be the common root and set \(w=\alpha^g\).  Then
\[
 Q_{A,B}(w^{r_0})=0,
 \qquad Q_{C,D}(w^{s_0})=0.
\]
Let \(k_w\) be the minimal polynomial of \(w\), and put \(M=C+D\).
Because \(|w|=|\alpha|^g\neq1\), the polynomial \(k_w\) is noncyclotomic.
It divides both normalized package polynomials
\[
 Q_{A,B}(x^{r_0}),
 \qquad Q_{C,D}(x^{s_0}).
\]
The associated exponent-pair gcds are \(r_0,s_0\), which are coprime.
Theorem~\ref{thm:schinzelbound} gives
\[
 \deg k_w\leq\min(N,M)-2\leq N-2.
\]
On the other hand, \(\beta=w^{r_0}\) is a root of the irreducible polynomial
\(Q_{A,B}\) of degree \(N-2\), so Lemma~\ref{lem:kummerdegree} gives the
reverse inequality.  Hence \(\Q(w)=\Q(\beta)\).

If \(r_0=1\), then \(Q_{A,B}\) divides \(Q_{s_0C,s_0D}\), and
Theorem~\ref{thm:nodivision} forces equality of the oriented pairs and hence
of the packages.  If \(r_0>1\), the field equality gives
\(w\in\Q(\beta)\) and \(w^{r_0}=\beta\), as asserted.
\end{proof}

\begin{theorem}[Global isolation of every lift of an irreducible shape]
\label{thm:allliftsisolated}
Let \(A,B\) be coprime and unequal, and suppose \(Q_{A,B}\) is irreducible.
Then for every \(r\geq1\), the package \(\{rA,rB\}\) is globally isolated.
\end{theorem}

\begin{proof}
An unequal-package collision would give the second alternative in
Theorem~\ref{thm:kummerreduction}.  Choose a prime \(p\mid r_0\); then
\(\beta=(w^{r_0/p})^p\) inside \(\Q(\beta)\), contradicting
Theorem~\ref{thm:selfpower}.
\end{proof}

\begin{corollary}[Individual composition irreducibility also isolates]
\label{cor:individualirreducible}
Let \(A,B\) be coprime and unequal, and \(r\geq1\).  If
\(Q_{A,B}(x^r)\) is irreducible over \(\Q\), then \(\{rA,rB\}\) is
globally isolated.
\end{corollary}

\begin{proof}
If a common factor existed, a source root \(\alpha\) would have degree
\(r(A+B-2)\).  With the notation in the preceding proof,
\([\Q(\alpha):\Q(w)]\leq g\), so
\[
 [\Q(w):\Q]\geq r_0(A+B-2).
\]
The Schinzel bound is at most \(A+B-2\), forcing \(r_0=1\); proper-divisibility
rigidity then forces equality of the packages.
\end{proof}

\begin{theorem}[Density-one global isolation]
\label{thm:densityisolation}
Let \(I(X)\) denote the number of unordered pairs \(1\leq a<b\),
\(a+b\leq X\), whose reciprocal packages are not globally isolated.  Then
\begin{equation}
 I(X)\ll X^{20/11}\log(2X).
 \tag{14.10}
\end{equation}
\end{theorem}

\begin{proof}
Write uniquely
\[
 a=rA,\qquad b=rB,
 \qquad\gcd(A,B)=1,
 \qquad A<B.
\]
By Theorem~\ref{thm:allliftsisolated}, a non-isolated package must have
reducible primitive base \(Q_{A,B}\).

To connect precisely with Borisov's notation, his triple is
\[
 a_{\rm Bor}=A+B,
 \qquad b_{\rm Bor}=A,
 \qquad c_{\rm Bor}=B.
\]
For a fixed small \(\varepsilon>0\), the contrapositive of his
Theorem~3.8 places every sufficiently large reducible primitive shape in the
complementary numerical range counted by Theorem~3.9 \cite{Borisov1998}.  The finitely many
small totals are absorbed into the constant.  Thus the number \(R(Y)\) of
reducible primitive shapes of total at most \(Y\) satisfies
\begin{equation}
 R(Y)\ll Y^{20/11}\log(2Y).
 \tag{14.11}
\end{equation}
Therefore
\[
 I(X)\leq\sum_{r\leq X/3}R(X/r)
 \ll X^{20/11}\log(2X)
 \sum_{r\geq1}r^{-20/11}
 \ll X^{20/11}\log(2X).
\]
Since the total number of unordered unequal pairs of total at most \(X\) is
asymptotic to \(X^2/4\), a density-one set of package vertices is globally
isolated.
\end{proof}

\begin{remark}
The estimate is global, not a bounded-box assertion: every certified package
is coprime to every distinct package in the entire infinite family.  The
remaining set is merely the set not eliminated by the present proof; the
theorem does not assert that its members actually share factors.
\end{remark}

\section{Factor-level power sieves, residual splitting, and exact signatures}

The previous section handles irreducible primitive bases.  For reducible bases, the same power questions must be tracked factor by factor.  The purpose here is to record exactly where endpoint primes can still create internal powers and which congruence signatures survive.
\label{sec:degreesieve}

The local cluster argument constrains proper factors without assuming
irreducibility of the whole primitive polynomial.  We first retain the
radical sieve, then resolve the equal-slope \(2\)-adic block by its residual
polynomial.

\begin{theorem}[Primitive total-radical sieve]
\label{thm:radicalsieve}
Let \(A,B\) be coprime positive integers, put \(N=A+B\), and let
\(h\in\Q[x]\) be an irreducible factor of \(Q_{A,B}\) of degree \(e\).
For every prime \(p\mid N\),
\begin{equation}
 e\equiv0\quad\text{or}\quad-2\pmod p.
 \tag{15.1}
\end{equation}
Consequently,
\begin{equation}
 \operatorname{rad}(N)\mid e(e+2).
 \tag{15.2}
\end{equation}
If the same reciprocal irreducible orbit occurs in primitive packages of
totals \(N\) and \(M\), then
\begin{equation}
 \operatorname{lcm}\bigl(\operatorname{rad}N,
                          \operatorname{rad}M\bigr)
 \mid e(e+2).
 \tag{15.3}
\end{equation}
\end{theorem}

\begin{proof}
Fix \(p^k\Vert N\).  Since \(p\nmid AB\), normalize a rational factor to
be monic over \(\Z_p\) with unit constant term, and work over the unramified
extension containing the prime-to-\(p\) cluster centers.

Assume first that \(p\) is odd.  At every nontrivial center, Borisov's
Lemma~2.8 gives nonhorizontal Newton edges of lengths
\[
 p,\quad p^2-p,\quad\ldots,\quad p^k-p^{k-1},
\]
with reduced slope denominators equal to those lengths.  At the center
\(1\), Lemma~\ref{lem:centeroneedge} gives a side of length \(p-2\) and
reduced slope denominator \(p-2\); Borisov's Lemma~2.9 gives the subsequent
\(p\)-divisible lengths.  A monic local
factor has a Newton polygon with integral vertices.  Its horizontal
contribution on an edge of reduced slope denominator \(\ell\) is divisible
by \(\ell\); here \(\ell\) equals the full edge length, so the factor takes
that edge wholly or not at all.  Contributions at nontrivial centers are
divisible by \(p\), while the center \(1\) contributes either \(0\) or
\(p-2\) modulo \(p\).  This proves \((15.1)\) for odd \(p\).

For \(p=2\), one must not separate Borisov's first two formal layers when
\(k\geq2\): both have slope \(1/2\) and form a single canonical side of
length four.  Its reduced slope denominator is nevertheless \(2\), so any
local factor takes an even number of roots from that side.  Every later
nontrivial side has even reduced slope denominator, and the same is true for
every nonhorizontal side at the center \(1\).  Hence every factor degree is
even.  This is precisely \((15.1)\) modulo \(2\).

Thus every \(p\mid N\) divides \(e(e+2)\), proving \((15.2)\).  For a
common reciprocal orbit, apply the result to whichever of \(h,h^\vee\)
divides the second orientation; the reciprocal has the same degree.  This
gives \((15.3)\).
\end{proof}

\begin{corollary}[A numerical coprimality criterion]
\label{cor:radicalcriterion}
Let two primitive packages have totals \(N,M\), put
\[
 \nu=\min(N,M),\qquad
 \rho=\operatorname{lcm}\bigl(\operatorname{rad}N,
                                \operatorname{rad}M\bigr).
\]
If
\begin{equation}
 \rho>\nu(\nu-2),
 \tag{15.4}
\end{equation}
then the packages are coprime.  In particular, distinct primitive packages
with coprime squarefree totals are coprime.
\end{corollary}

\begin{proof}
A common factor has degree \(1\leq e\leq\nu-2\), so
\(e(e+2)\leq\nu(\nu-2)\), contradicting \((15.3)\)--\((15.4)\).
For coprime squarefree totals, \(\rho=NM>\nu(\nu-2)\).
\end{proof}

\subsection{The merged \texorpdfstring{\(2\)}{2}-adic side is residual-separable}

We use the first-order theorem of the residual polynomial in the following
standard form; see \cite[Theorem~1.21]{GuardiaMontesNart2008}.

\begin{lemma}[Newton--Ore residual splitting]
\label{lem:oreresidual}
Let \(L/\Q_2\) be unramified and let a polynomial over \(L\) have a
Newton side of slope \(-1/2\), horizontal length four, and separable
quadratic residual polynomial.  After a finite unramified extension
\(L'/L\), the four roots belonging to that side split into two packets of
two.  Each root \(\gamma\) in either packet satisfies
\[
 [L'(\gamma):L']\leq2.
\]
\end{lemma}

\begin{proof}
The denominator of the reduced slope is \(2\), and the side has residual
degree two.  After an unramified extension splitting the separable residual
polynomial, its two distinct linear residual factors give, by the theorem of
the residual polynomial, two factors of degree \(2\).  Each associated root
therefore has local degree at most two over the enlarged unramified base.
\end{proof}

\begin{proposition}[Residual polynomials of the merged block]
\label{prop:mergedtwo}
Let \(A,B\) be coprime, \(N=A+B\), and suppose a nontrivial
prime-to-\(2\) cluster center exists.
\begin{enumerate}[label=(\alph*)]
\item If \(N=2^\kappa n\), \(\kappa\geq2\), \(n\) odd, and
\(\zeta\neq1\) is an \(n\)-th root of unity, then the length-four
slope-\(1/2\) side at \(\zeta\) has residual polynomial
\[
 R_N(Z)=(1-\bar\zeta^A)+\bar\zeta^{-2}Z+
        \bar\zeta^{-4}Z^2.
\]
\item If \(B=2^\kappa b\), \(\kappa\geq2\), \(b\) odd, and
\(\zeta\neq1\) is a \(b\)-th root of unity, then the corresponding unit
side has residual polynomial
\[
 R_B(Z)=(1-\bar\zeta^A)+\bar\zeta^{A-2}Z+
        \bar\zeta^{A-4}Z^2.
\]
\end{enumerate}
Both polynomials are separable.  The case \(2^\kappa\Vert A\) follows by
reciprocity.
\end{proposition}

\begin{proof}
Work in the unramified field generated by \(\zeta\), and write
\[
 H_{A,N}(\zeta+Y)=\sum_{j\geq0}c_jY^j.
\]
At a nontrivial center the denominator \((x-1)^2\) is a unit, so the
relevant Newton data are unchanged on passing to \(Q_{A,B}\).

In case (a), \(\zeta^N=1\), \(A\) is odd, and
\[
 c_0=N(1-\zeta^A),
\]
\[
 c_2=\frac{AN}{2}\zeta^{-2}
       \bigl((N-1)-(A-1)\zeta^A\bigr),
\]
\[
 c_4=A\binom N4\zeta^{-4}
      -N\binom A4\zeta^{A-4}.
\]
The three valuations are respectively
\(\kappa,\kappa-1,\kappa-2\).  After division by the powers of two on the
side and reduction, the coefficients are
\[
 1-\bar\zeta^A,
 \qquad \bar\zeta^{-2},
 \qquad \bar\zeta^{-4}.
\]
Here the second term in \(c_4\) has larger valuation, while
\(2^{2-\kappa}\binom N4\) is odd.

In case (b), \(\zeta^B=1\), and direct cancellation gives
\[
 c_0=B(1-\zeta^A),
 \qquad
 c_2=\frac{ANB}{2}\zeta^{A-2},
\]
\[
 c_4=\frac{ANB}{24}
 \bigl(N^2+NA+A^2-6(N+A)+11\bigr)\zeta^{A-4}.
\]
Since \(A,N\) are odd and \(N\equiv A\pmod4\), the integer in parentheses
is congruent to \(2\pmod4\).  The normalized residual coefficients are
therefore
\[
 1-\bar\zeta^A,
 \qquad \bar\zeta^{A-2},
 \qquad \bar\zeta^{A-4}.
\]
Borisov's polygon identifies these points as the merged canonical side in
both cases.  In characteristic two,
\[
 R_N'(Z)=\bar\zeta^{-2}\neq0,
 \qquad
 R_B'(Z)=\bar\zeta^{A-2}\neq0.
\]
Thus both residual polynomials are separable.
\end{proof}

\begin{corollary}[Internal powers at every total prime]
\label{cor:totalprimepower}
Let \(h\) be an irreducible factor of a primitive \(Q_{A,B}\), and let
\(\beta\) be a root.  If \(p\mid A+B\), then
\[
 \beta\notin\Q(\beta)^p.
\]
This includes \(p=2\) for every value of \(v_2(A+B)\).
\end{corollary}

\begin{proof}
For odd \(p\), a root lies on a Newton side of some length \(\ell\) and
slope denominator \(\ell\).  Borisov's exact count gives at most \(\ell\)
roots at that center and distance, while the valuation denominator forces
the local ramification degree to be at least \(\ell\).
Lemma~\ref{lem:edgepower} applies.

For \(p=2\), sides of length two with no equal-slope merger are handled in
the same way.  If the primitive total has odd part greater than one and
\(v_2(A+B)\geq2\), Proposition~\ref{prop:mergedtwo}(a) and
Lemma~\ref{lem:oreresidual} split the four-root side, after an unramified
extension, into two degree-two packets.  Lemma~\ref{lem:edgepower} then
applies to each packet.  If the primitive total is a pure power of two,
there are no nontrivial centers; at the center \(1\), the nonzero sides have
successive lengths \(2,4,8,\ldots\) and distinct slopes, so the original
edge argument applies directly.  These cases exhaust all roots.
\end{proof}

\subsection{Endpoint and good-prime power sieves}

\begin{theorem}[Endpoint-prime internal-power sieve]
\label{thm:endpointpower}
Let \(h\) be an irreducible factor of primitive \(Q_{A,B}\), of degree
\(e\), and let \(\beta\) be a root.
\begin{enumerate}[label=(\alph*)]
\item Suppose \(p^\kappa\Vert B\), for an arbitrary prime \(p\).  If
\(\beta\in\Q(\beta)^p\), then
\begin{equation}
 e\leq A,\qquad pA\mid e\kappa.
 \tag{15.5}
\end{equation}
\item Suppose \(p^\kappa\Vert A\).  Then an internal \(p\)-th power
forces
\begin{equation}
 e\leq B,\qquad pB\mid e\kappa.
 \tag{15.6}
\end{equation}
\end{enumerate}
Equivalently, in (a),
\[
 \frac{pA}{\gcd(pA,\kappa)}\mid e;
\]
in particular \(\kappa<p\) is impossible.  The reciprocal statement holds
in (b).
\end{theorem}

\begin{proof}
We prove (a); reciprocity gives (b).  Borisov's endpoint description
\cite[Lemma~2.10]{Borisov1998} gives exactly \(A\) nonunit roots, each of
valuation \(\kappa/A\).  Every other root lies in a prime-to-\(p\)
root-of-unity cluster; the reciprocal statement uses
\cite[Lemma~2.11]{Borisov1998}.

For odd \(p\), every unit-cluster side has exact root count equal to its
reduced slope denominator, so Lemma~\ref{lem:edgepower} excludes an internal
\(p\)-th power there.  For \(p=2\), the unmerged sides are treated
identically.  If \(\kappa\geq2\) and the odd part of \(B\) is greater than
one, Proposition~\ref{prop:mergedtwo}(b) and
Lemma~\ref{lem:oreresidual} split the merged side into degree-two packets.
If the odd part is one, there are no nontrivial centers and the center-\(1\)
sides have distinct slopes.  Thus unit-cluster roots are excluded for every
\(p\) and every \(\kappa\).

If \(\beta=w^p\), every rational conjugate of \(\beta\) is likewise a
\(p\)-th power in its conjugate field.  Hence all \(e\) roots of \(h\) must
belong to the nonunit block, proving \(e\leq A\).  The leading coefficient
of \(h\) is a \(p\)-adic unit, and therefore
\[
 v_p\!\left(N_{\Q(\beta)/\Q}(\beta)\right)
 =\frac{e\kappa}{A}.
\]
The norm is a rational \(p\)-th power, so this integer is divisible by
\(p\).  Thus \(pA\mid e\kappa\).  If \(\kappa<p\), then
\(e\kappa\leq A\kappa<pA\), a contradiction.
\end{proof}

\begin{corollary}[Good-prime internal-power sieve]
\label{cor:goodpower}
Under the hypotheses of Theorem~\ref{thm:frobeniusfactor},
\[
 \operatorname{rad}(A+B)\leq e(e+2)\leq p(p-2).
\]
In particular, \(\operatorname{rad}(A+B)>p(p-2)\) rules out an internal
\(p\)-th power for every factor.
\end{corollary}

\begin{proof}
Combine Theorems~\ref{thm:radicalsieve} and
\ref{thm:frobeniusfactor}.
\end{proof}

We next refine the Kummer trigger prime by prime.  We use the prime-degree
case of Capelli's binomial criterion
\cite[Chapter~VI, Section~9]{LangAlgebra}: over a characteristic-zero field
\(K\), the polynomial \(x^p-a\), with \(p\) prime, is irreducible unless
\(a\in K^p\).

\begin{proposition}[Prime-by-prime Kummer trigger]
\label{prop:largetrigger}
Consider selected orientations of packages with primitive shapes \((A,B)\)
and \((C,D)\), scaling factors \(r,s\), and a common root \(\alpha\).  Put
\[
 g=\gcd(r,s),\quad r_0=r/g,\quad s_0=s/g,
 \quad w=\alpha^g,
 \quad L=\min(A+B,C+D)-2.
\]
Let \(\beta=w^{r_0}\) be a root of an irreducible factor
\(h\mid Q_{A,B}\) of degree \(e\).  For every prime \(p\mid r_0\), at
least one of the following holds:
\begin{enumerate}[label=(\roman*)]
\item \(pe\leq L\);
\item \(\beta\in\Q(\beta)^p\).
\end{enumerate}
If \(pe>L\), the second alternative holds and hence all applicable
conclusions of Corollary~\ref{cor:totalprimepower},
Theorem~\ref{thm:endpointpower}, and
Theorem~\ref{thm:frobeniusfactor} apply.
\end{proposition}

\begin{proof}
The minimal polynomial \(k_w\) is noncyclotomic because
\(|w|=|\alpha|^g\neq1\).  It divides
\(Q_{A,B}(x^{r_0})\) and \(Q_{C,D}(x^{s_0})\), whose exponent-pair gcds
are coprime.  Theorem~\ref{thm:schinzelbound} gives
\[
 [\Q(w):\Q]\leq L.
\]
Fix \(p\mid r_0\), put \(u=w^{r_0/p}\), and set
\(K=\Q(\beta)\).  If \(\beta\notin K^p\), Capelli's criterion makes
\(x^p-\beta\) irreducible over \(K\), so
\[
 [K(u):K]=p.
\]
Since \(K(u)\subseteq\Q(w)\), it follows that
\(pe\leq[\Q(w):\Q]\leq L\).  The contrapositive gives the second
alternative when \(pe>L\).
\end{proof}

\begin{remark}
The former condition \(2e>L\) is the special case in which the
internal-power alternative is forced simultaneously for every prime divisor
of \(r_0\).  The prime-by-prime form is stronger when the relative scaling
exponent has a large prime divisor.
\end{remark}

\subsection{Role assignments and exact factor partitions}

\begin{proposition}[Odd-prime degree signatures]
\label{prop:rolesignature}
Suppose a reciprocal irreducible orbit of degree \(e\) occurs in primitive
packages of totals \(N\) and \(M\).  Put
\[
 \rho_{\mathrm{odd}}
 =\operatorname{lcm}\bigl(\operatorname{rad}N,
                           \operatorname{rad}M\bigr)_{\mathrm{odd}}.
\]
Then there is a unique factorization
\begin{equation}
 \rho_{\mathrm{odd}}=UV,
 \qquad (U,V)=1,
 \qquad U\mid e,
 \qquad V\mid e+2,
 \tag{15.7}
\end{equation}
obtained by assigning each odd prime divisor of \(\rho_{\mathrm{odd}}\) to
exactly one of the two roles \(e\equiv0\) or \(e\equiv-2\).  If
\(2\mid NM\), then \(e\) is even.
\end{proposition}

\begin{proof}
For every odd \(q\mid\rho_{\mathrm{odd}}\),
Theorem~\ref{thm:radicalsieve} says that \(q\mid e(e+2)\).  Since
\(\gcd(e,e+2)\mid2\), the prime \(q\) divides exactly one of the two
integers.  Multiplying the primes in their two roles gives \((15.7)\).  The
assertion at \(2\) is the parity part of the same theorem.
\end{proof}

\begin{corollary}[Exact factor-partition signature]
\label{cor:partitionsignature}
Let \(N=qk\), where \(q\) is an odd prime, and let the irreducible factor
degrees of a primitive \(Q_{A,B}\) be \(e_1,\ldots,e_t\).  Write uniquely
\[
 e_i=qa_i-2\varepsilon_i,
 \qquad a_i\geq1,
 \qquad \varepsilon_i\in\{0,1\},
\]
and put \(t_q=\sum_i\varepsilon_i\).  Then for an integer \(j\geq0\),
\begin{equation}
 t_q=1+qj,
 \qquad
 \sum_i a_i=k+2j,
 \qquad
 (q-2)j\leq k-1.
 \tag{15.8}
\end{equation}
In particular, if \(q>k+1\), then \(j=0\): exactly one factor has the
\(-2\)-role at \(q\), and the total number of irreducible factors is at
most \(k=N/q\).  If \(N=2q\) with \(q\geq5\), the \(q\)-signature
alone would force any reducible factorization to have two factors of degrees
\(q-2\) and \(q\).
\end{corollary}

\begin{proof}
Summing the factor degrees gives
\[
 q\sum_i a_i-2t_q=qk-2.
\]
Reduction modulo \(q\) yields \(t_q\equiv1\pmod q\), so
\(t_q=1+qj\), and substitution gives \(\sum_i a_i=k+2j\).  Since each
index counted by \(t_q\) contributes at least one to \(\sum_i a_i\),
\[
 1+qj=t_q\leq\sum_i a_i=k+2j,
\]
which is the final inequality in \((15.8)\).  If \(q>k+1\), then
\(q-2>k-1\), forcing \(j=0\).  For \(k=2\), reducibility forces two
positive \(a_i\)'s summing to two, and exactly one \(\varepsilon_i\) is
one; the degrees are therefore \(q-2\) and \(q\).
\end{proof}

\begin{corollary}[Prime and twice-prime totals]
\label{cor:prime-totals}
Let \(A,B\) be coprime and unequal.  If \(A+B=q\), or if
\(A+B=2q\), where \(q\) is an odd prime, then \(Q_{A,B}\) is
irreducible over \(\Q\).  Consequently every positive integral lift of
such a primitive shape is globally isolated.
\end{corollary}

\begin{proof}
If the total is \(q\), Theorem~\ref{thm:radicalsieve} says that a factor
degree \(e\), with \(1\leq e\leq q-2\), is congruent to \(0\) or
\(-2\pmod q\).  The only positive possibility in this range is
\(e=q-2\), the full degree.

Now let the total be \(2q\).  The prime \(2\) in the radical sieve makes
every factor degree even.  The \(q\)-congruence shows that a proper factor
degree in the range \(1\leq e<2q-2\) can only be \(q-2\) or \(q\), both
odd.  Thus no proper factor exists.  Global isolation of all lifts follows
from Theorem~\ref{thm:allliftsisolated}.
\end{proof}

The following role rigidity is useful when a common factor is much smaller
than a prime in the shared additive-triple support.

\begin{proposition}[Large-prime endpoint-role rigidity]
\label{prop:largerole}
Let an irreducible polynomial \(h\) of degree \(e\) divide selected
primitive orientations \(Q_{A,B}\) and \(Q_{C,D}\).  Let \(p\geq5\) be
prime with \(p>e+2\), and suppose \(p^\kappa\Vert B\).  Then
\(p\nmid C+D\), and necessarily \(p^\lambda\Vert D\) for some
\(\lambda\geq1\).  Moreover
\begin{equation}
 \frac{\kappa}{A}=\frac{\lambda}{C},
 \qquad
 \frac{A}{\gcd(A,\kappa)}\mid e,
 \qquad
 \frac{C}{\gcd(C,\lambda)}\mid e.
 \tag{15.9}
\end{equation}
In particular, \(\kappa=\lambda=1\) is impossible for two distinct
packages.
\end{proposition}

\begin{proof}
Support compatibility, Corollary~\ref{cor:supportcompatibility}, forces
\(p\mid CD(C+D)\).  It cannot divide \(C+D\), because the radical sieve
would give \(e\equiv0\) or \(-2\pmod p\), impossible when
\(0<e<p-2\).

At an endpoint prime, every unit-cluster side has reduced slope denominator
at least \(p-2>e\).  A local factor of the degree-\(e\) polynomial \(h\)
cannot contain such a root.  Hence every root of \(h\) belongs to the
nonunit block.  In \(Q_{A,B}\) these valuations are positive and equal to
\(\kappa/A\); therefore the second orientation must also put \(p\) in its
constant endpoint, namely \(D\), rather than in its leading endpoint
\(C\), whose nonunit valuations have the opposite sign.  Equality of the
valuations gives the first relation in \((15.9)\).  Integrality of the
valuation of the norm gives the two denominator divisibilities.

If \(\kappa=\lambda=1\), then \(A\mid e\) and \(C\mid e\).  The two
nonunit blocks also give \(e\leq A,C\), so \(A=C=e\), contradicting the
disjointness of the two additive triples.
\end{proof}

The merged \(2\)-adic slope is therefore no longer an unresolved local
case.  The remaining bad-endpoint room is precise: an internal \(p\)-th
power must lie entirely in a nonunit endpoint block and requires an endpoint
valuation exponent at least \(p\), together with the divisibilities in
\((15.5)\)--\((15.6)\).

\section{Reciprocal resultants as a supplementary sieve}

The reciprocal resultant gives a compact numerical check on proposed factor sharing within a single package.  It is not strong enough to prove package coprimality by itself, but it supplies useful endpoint information and an independent consistency test.

For primitive \(a,b\), put \([k]=1+x+\cdots+x^{k-1}\).  Reciprocal
polynomials are normalized as in Section~\ref{sec:finite-support}.  One has
\begin{equation}
 Q_{a,b}(x)+Q_{b,a}(x)=(a+b)[a](x)[b](x).
 \tag{16.1}
\end{equation}

\begin{proposition}[Exact reciprocal resultant]
If \(\gcd(a,b)=1\), then
\begin{equation}
 \left|\Res(Q_{a,b},Q_{b,a})\right|
 =a^{a-2}b^{b-2}(a+b)^{a+b-2},
 \tag{16.2}
\end{equation}
where a factor \(1^{-1}\) is interpreted as \(1\).
\end{proposition}

\begin{proof}
At a root \(\alpha\) of \(Q_{a,b}\), equation \((16.1)\) gives
\[
 Q_{b,a}(\alpha)=(a+b)[a](\alpha)[b](\alpha).
\]
At a nontrivial \(a\)-th root of unity \(\zeta\),
\[
 Q_{a,b}(\zeta)=a\frac{\zeta^b-1}{(\zeta-1)^2}.
\]
Because multiplication by \(b\) permutes the nontrivial \(a\)-th roots,
\[
 |\Res(Q_{a,b},[a])|=a^{a-2}.
\]
Similarly, \(|\Res(Q_{a,b},[b])|=b^{b-2}\).  More explicitly,
\begin{align*}
 \Res(Q_{a,b},Q_{b,a})
 &=\Res\!\left(Q_{a,b},(a+b)[a][b]-Q_{a,b}\right)\\
 &=(a+b)^{a+b-2}
   \Res(Q_{a,b},[a])\Res(Q_{a,b},[b]).
\end{align*}
The preceding two evaluations give the formula.
\end{proof}

If \(a=gA\), \(b=gB\), \(\gcd(A,B)=1\), then composition and scaling give
\begin{equation}
 \left|\Res(Q_{a,b},Q_{b,a})\right|
 =g^{2g(A+B-2)}
 \left(A^{A-2}B^{B-2}(A+B)^{A+B-2}\right)^g.
 \tag{16.3}
\end{equation}
These formulas are useful as a sieve for a proposed common irreducible factor,
but cannot by themselves prove coprimality: a nonreciprocal irreducible
polynomial may have reciprocal resultant \(1\).

\section{Sparse Wronskians, explicit absolute boundedness, and arithmetic sieves}
\label{sec:wronskian}

The proof now turns the structural restrictions into absolute bounds.  Sparse Wronskians, torus-intersection rigidity, height estimates, and endpoint divisibility together force any normalized counterexample into an explicit finite range.

Let
\[
 W_{m,n;p,q}(x)
 =H'_{m,n}(x)H_{p,q}(x)-H_{m,n}(x)H'_{p,q}(x).
\]

\begin{proposition}[Sparse Wronskian]\label{prop:wronskian}
If a noncyclotomic irreducible polynomial \(h\) divides both
\(H_{m,n}\) and \(H_{p,q}\), then
\begin{equation}
 h(x)^2\mid \frac{W_{m,n;p,q}(x)}{(x-1)^4}.
 \tag{17.1}
\end{equation}
Moreover the numerator has at most eight monomials before equal exponents are
merged:
\begin{align}
 W_{m,n;p,q}(x)={}&mp(n-q)x^{n+q-1}
 +mq(p-n)x^{n+p-1} \notag\\
 &+np(q-m)x^{m+q-1}
 +nq(m-p)x^{m+p-1} \notag\\
 &+mn(q-p)(x^{n-1}-x^{m-1}) \notag\\
 &-pq(n-m)(x^{q-1}-x^{p-1}).
 \tag{17.2}
\end{align}
If the ordered edges are distinct, then \(W_{m,n;p,q}\not\equiv0\).
\end{proposition}

\begin{proof}
Write \(H_{m,n}=(x-1)^2K_1\) and
\(H_{p,q}=(x-1)^2K_2\).  Then
\[
 W_{m,n;p,q}=(x-1)^4(K_1'K_2-K_1K_2').
\]
If \(K_1=hA\) and \(K_2=hB\), the terms containing \(h'\) cancel:
\[
 K_1'K_2-K_1K_2'=h^2(A'B-AB').
\]
This proves \((17.1)\), and direct expansion gives \((17.2)\).
If the Wronskian vanished identically, then
\((H_{m,n}/H_{p,q})'=0\), so the two trinomials would be scalar multiples.
Comparison of their three nonzero supports \(\{0,m,n\}\) and
\(\{0,p,q\}\) would force \((m,n)=(p,q)\).
\end{proof}

The Wronskian supplies a square-factor obstruction for each fixed exponent
tuple.  A uniform structural reduction comes from applying the low-height
sparse-gcd theorem of Amoroso--Sombra--Zannier directly to the two collision
trinomials.  We first make the two geometric inputs completely formal.

\begin{lemma}[Four-unit divisor rigidity]\label{lem:fourunit}
Let \(C\) be a smooth projective curve over an algebraically closed field
\(K\) of characteristic zero, and let \(u,v\in K(C)\) be nonconstant.
If
\[
 \dim_{\Q}\operatorname{span}_{\Q}
 \{\operatorname{div}(u),\operatorname{div}(u-1),
   \operatorname{div}(v),\operatorname{div}(v-1)\}\leq2,
\]
then
\[
 v\in
 \left\{u,1-u,u^{-1},(1-u)^{-1},
          \frac{u}{u-1},\frac{u-1}{u}\right\}.
\]
\end{lemma}

\begin{proof}
Put
\[
 \Delta_0=\operatorname{div}(u),\quad
 \Delta_1=\operatorname{div}(u-1),\quad
 E_0=\operatorname{div}(v),\quad
 E_1=\operatorname{div}(v-1).
\]
A nonconstant function on a smooth projective curve is surjective onto
\(\mathbb P^1\).  Hence the fibres of \(u\) over \(0,1,\infty\) are all
nonempty.  The divisors \(\Delta_0,\Delta_1\) are linearly independent:
at a point over \(u=0\), only \(\Delta_0\) has positive valuation, and at
a point over \(u=1\), only \(\Delta_1\) does.  Thus
\[
 E_0=a\Delta_0+b\Delta_1,
 \qquad
 E_1=c\Delta_0+d\Delta_1
\]
for rational \(a,b,c,d\).  In particular, the zero, one, and pole fibres of
\(v\) are supported over the three special fibres of \(u\).  Moreover the
role is uniform on each such fibre.  Indeed, at a point \(P\) over \(u=0\),
\[
 \bigl(\ord_P(v),\ord_P(v-1)\bigr)
 =\ord_P(u)(a,c),
\]
while over \(u=1\) it is \(\ord_P(u-1)(b,d)\), and over \(u=\infty\) it is
\(-\ord_P(u^{-1})(a+b,c+d)\).  The coefficient pair is independent of the
chosen point in that fibre.

Because \(v-(v-1)=1\), the normalized valuation pair on each special fibre
is therefore uniformly one of
\[
 (r,0),\quad(0,r),\quad(-r,-r)\quad(r>0),
 \qquad\text{or}\qquad(0,0).
\]
Since the nonconstant function \(v\) assumes all three values
\(0,1,\infty\), the three nonneutral roles occur, and they occur
bijectively on the three special fibres of \(u\).  Replacing \(v\) by one
of the six Möbius transformations permuting \(\{0,1,\infty\}\) preserves
the same two-dimensional divisor span.  We may therefore arrange that the
zero, one, and pole fibres of \(v\) lie respectively over those of \(u\).
Then, for positive rationals \(r_0,r_1\),
\[
 \operatorname{div}(v)=r_0\operatorname{div}(u),
 \qquad
 \operatorname{div}(v-1)=r_1\operatorname{div}(u-1).
\]
At the common pole fibre, \(v\) and \(v-1\) have the same negative
valuation, so \(r_0=r_1=:r\).

Write \(r=A/B\) with \(A,B\) positive coprime integers.  The divisor
equalities give constants \(\kappa,\lambda\neq0\) such that
\[
 v^B=\kappa u^A,
 \qquad
 (v-1)^B=\lambda(u-1)^A.
\]
Logarithmic differentiation gives
\[
 B\frac{\dd v}{v}=A\frac{\dd u}{u},
 \qquad
 B\frac{\dd v}{v-1}=A\frac{\dd u}{u-1}.
\]
In characteristic zero, \(\dd u\not\equiv0\), and hence also
\(\dd v\not\equiv0\).  Comparing the two identities gives
\[
 \frac vu=\frac{v-1}{u-1},
\]
so \(v=u\).  Undoing the initial Möbius permutation gives the six
possibilities.
\end{proof}

\begin{proposition}[Rank-two torus-intersection rigidity]
\label{prop:ranktworigidity}
Let \(0<m<n\), \(0<p<q\), assume
\[
 \gcd\bigl(\gcd(m,n),\gcd(p,q)\bigr)=1,
\]
and suppose the two unordered packages are distinct.  Put
\[
 X=V\bigl((n-m)-nx_1+mx_2,\ (q-p)-qx_3+px_4\bigr)
 \subset\mathbb G_m^4.
\]
There is no geometric irreducible curve contained in the intersection of
\(X\) with a two-dimensional subtorus and containing a point
\[
 (\xi^m,\xi^n,\xi^p,\xi^q)
\]
for which \(\xi\notin\mu_\infty\) is a common root of
\(H_{m,n}\) and \(H_{p,q}\).
\end{proposition}

\begin{proof}
Base-change to an algebraic closure and work on the smooth projective model
of the normalization of such a curve \(C\).  Introduce
\[
 u=\frac{n}{n-m}x_1,
 \qquad u-1=\frac{m}{n-m}x_2,
\]
\[
 v=\frac{q}{q-p}x_3,
 \qquad v-1=\frac{p}{q-p}x_4.
\]

Suppose first that one projection of \(C\) to a punctured factor line is
constant.  The other projection is then nonconstant and hence dominant.  A
monomial \(x_3^a x_4^b\) constant on the second line must have
\(a=b=0\): parameterize the line by
\[
 x_3=t,
 \qquad
 x_4=\frac{qt-(q-p)}p,
\]
and compare the valuations at \(t=0\), \(t=(q-p)/q\), and infinity.
Thus, if the first projection is constant, the two independent characters
of the containing subtorus have zero last two coordinates.  They remain
independent in the first two coordinates, so the fixed point in
\(\mathbb G_m^2\) lies in a finite subgroup and is torsion.  At the
specified point, both \(\xi^m\) and \(\xi^n\) are roots of unity, forcing
\(\xi\) to be a root of unity.  This is a contradiction.  The other
constant projection is symmetric.

Both \(u\) and \(v\) are therefore nonconstant.  The annihilator lattice
of the two-dimensional subtorus has rank two.  Its two independent
characters lie in the kernel of the divisor map
\[
 \Q^4\longrightarrow\operatorname{Div}(C)\otimes\Q
\]
sending the standard basis to
\(\operatorname{div}(u),\operatorname{div}(u-1),
\operatorname{div}(v),\operatorname{div}(v-1)\).  Hence the four divisors
span a space of dimension at most two.  Lemma~\ref{lem:fourunit} shows that
\(v=\sigma(u)\) for one of its six Möbius transformations.

At the point coming from \(\xi\), this gives two binomial equations
\[
 \xi^{p-r_\sigma}=c_1,
 \qquad
 \xi^{q-s_\sigma}=c_2,
 \qquad c_1,c_2\in\Q^\times,
\]
where
\[
\begin{array}{c|cc}
 \sigma(u)&r_\sigma&s_\sigma\\ \hline
 u&m&n\\
 1-u&n&m\\
 u^{-1}&-m&n-m\\
 (1-u)^{-1}&-n&m-n\\
 u/(u-1)&m-n&-n\\
 (u-1)/u&n-m&-m
\end{array}
\]
The two exponent differences cannot both vanish unless the ordered edges,
and hence the packages, agree.  Put
\[
 d=\gcd\bigl(|p-r_\sigma|,|q-s_\sigma|\bigr)>0.
\]
Bézout gives \(\xi^d=c\in\Q^\times\).  If a prime divided both \(d\)
and \(\gcd(m,n)\), it would divide \(r_\sigma,s_\sigma,p,q\), contrary
to the coprimality hypothesis.  Thus
\[
 \gcd\bigl(d,\gcd(m,n)\bigr)=1.
\]
The minimal polynomial of \(\xi\) divides both \(H_{m,n}\) and
\(x^d-c\).  Lemma~\ref{lem:schinzelfibre} gives degree at most one, and
Theorem~\ref{thm:rational} then permits only the two orientations of the
single package \(\{1,2\}\).  Two distinct packages are impossible.
\end{proof}

We now state the sparse-gcd handoff in precisely the form used below.  For a
torus homomorphism \(\psi\), write \(\psi^\#\) for pullback on Laurent
polynomials.

\begin{proposition}[Specialized Amoroso--Sombra--Zannier output]
\label{prop:aszspecial}
Let
\[
 \varphi(t)=(t^m,t^n,t^p,t^q),
 \qquad
 D=\max\{m,n,p,q\},
\]
where \(\gcd(m,n,p,q)=1\), and let \(\xi\notin\mu_\infty\) be a common
root of \(H_{m,n}\) and \(H_{p,q}\).  Put
\[
 L_1=(n-m)-nx_1+mx_2,
 \qquad
 L_2=(q-p)-qx_3+px_4.
\]
Let \(B_4=B'(4,1)\) be the effective constant in Theorem~2.6 of
\cite{AmorosoSombraZannier2017}.  If
\begin{equation}
 D^{1/6}>B_4\bigl(1+\log(2D)\bigr),
 \tag{17.3}
\end{equation}
then there are \(0\leq k\leq3\), an injective torus homomorphism
\[
 \psi:\mathbb G_m^{4-k}\longrightarrow\mathbb G_m^4
\]
of size at most \(B_4\), and
\(\varphi_1:\mathbb G_m\to\mathbb G_m^{4-k}\), such that
\(\varphi=\psi\circ\varphi_1\).  Put
\[
 F_i=\psi^\#L_i,
 \qquad y_0=\varphi_1(\xi).
\]
Either both \(F_i\) vanish identically, or, with
\[
 G=\gcd(F_1,F_2),
 \qquad g=\varphi_1^\#G,
\]
one has \(G(y_0)=0\).
\end{proposition}

\begin{proof}
The coefficient height is at most \(\log(2D)\), and the exponent vector is
primitive, so Theorem~2.6 applies under \((17.3)\).  It gives the displayed
factorization and says that \(g\mid\gcd(H_{m,n},H_{p,q})\).  Moreover, a
root of the quotient of this gcd by \(g\) is either torsion or is a common
root of the two proper-subsum equations obtained by retaining a nonempty
proper subset \(\Lambda\subset\{1,2,3,4\}\) of the four nonconstant
monomials, while retaining the constant terms.

The point \(\xi\) is not torsion.  Let
\(I=\Lambda\cap\{1,2\}\).  If \(I\) is empty, the first partial equation
is the nonzero constant \(n-m\).  If \(I=\{1\}\), it is
\((n-m)-n\xi^m=0\); subtracting it from the full first equation gives
\(m\xi^n=0\), impossible.  If \(I=\{2\}\), subtraction gives
\(-n\xi^m=0\), also impossible.  Thus \(I=\{1,2\}\).  The second
equation similarly forces
\(\Lambda\cap\{3,4\}=\{3,4\}\).  Hence \(\Lambda\) would be the full
set, contradicting properness.  Therefore \(g(\xi)=0\), which is exactly
\(G(y_0)=0\).  The possibility that both pulled-back forms vanish
identically is recorded separately.
\end{proof}

\begin{theorem}[Effective absolute boundedness of normalized counterexamples]
\label{thm:hyperplane}
There is an effectively computable absolute constant \(D_0\) such that a
normalized four-index counterexample
\[
 0<m<n,\qquad0<p<q,\qquad\gcd(m,n,p,q)=1
\]
must satisfy
\[
 \max\{m,n,p,q\}\leq D_0.
\]
More explicitly, with \(B_4\) as in
Proposition~\ref{prop:aszspecial}, define
\begin{equation}
 D_{\mathrm{ASZ}}=
 \max\!\left(
 \{1\}\cup
 \left\{D\in\mathbb N:
 D^{1/6}\leq B_4(1+\log(2D))\right\}
 \right),
 \qquad
 D_0=\max\{\lceil B_4\rceil,D_{\mathrm{ASZ}}\}.
 \tag{17.4}
\end{equation}
Then this \(D_0\) has the stated property.
\end{theorem}

\begin{proof}
Assume \(D>D_0\) and use Proposition~\ref{prop:aszspecial}.  The two
pulled-back forms cannot both vanish identically.  Indeed, if
\[
 (n-m)-n\chi_1+m\chi_2=0
\]
as a Laurent polynomial in characters of a positive-dimensional torus,
linear independence of distinct characters forces
\(\chi_1=\chi_2=1\).  The second form similarly forces
\(\chi_3=\chi_4=1\), making \(\psi\) trivial, contrary to injectivity.
Thus \(G(y_0)=0\).  Choose a geometric irreducible factor \(P\) of \(G\)
vanishing at \(y_0\); if exactly one \(F_i\) is zero, take \(P\) from the
other one.

If \(k=3\), write the primitive exponent column of \(\psi\) as
\(B=(b_1,b_2,b_3,b_4)^{\mathsf T}\), so
\[
 (m,n,p,q)^{\mathsf T}=B\theta.
\]
Injectivity gives \(\gcd(b_1,b_2,b_3,b_4)=1\), while normalization gives
\(\gcd(m,n,p,q)=1\).  Hence \(|\theta|=1\) and
\(D=\|B\|_\infty\leq B_4\), contradicting \(D>D_0\).

If \(k=0\), the injective endomorphism \(\psi\) of a torus of the same
dimension is an automorphism.  Pullback therefore preserves coprimality,
whereas \(L_1,L_2\) are coprime Laurent linear forms.  Thus \(G\) is a
unit, contradicting \(G(y_0)=0\).

Suppose \(k=1\).  The hypersurface \(V(P)\subset\mathbb G_m^3\) is a
geometric irreducible surface.  An injective homomorphism of tori is a closed
immersion, so \(\psi(V(P))\) is a closed irreducible surface contained in
\[
 X=V(L_1,L_2),
\]
the irreducible product of two punctured lines.  Since \(\dim X=2\), the
image must equal \(X\), placing \(X\) inside the proper subtorus
\(\operatorname{im}\psi\).

This is impossible even for a mixed monomial relation.  If
\[
 x_1^a x_2^b x_3^c x_4^d
\]
were constant on the product, fixing the second factor would make
\(x_1^a x_2^b\) constant on the first line, and fixing the first would do
the same for \(x_3^c x_4^d\) on the second.  Parameterizing the first line
by
\[
 x_1=t,
 \qquad x_2=\frac{nt-(n-m)}m
\]
and comparing the valuations at \(0,(n-m)/n,\infty\) gives \(a=b=0\);
symmetrically \(c=d=0\).  Thus \(X\) lies in no proper subtorus.

Finally let \(k=2\).  The geometric curve \(V(P)\subset\mathbb G_m^2\)
has closed image contained in
\(X\cap\operatorname{im}\psi\) and passing through the specified point.
The global normalization gives
\[
 \gcd\bigl(\gcd(m,n),\gcd(p,q)\bigr)=1.
\]
Proposition~\ref{prop:ranktworigidity} excludes this curve.  Every possible
value of \(k\) is contradictory, proving the bound.
\end{proof}

\begin{corollary}[Finite-template reduction]\label{cor:finiteglobal}
Every counterexample is a simultaneous positive integral dilation of one of
finitely many normalized configurations.  In particular, the full package
coprimality conjecture is reduced to an effective finite exact computation.
\end{corollary}

\begin{proof}
Divide the four edge endpoints by their common gcd and replace the common
root \(z\) by the corresponding power, as in \((2.6)\).  The normalized
configuration is bounded by Theorem~\ref{thm:hyperplane}.  Conversely, the
original configuration is its simultaneous dilation.
\end{proof}

\begin{remark}[The symbolic ASZ threshold]
The constant \(B_4\) is effective in the cited theorem, but no numerical
value is supplied there.  Thus Theorem~\ref{thm:hyperplane}, by itself, does
not furnish a practical endpoint.  The direct argument in the next subsection
bypasses \(B_4\) and produces an explicit numerical bound.
\end{remark}

\subsection{A direct numerical bound from height, geometry of numbers, and Smyth's theorem}
\label{subsec:explicitbound}

The preceding Amoroso--Sombra--Zannier argument proves effective absolute
boundedness without displaying a numerical constant.  In the present
four-variable situation one can obtain a substantially more concrete bound
by combining an elementary height estimate with a short rank-two subtorus.
We write \(\mathrm h(\alpha)\) for the absolute logarithmic Weil height and
\(M(f)\) for the Mahler measure of a primitive polynomial \(f\in\Z[x]\).

\begin{lemma}[Height of a collision root]\label{lem:collisionheight}
Let \(0<u<v\), and suppose \(H_{u,v}(\alpha)=0\).  Then
\begin{equation}
 (v-u)\,\mathrm h(\alpha)\leq\log(2v),
 \qquad
 u\,\mathrm h(\alpha)\leq\log(2v).
 \tag{17.5}
\end{equation}
Consequently,
\begin{equation}
 v\,\mathrm h(\alpha)\leq2\log(2v).
 \tag{17.6}
\end{equation}
\end{lemma}

\begin{proof}
We use the elementary affine-height inequality
\[
 \mathrm h\!\left(\frac{a\beta+b}{c}\right)
 \leq \mathrm h(\beta)
      +\log\!\bigl(2\max\{|a|,|b|,|c|\}\bigr),
 \qquad a,b,c\in\Z,\ c\neq0.
\]
It follows place by place from the triangle inequality, with the factor two
needed only at the archimedean places.  The collision equation gives
\[
 \alpha^v=\frac{v\alpha^u-(v-u)}u,
\]
and hence
\[
 v\,\mathrm h(\alpha)
 \leq u\,\mathrm h(\alpha)+\log(2v),
\]
which is the first inequality in \((17.5)\).

The reciprocal identity
\[
 x^vH_{v-u,v}(x^{-1})=H_{u,v}(x)
\]
shows that \(\alpha^{-1}\) is a root of \(H_{v-u,v}\).  Applying the first
inequality to that edge and using
\(\mathrm h(\alpha^{-1})=\mathrm h(\alpha)\) gives the second inequality.
Adding them proves \((17.6)\).
\end{proof}

For a Euclidean lattice \(\Lambda\), write \(\det\Lambda\) for its
covolume in its real span.

\begin{lemma}[Degree of a rank-two subtorus]\label{lem:toricdegree}
Let \(R\subset\Z^4\) be a saturated rank-two lattice, and let
\(T\subset\mathbb G_m^4\) be the two-dimensional subtorus annihilated by
\(R\).  Under the standard embedding
\[
 \mathbb G_m^4\longrightarrow\mathbb P^4,
 \qquad
 (x_1,x_2,x_3,x_4)\longmapsto[1:x_1:x_2:x_3:x_4],
\]
the projective closure \(\overline T\) satisfies
\begin{equation}
 \deg\overline T\leq2\det R.
 \tag{17.7}
\end{equation}
\end{lemma}

\begin{proof}
Put \(M=\Z^4/R\).  Saturation makes \(M\) a free rank-two lattice.  Choose
a lattice basis of \(M\), and let \(u_1,\ldots,u_4\in\Z^2\) be the images
of the four standard basis vectors.  With
\[
 \mathcal A=\{0,u_1,u_2,u_3,u_4\},
 \qquad
 P=\operatorname{conv}(\mathcal A),
\]
the closure \(\overline T\) is the projective toric surface
\(X_{\mathcal A}\).  The vectors \(u_i\) generate \(M\), so there is no
lattice-index correction.  The normalization of \(X_{\mathcal A}\) is the
normal toric surface associated with \(P\), and finite birational
normalization preserves the hyperplane degree.  Hence the standard
degree--volume theorem gives
\[
 \deg\overline T=\operatorname{Vol}_M(P),
\]
where \(\operatorname{Vol}_M\) is normalized lattice area: it is twice
ordinary Euclidean area in coordinates in which a fundamental
parallelogram of \(M\) has area one
\cite[Theorems~3.A.5 and~13.4.3]{CoxLittleSchenck2011}.

Let \(U\) be the \(2\times4\) integer matrix with columns \(u_i\), and let
\(S\subset\Z^4\) be its row lattice.  The quotient map
\(U:\Z^4\to\Z^2\) is surjective and has kernel \(R\).  Thus \(S\) is
primitive; since its real span is \(R^\perp\),
\[
 S=R^\perp\cap\Z^4.
\]
For complementary primitive rank-two lattices in the unimodular lattice
\(\Z^4\), the covolumes agree.  Concretely, the Euclidean Hodge star sends
a primitive generator of \(\bigwedge^2R\) to a primitive generator of
\(\bigwedge^2S\) and preserves norm.  Hence \(\det S=\det R\).  The two
rows of \(U\) form a basis of \(S\), so Cauchy--Binet gives
\begin{equation}
 (\det R)^2=(\det S)^2=\det(UU^{\mathsf T})
 =\sum_{1\leq i<j\leq4}\det(u_i,u_j)^2.
 \tag{17.8}
\end{equation}

Order the vertices of \(P\) cyclically.  In the shoelace sum, every edge
incident with the vertex \(0\) contributes zero.  All other vertices belong
to \(\{u_1,\ldots,u_4\}\), so at most four nonzero determinant terms
remain.  Therefore Cauchy--Schwarz and \((17.8)\) give
\[
 \deg\overline T
 \leq2\left(\sum_{i<j}\det(u_i,u_j)^2\right)^{1/2}
 =2\det R.
\]
\end{proof}

\begin{proposition}[A short rank-two subtorus]\label{prop:shorttorus}
Let
\[
 \boldsymbol a=(m,n,p,q)\in\Z^4
\]
be primitive, and put \(D=\|\boldsymbol a\|_\infty\).  There is a
two-dimensional subtorus \(T\subset\mathbb G_m^4\) containing the
one-parameter subgroup
\[
 t\longmapsto(t^m,t^n,t^p,t^q)
\]
and satisfying
\begin{equation}
 \deg\overline T<4D^{2/3}.
 \tag{17.9}
\end{equation}
\end{proposition}

\begin{proof}
Let
\[
 L=\boldsymbol a^\perp\cap\Z^4.
\]
Primitivity of \(\boldsymbol a\) means that the homomorphism
\(z\mapsto\langle\boldsymbol a,z\rangle\) maps \(\Z^4\) onto \(\Z\).
Thus \(L\) is saturated of rank three.  Choose \(c\in\Z^4\) with
\(\langle\boldsymbol a,c\rangle=1\).  A basis of \(L\), together with
\(c\), is a basis of \(\Z^4\); comparing base volume with the height
\(1/\|\boldsymbol a\|_2\) above \(\boldsymbol a^\perp\) gives
\[
 \det L=\|\boldsymbol a\|_2\leq2D.
\]

Its Euclidean dual lattice \(L^\vee\) has determinant
\((\det L)^{-1}\).  Minkowski's first theorem for the three-dimensional
Euclidean ball \cite[Chapter~III, \S2]{Cassels1997} provides a nonzero
\(y\in L^\vee\) with
\[
 \|y\|_2\leq
 \left(\frac6\pi\right)^{1/3}(\det L)^{-1/3}.
\]
Divide by the content of \(y\) in the free abelian group \(L^\vee\).
This can only decrease its norm, and it makes \(y\) primitive, meaning that
\[
 L\longrightarrow\Z,
 \qquad b\longmapsto\langle b,y\rangle
\]
is surjective.  Set
\[
 R=\{b\in L:\langle b,y\rangle=0\}.
\]
Then \(L/R\simeq\Z\), so \(R\) is saturated in \(L\).  It is also
saturated in \(\Z^4\): in the exact sequence
\[
 0\longrightarrow L/R\longrightarrow\Z^4/R
 \longrightarrow\Z^4/L\longrightarrow0
\]
both outer groups are torsion-free, hence so is \(\Z^4/R\).

Choose a basis \(b_1,b_2\) of \(R\) and \(b_3\in L\) with
\(\langle b_3,y\rangle=1\).  The Euclidean height of \(b_3\) above
\(R\otimes_{\Z}\mathbb R\) is \(1/\|y\|_2\).  Therefore
\[
 \det L=\frac{\det R}{\|y\|_2},
 \qquad
 \det R=(\det L)\|y\|_2
 \leq\left(\frac6\pi\right)^{1/3}(2D)^{2/3}.
\]
Let \(T\) be the subtorus annihilated by \(R\).  Since
\(R\subset\boldsymbol a^\perp\), every character in \(R\) is trivial on
the displayed one-parameter subgroup; hence \(T\) contains that subgroup.
Lemma~\ref{lem:toricdegree} gives
\[
 \deg\overline T
 \leq2\left(\frac6\pi\right)^{1/3}(2D)^{2/3}
 =2\left(\frac{24}{\pi}\right)^{1/3}D^{2/3}
 <4D^{2/3},
\]
because \(\pi>3\).
\end{proof}

\begin{proposition}[Degree of a normalized common-root field]
\label{prop:fielddegree}
Let
\[
 0<m<n,\qquad0<p<q,\qquad\gcd(m,n,p,q)=1,
\]
let the two packages be distinct, and suppose that
\(\alpha\notin\mu_\infty\) is a common root of
\(H_{m,n}\) and \(H_{p,q}\).  Put
\[
 D=\max\{m,n,p,q\},
 \qquad d=[\Q(\alpha):\Q].
\]
Then
\begin{equation}
 d<4D^{2/3}.
 \tag{17.10}
\end{equation}
\end{proposition}

\begin{proof}
Take the subtorus \(T\) from Proposition~\ref{prop:shorttorus}.  It contains
\[
 P_\alpha=(\alpha^m,\alpha^n,\alpha^p,\alpha^q).
\]
The variety cut out in \(\mathbb G_m^4\) by the two collision lines is the
variety \(X\) of Proposition~\ref{prop:ranktworigidity}.  That proposition
shows that \(P_\alpha\) is isolated in \(T\cap X\): a positive-dimensional
local component would contain a geometric irreducible curve through the
point.  Since \(P_\alpha\) lies in the open torus, a positive-dimensional
component of \(\overline T\cap\overline X\) through it would restrict near
that point to one in \(T\cap X\).  Hence it is also an isolated projective
intersection point.

Both \(T\) and \(X\), and therefore their projective closures, are defined
over \(\Q\).  The closure \(\overline X\) is the intersection of the two
rational hyperplanes
\[
 (n-m)x_0-nx_1+mx_2=0,
 \qquad
 (q-p)x_0-qx_3+px_4=0
\]
in \(\mathbb P^4\), so \(\deg\overline X=1\).  The generalized Bezout
inequality for isolated components
\cite[\S8.4]{Fulton1998} gives
\[
 \deg\bigl((\overline T\cap\overline X)_{\mathrm{isol}}\bigr)
 \leq\deg\overline T\deg\overline X=\deg\overline T.
\]
Every Galois conjugate of \(P_\alpha\) is again isolated and contributes
positive intersection multiplicity.  The closed point determined by this
orbit therefore contributes at least \([\Q(P_\alpha):\Q]\), whence
\[
 [\Q(P_\alpha):\Q]\leq\deg\overline T<4D^{2/3}.
\]
Finally, choose integers \(r_1,r_2,r_3,r_4\) with
\(r_1m+r_2n+r_3p+r_4q=1\).  Since all four coordinates are nonzero,
\[
 \alpha=(\alpha^m)^{r_1}(\alpha^n)^{r_2}
         (\alpha^p)^{r_3}(\alpha^q)^{r_4}.
\]
Thus \(\Q(P_\alpha)=\Q(\alpha)\), proving \((17.10)\).
\end{proof}

Let
\[
 \theta_0=1.3247179572447460259\ldots
\]
be the plastic number, the real root greater than one of
\(T^3-T-1\).

\begin{proposition}[A uniform Mahler-measure gap]\label{prop:mahlergap}
Let \(h\in\Z[x]\) be the primitive irreducible polynomial, with positive
leading coefficient, of a noncyclotomic common root of two collision
trinomials.  Then \(h\) is nonreciprocal and
\begin{equation}
 M(h)\geq\theta_0.
 \tag{17.11}
\end{equation}
\end{proposition}

\begin{proof}
The polynomial \(h\) divides an orientation \(Q_{a,b}\).  If it were
self-reciprocal, it would also divide the reciprocal orientation
\(Q_{b,a}\), contradicting the nonzero reciprocal resultant \((16.3)\).
Thus \(h\) is nonreciprocal.

If \(h\) is monic, Smyth's theorem \cite{Smyth1971} gives
\(M(h)\geq\theta_0\).  If it is not monic, its positive integral leading
coefficient is at least two, and
\(M(h)\geq2>\theta_0\).
\end{proof}

\begin{theorem}[Explicit absolute bound for normalized counterexamples]
\label{thm:explicitbound}
Every normalized four-index counterexample
\[
 0<m<n,\qquad0<p<q,\qquad\gcd(m,n,p,q)=1
\]
satisfies
\begin{equation}
 \boxed{\max\{m,n,p,q\}\leq175\,394\,637.}
 \tag{17.12}
\end{equation}
\end{theorem}

\begin{proof}
Let \(\alpha\) be a common root, let \(h\) be its primitive irreducible
polynomial, let \(d=\deg h\), and put
\(D=\max\{m,n,p,q\}\).  One of the two collision edges has upper endpoint
\(D\).  Lemma~\ref{lem:collisionheight} gives
\begin{equation}
 \mathrm h(\alpha)\leq\frac{2\log(2D)}D.
 \tag{17.13}
\end{equation}
By Proposition~\ref{prop:fielddegree},
\(d<4D^{2/3}\).  Since
\(\log M(h)=d\,\mathrm h(\alpha)\),
Proposition~\ref{prop:mahlergap} yields
\begin{equation}
 \log\theta_0
 \leq\log M(h)
 <\frac{8\log(2D)}{D^{1/3}}.
 \tag{17.14}
\end{equation}
The function on the right is strictly decreasing for \(D\geq11\).  Certified
real-ball arithmetic gives
\[
 \frac{8\log(2\cdot175394637)}{175394637^{1/3}}
   >\log\theta_0,
\]
whereas
\[
 \frac{8\log(2\cdot175394638)}{175394638^{1/3}}
   <\log\theta_0.
\]
Thus \((17.14)\) is impossible for \(D\geq175394638\).  The smaller values
\(D<11\) are already below the asserted bound.
\end{proof}

\begin{corollary}[The upper range is unit-normalized]\label{cor:unitrange}
In a normalized counterexample with
\[
 D=\max\{m,n,p,q\}\geq6\,816\,242,
\]
the primitive irreducible polynomial \(h\) of a common root is monic and has
constant coefficient \(\pm1\).
\end{corollary}

\begin{proof}
The same decreasing function satisfies
\[
 \frac{8\log(2\cdot6816242)}{6816242^{1/3}}<\log2.
\]
Hence \((17.13)\) and Proposition~\ref{prop:fielddegree} give
\(M(h)<2\).  Both the absolute leading coefficient and the absolute constant
coefficient of a primitive polynomial are at most its Mahler measure.  They
are nonzero integers, so both equal one.
\end{proof}

\begin{remark}
Theorem~\ref{thm:explicitbound} removes the practical indeterminacy in the
constant \(B_4\).  The Amoroso--Sombra--Zannier route remains an independent
structural proof of effective boundedness, while the direct argument gives a
fully numerical endpoint.  The two tiny numerical comparisons above can be
certified by the real-ball script recorded in
Section~\ref{sec:computation}.
\end{remark}

\subsection{Moment, discriminant, and endpoint sieves}
\label{subsec:momentsieves}

The explicit box is still large.  The following additional restrictions are
intended for the finite computation.  They are extracted from the all-moment
form of Borisov's cluster-count argument.

Let
\[
 h(x)=ux^e+c_{e-1}x^{e-1}+\cdots+c_1x+v\in\Z[x],
 \qquad u>0,
\]
be primitive and irreducible.  For its roots \(\beta\), counted with
multiplicity, put
\begin{equation}
 A_j(h)=\sum_{h(\beta)=0}(1-\beta)\beta^j
 \qquad(j\geq0).
 \tag{17.15}
\end{equation}

\begin{proposition}[All-moment congruence]\label{prop:momentcongruence}
Suppose \(h\) divides two selected primitive orientations
\(Q_{A,B}\) and \(Q_{C,D}\), where
\[
 (A,B)=(C,D)=1,
 \qquad N=A+B,
 \qquad M=C+D.
\]
If a prime \(\ell\) divides \(ABNCDM\) but does not divide \(uv\), then
\begin{equation}
 A_j(h)\in\ell\Z_{(\ell)}
 \qquad\text{for every }j\geq0.
 \tag{17.16}
\end{equation}
\end{proposition}

\begin{proof}
Use a package whose additive triple contains \(\ell\).  Since
\(\ell\nmid uv\), the leading and constant coefficients of \(h\) are
\(\ell\)-adic units.  The Newton polygon of \(h\) is therefore horizontal,
so all roots of \(h\) are \(\ell\)-adic units.

Borisov's cluster argument \cite[Lemma~3.1 and Remark~3.5]{Borisov1998}
shows that, at every nontrivial prime-to-\(\ell\) root-of-unity center, the
number of roots belonging to a rational factor is divisible by \(\ell\).
Modulo \(\ell\), the summand in \((17.15)\) is constant on such a cluster,
namely \((1-\zeta)\zeta^j\), so the cluster contribution vanishes.  At the
center \(1\), every summand already vanishes modulo \(\ell\).  If \(\ell\)
is an endpoint prime, the nonunit endpoint block is absent because every root
of \(h\) is a unit.  These contributions exhaust the roots and prove
\((17.16)\).  The argument also covers \(\ell=2\); the merged side still
contributes an even number of roots to every rational factor, as in
Proposition~\ref{prop:mergedtwo}.
\end{proof}

Define the reciprocal coefficient polynomial
\[
 \widetilde h(t)=t^eh(t^{-1})
 =u+c_{e-1}t+\cdots+c_1t^{e-1}+vt^e.
\]

\begin{corollary}[Differential moment congruence]
\label{cor:momentdifferential}
Under the hypotheses of Proposition~\ref{prop:momentcongruence}, put
\[
 \mathcal S=\operatorname{rad}(ABNCDM),
 \qquad
 \mathcal R=\prod_{\substack{\ell\mid\mathcal S\\ \ell\nmid uv}}\ell.
\]
Then every coefficient of
\begin{equation}
 (1-t)\widetilde h'(t)+e\widetilde h(t)
 \tag{17.17}
\end{equation}
is divisible by \(\mathcal R\).  In particular,
\begin{equation}
 \mathcal R\mid eu+c_{e-1},
 \qquad
 \mathcal R\mid ev+c_1.
 \tag{17.18}
\end{equation}
\end{corollary}

\begin{proof}
For each prime in \(\mathcal R\), the constant coefficient
\(\widetilde h(0)=u\) is a unit.  The generating series of the moments is
\[
 \sum_{j\geq0}A_j(h)t^j
 =e+(1-t)\frac{\widetilde h'(t)}{\widetilde h(t)}.
\]
Proposition~\ref{prop:momentcongruence} makes the left side zero modulo the
prime.  Multiplication by \(\widetilde h(t)\) gives \((17.17)\), and the
constant and \(t^{e-1}\) coefficients give \((17.18)\).
\end{proof}

\begin{proposition}[Hankel--discriminant divisibility]
\label{prop:hankeldisc}
With the same notation,
\begin{equation}
 \boxed{\mathcal R^e\mid h(1)\Disc(h).}
 \tag{17.19}
\end{equation}
Consequently,
\begin{equation}
 \mathcal R\leq2eM(h)^{2-1/e}.
 \tag{17.20}
\end{equation}
\end{proposition}

\begin{proof}
Form the \(e\times e\) Hankel matrix
\[
 \mathcal H=(A_{i+j}(h))_{0\leq i,j<e}.
\]
Writing \(w_k=1-\beta_k\), its Vandermonde factorization is
\[
 \mathcal H=V\,\operatorname{diag}(w_1,\ldots,w_e)V^{\mathsf T}.
\]
Hence, up to sign,
\begin{equation}
 \det\mathcal H
 =\prod_k(1-\beta_k)
  \prod_{i<j}(\beta_i-\beta_j)^2
 =\frac{h(1)\Disc(h)}{u^{2e-1}}.
 \tag{17.21}
\end{equation}
Every entry is divisible by every prime in \(\mathcal R\), while \(u\) is a
unit at those primes.  Thus \((17.19)\) follows.

The elementary estimates
\[
 |h(1)|\leq2^eM(h),
 \qquad
 |\Disc(h)|\leq e^eM(h)^{2e-2}
\]
(the second is Mahler's discriminant inequality \cite{Mahler1964}) imply
\[
 \mathcal R^e
 \leq2^ee^eM(h)^{2e-1}.
\]
Taking \(e\)-th roots proves \((17.20)\).
\end{proof}

\begin{corollary}[Support bound for two colliding primitive packages]
\label{cor:supportbound}
Under the same hypotheses,
\begin{equation}
 \mathcal S\leq2eM(h)^{4-1/e}.
 \tag{17.22}
\end{equation}
If \(h\) is monic with constant coefficient \(\pm1\), then
\begin{equation}
 \mathcal S\leq2eM(h)^{2-1/e}.
 \tag{17.23}
\end{equation}
\end{corollary}

\begin{proof}
The primes of \(\mathcal S\) omitted from \(\mathcal R\) divide \(uv\).
Since
\[
 \operatorname{rad}(uv)\leq|uv|\leq M(h)^2,
\]
formula \((17.20)\) gives \((17.22)\).  In the unit-normalized case no prime
is omitted, so \(\mathcal S=\mathcal R\), giving \((17.23)\).
\end{proof}

\begin{proposition}[Endpoint divisibilities and slope matching]
\label{prop:endpointarith}
Under the same hypotheses,
\begin{equation}
 |v|^A\mid B^e,
 \qquad u^B\mid A^e,
 \qquad |v|^C\mid D^e,
 \qquad u^D\mid C^e.
 \tag{17.24}
\end{equation}
Moreover, for every prime \(\ell\mid u\),
\begin{equation}
 \frac{v_\ell(A)}B=\frac{v_\ell(C)}D,
 \tag{17.25}
\end{equation}
and, for every prime \(\ell\mid v\),
\begin{equation}
 \frac{v_\ell(B)}A=\frac{v_\ell(D)}C.
 \tag{17.26}
\end{equation}
\end{proposition}

\begin{proof}
Because \(Q_{A,B}\) is primitive with leading coefficient \(A\) and
constant coefficient \(B\), Gauss's lemma gives an integral complementary
factor.  Hence \(u\mid A\), \(|v|\mid B\), and \((u,v)=1\).
At a root \(\beta\) of \(h\),
\[
 \beta^A(N-A\beta^B)=B.
\]
Taking the resultant with \(h\) gives
\begin{equation}
 \Res(h,N-Ax^B)
 =\pm\frac{B^eu^N}{v^A}\in\Z.
 \tag{17.27}
\end{equation}
Since \((u,v)=1\), this proves \(|v|^A\mid B^e\).  Applying the same
argument to the reciprocal factor in \(Q_{B,A}\) proves
\(u^B\mid A^e\).  The second package gives the remaining two divisibilities.

Now let \(\ell\mid u\).  Then \(\ell\mid A,C\) and \(\ell\nmid v\).
The product of the roots of \(h\) has negative \(\ell\)-adic valuation, so
at least one root is nonunit.  Borisov's endpoint polygons say that every
negative-valuation root of \(Q_{A,B}\) has valuation
\(-v_\ell(A)/B\), and every such root of \(Q_{C,D}\) has valuation
\(-v_\ell(C)/D\).  The same root lies in both polynomials, proving
\((17.25)\).  The positive-valuation endpoint description gives
\((17.26)\) in exactly the same way when \(\ell\mid v\).
\end{proof}

\begin{remark}
For a finite search, the support bound, the four endpoint divisibilities, and
the slope equalities should be applied before any polynomial factorization.
For a surviving exponent tuple, let \(S\in\Z[x]\) be the primitive part of
\(W_{m,n;p,q}/(x-1)^4\).  Proposition~\ref{prop:wronskian} then permits a
cheap modular certificate: if an auxiliary prime does not divide the leading
coefficient of \(S\)---equivalently, reduction preserves its degree---and
\(\overline S\) is square-free over the corresponding finite field, then
\(S\) is square-free over \(\Q\), so no common irreducible factor can exist.
The good-reduction condition is essential: at a prime where the degree drops,
a square factor can disappear after reduction.
\end{remark}


\section{Exact upper-range elimination and reproducibility}
\label{sec:computation}

This is the first decisive computational section.  All filters used here are exact integer filters derived above; the computation does not rely on numerical root finding or on factoring the large collision polynomials.

Put
\[
 D_0=6\,816\,242,
 \qquad
 D_1=175\,394\,637.
\]
This section gives the mathematical reductions and the exact computation
which exclude every normalized counterexample with \(D_0\leq D\leq D_1\).
Write the two unordered packages in primitive form as
\[
 \{rA,rB\},\qquad \{sC,sE\},
 \qquad
 \gcd(A,B)=\gcd(C,E)=1,
\]
put \(N=A+B\), \(M=C+E\), and order the shapes by \(A<B\), \(C<E\).
The corresponding collision upper endpoints are \(rN\) and \(sM\), so
\[
 D=\max(rN,sM).
\]
Normalization gives \(\gcd(r,s)=1\).  Let \(\alpha\) be a common root,
let \(h_\alpha\) be its primitive minimal polynomial, and put
\[
 d=[\Q(\alpha):\Q],\qquad
 \beta=\alpha^r,\qquad \gamma=\alpha^s.
\]
The elements \(\beta\) and \(\gamma\) are roots of selected reciprocal
orientations of the two primitive shapes.

\subsection{An exact uniform degree gap and small scales}

\begin{lemma}[An algebraic certificate for the plastic number]
\label{lem:plasticpower}
Let \(\theta_0\) be the real root greater than one of \(T^3-T-1\).  Then
\begin{equation}
 \theta_0^{70}
 =62\,608\,681+109\,870\,576\theta_0
   +82\,938\,844\theta_0^2
 >2D_1.
 \tag{18.1}
\end{equation}
\end{lemma}

\begin{proof}
Reduction of \(T^{70}\) modulo \(T^3-T-1\) gives the displayed identity.
The polynomial \(T^3-T-1\) is strictly increasing for \(T\geq1\), and
\[
 \left(\frac{33}{25}\right)^3-\frac{33}{25}-1
 =-\frac{313}{15625}<0.
\]
Hence \(\theta_0>33/25\).  Substitution in the positive-coefficient
right-hand side of \((18.1)\) gives the exact rational lower bound
\[
 \theta_0^{70}>
 \frac{220\,094\,051\,941}{625}
 >350\,789\,274=2D_1.
\]
\end{proof}

\begin{proposition}[Uniform degree and scale bounds]
\label{prop:uniformdegree}
For a normalized counterexample with \(D\leq D_1\),
\begin{equation}
 d>\frac D{140}.
 \tag{18.2}
\end{equation}
In the primitive decomposition above,
\begin{equation}
 1\leq r,s\leq139.
 \tag{18.3}
\end{equation}
In particular, throughout the unit range \(D\geq D_0\),
\[
 d\geq\left\lfloor\frac{D_0}{140}\right\rfloor+1=48\,688.
\]
\end{proposition}

\begin{proof}
The height estimate \((17.13)\) and the Mahler gap \((17.11)\) give
\[
 d\geq\frac{D\log\theta_0}{2\log(2D)}.
\]
Lemma~\ref{lem:plasticpower} gives
\(70\log\theta_0>\log(2D_1)\geq\log(2D)\), proving \((18.2)\).
The confluent degree bound, Theorem~\ref{thm:schinzelbound}, gives
\[
 d\leq\min(N,M)-2<N,M.
\]
Since \(rN\leq D\) and \(sM\leq D\), the inequalities
\(D/140<d<N\leq D/r\) and \(D/140<d<M\leq D/s\) imply
\(r,s<140\).
\end{proof}

\subsection{Full Kummer degree in the unit range}

The prime-degree form of Capelli's criterion used earlier has the following
standard composite-degree form: over a characteristic-zero field \(K\),
\(X^n-a\) is irreducible if \(a\notin K^p\) for every prime \(p\mid n\),
and, when \(4\mid n\), also \(a\notin-4K^4\); see
\cite[Chapter~VI, Section~9]{LangAlgebra}.

\begin{proposition}[Full Kummer degree]
\label{prop:unitkummer}
Assume \(D_0\leq D\leq D_1\).  Let \(e_1=[\Q(\beta):\Q]\) and
\(e_2=[\Q(\gamma):\Q]\).  Then
\[
 X^r-\beta\quad\text{is irreducible over }\Q(\beta),
 \qquad
 X^s-\gamma\quad\text{is irreducible over }\Q(\gamma).
\]
Consequently there is an integer \(k\geq1\) such that
\begin{equation}
 d=re_1=se_2=rsk,
 \qquad e_1=sk,
 \qquad e_2=rk.
 \tag{18.4}
\end{equation}
\end{proposition}

\begin{proof}
Corollary~\ref{cor:unitrange} makes \(\alpha\), and hence \(\beta,\gamma\),
algebraic units.  Fix a prime \(p\mid r\), put \(K=\Q(\beta)\), and suppose
first that \(\beta\in K^p\).

If \(p\mid N\), Corollary~\ref{cor:totalprimepower} gives an immediate
contradiction.  If \(p\mid AB\), the proof of
Theorem~\ref{thm:endpointpower} shows that an internal \(p\)-th power factor
must lie entirely in the nonunit endpoint block, again impossible for the
unit \(\beta\).  The only remaining case is \(p\nmid ABN\).  Then \(p\) is
odd, and Theorem~\ref{thm:frobeniusfactor} gives \(e_1\leq p-2\).  Since
\(\alpha\) satisfies \(X^r-\beta\) over \(K\),
\[
 d\leq re_1\leq139\cdot137=19\,043,
\]
contrary to \(d\geq48\,688\).  Thus \(\beta\notin K^p\) for every
prime \(p\mid r\).

It remains to exclude Capelli's exceptional fourth-power case.  If
\(4\mid r\), then \(2\) occupies exactly one role in the primitive triple
\((A,B,N)\).  The local proofs of Corollary~\ref{cor:totalprimepower} and
Theorem~\ref{thm:endpointpower}, including the residual splitting of the
merged side, place every unit root, after an unramified base extension, in a
packet for which
\[
 v_2(\beta-\zeta)=\frac1\ell<1,
 \qquad [L(\beta):L]\leq\ell,
\]
with \(\zeta\) a root of unity of odd order and \(\ell\geq2\).  Since
\[
 -\beta-\zeta=-\bigl((\beta-\zeta)+2\zeta\bigr),
\]
the two summands have unequal valuations and
\(v_2(-\beta-\zeta)=1/\ell\).  Lemma~\ref{lem:edgepower} applied to
\(-\beta\) shows that \(-\beta\) is not a square in \(K\).  In particular
\(\beta\notin-4K^4\), because
\(\beta=-4\eta^4\) would give \(-\beta=(2\eta^2)^2\).
Capelli's criterion now proves the irreducibility of \(X^r-\beta\).
The proof for \(X^s-\gamma\) is identical.

It follows that \(d=re_1=se_2\).  Since \(\gcd(r,s)=1\), one has
\(s\mid e_1\) and \(r\mid e_2\), giving \((18.4)\).
\end{proof}

\subsection{One-package support and full-radical congruences}

\begin{proposition}[One-package unit support bound]
\label{prop:onepackagesupport}
Let \(f\in\Z[x]\) be a monic irreducible factor of a selected primitive
orientation of shape \((A,B)\), let \(e=\deg f\), and suppose
\(f(0)=\pm1\).  Then
\begin{equation}
 \operatorname{rad}(AB(A+B))
 \leq2eM(f)^{2-1/e}.
 \tag{18.5}
\end{equation}
\end{proposition}

\begin{proof}
For every prime dividing \(AB(A+B)\), the proof of
Proposition~\ref{prop:momentcongruence} uses only the one primitive package
containing that prime.  Since the leading and constant coefficients of \(f\)
are units, its endpoint nonunit block is absent.  The same all-moment
congruence, Hankel determinant, and Mahler discriminant estimate as in
Proposition~\ref{prop:hankeldisc} therefore apply with
\(\mathcal R=\operatorname{rad}(AB(A+B))\), giving \((18.5)\).
\end{proof}

\begin{corollary}[The exact support cap used by the computation]
\label{cor:unitsupportcap}
Under the hypotheses of Proposition~\ref{prop:unitkummer},
\begin{equation}
 r\operatorname{rad}(ABN)<4\,398\,937,
 \qquad
 s\operatorname{rad}(CEM)<4\,398\,937.
 \tag{18.6}
\end{equation}
\end{corollary}

\begin{proof}
Let \(f_\beta\) be the minimal polynomial of \(\beta\).  From \((18.4)\),
\[
 \log M(f_\beta)
 =e_1\,\mathrm h(\beta)
 =e_1r\,\mathrm h(\alpha)
 =d\,\mathrm h(\alpha)
 =\log M(h_\alpha).
\]
Since \(\beta\) is noncyclotomic, \(M(f_\beta)>1\).  By
Propositions~\ref{prop:fielddegree} and \ref{prop:onepackagesupport}, and
using \(e_1<4D^{2/3}/r\) together with \((17.13)\), we obtain
\[
 \begin{aligned}
 r\operatorname{rad}(ABN)
 &\leq2re_1M(f_\beta)^{2-1/e_1}\\
 &<8D^{2/3}
   \exp\!\left(\frac{16\log(2D)}{D^{1/3}}\right)
 =:C(D).
 \end{aligned}
\]
Indeed, \(\log M(f_\beta)=\log M(h_\alpha)\) and
\(\log M(h_\alpha)<8\log(2D)/D^{1/3}\).  Writing
\(x=D^{1/3}\),
\[
 \frac{\dd}{\dd x}\log C(x^3)
 =\frac{2x+48-16\log(2x^3)}{x^2}>0
\]
throughout \(D_0\leq D\leq D_1\): here \(x>189\) and
\(\log(2D)<20\), so the numerator is greater than \(106\).  Thus it is
enough to bound \(C(D_1)\).  The exact rational certificate in
\path{archive/Qab9_upper_bundle/code/certify_constants.py} uses
\[
 \frac{559\,764\,608}{10^6}<D_1^{1/3}
 <\frac{559\,764\,609}{10^6}
\]
and rational upper bounds from the atanh series for \(\log\) and the Taylor
series for \(\exp\).  It proves
\[
 C(D_1)<4\,398\,937.
\]
The second package is symmetric.
\end{proof}

\begin{proposition}[Full-radical degree signature in the unit case]
\label{prop:unitfullradical}
Under the hypotheses of Proposition~\ref{prop:onepackagesupport}, for every
prime \(p\mid AB(A+B)\),
\begin{equation}
 e\equiv0\quad\text{or}\quad-2\pmod p.
 \tag{18.7}
\end{equation}
Thus \(\operatorname{rad}(AB(A+B))\mid e(e+2)\).
\end{proposition}

\begin{proof}
For a total prime this is Theorem~\ref{thm:radicalsieve}.  At an endpoint
prime the unit hypothesis removes the endpoint nonunit block.  Borisov's
remaining root-of-unity clusters have contributions divisible by \(p\),
except that the center-\(1\) packet contributes either \(0\) or \(p-2\)
modulo odd \(p\).  For \(p=2\), every packet contribution is even, including
the residual-split merged side.  This proves \((18.7)\) in every role.
\end{proof}

\subsection{Role-binomial degree bounds}

\begin{lemma}[Reduction to a role binomial]
\label{lem:rolebinomial}
Let \(p\) divide exactly one member \(R\) of the primitive triple
\((A,B,N)\), and let \(h_\alpha\) be monic with constant coefficient
\(\pm1\).  If \(h_\alpha\) divides a selected orientation
\(Q_{A,B}(X^r)\) or \(Q_{B,A}(X^r)\), then
\begin{equation}
 \overline{h_\alpha}\mid X^{rR}-1
 \quad\text{in }\mathbb F_p[X],
 \qquad d\leq rR.
 \tag{18.8}
\end{equation}
If the same prime occupies role \(S\) in the second primitive triple, then
\[
 d\leq\gcd(rR,sS).
\]
\end{lemma}

\begin{proof}
For either orientation, direct reduction of the collision trinomial gives
\[
 H_{T,N}(X)\equiv cX^\nu(X^R-1)\pmod p
\]
for some \(c\in\mathbb F_p^\times\), \(\nu\geq0\), where
\(T\in\{A,B\}\) is the selected endpoint.  Hence
\[
 (X^r-1)^2Q_{T,N-T}(X^r)=H_{T,N}(X^r)
 \equiv cX^{r\nu}(X^{rR}-1)\pmod p.
\]
The monic unit polynomial \(h_\alpha\) has degree-preserving reduction and is
not divisible by \(X\), proving the first assertion.  If the prime appears in
both shapes, use
\(\gcd(X^a-1,X^b-1)=X^{\gcd(a,b)}-1\), valid in every characteristic.
\end{proof}

\subsection{The same-shape Bezout bound}

\begin{proposition}[Same primitive shape]
\label{prop:sameshapebezout}
Suppose distinct normalized packages have the same primitive shape
\(\{A,B\}\), with coprime scales \(r,s\), and share a noncyclotomic
irreducible factor of degree \(d\).  Then
\begin{equation}
 d\leq2rs.
 \tag{18.9}
\end{equation}
\end{proposition}

\begin{proof}
After replacing the common root by its inverse if necessary, put
\(x=\alpha^r\) and choose \(y=\alpha^s\) or \(\alpha^{-s}\) so that both
\(x,y\) are roots of \(Q_{A,B}\).  According as the two selected
orientations agree or disagree,
\[
 x^s=y^r
 \qquad\text{or}\qquad
 x^sy^r=1.
\]
Put \(N=A+B\) and
\[
 u=\frac A N x^B,
 \qquad 1-u=\frac B N x^{-A},
 \qquad
 v=\frac A N y^B,
 \qquad 1-v=\frac B N y^{-A}.
\]
Because \(\gcd(A,B)=1\), the map \(x\mapsto u\) is injective on the roots
under consideration: \(u\) determines both \(x^B\) and \(x^{-A}\), and an
integral Bezout combination recovers \(x\).

In the same-orientation case the point \((u,v)\) lies on
\[
 u^s=c_0v^r,
 \qquad
 (1-u)^s=c_1(1-v)^r;
\]
in the opposite-orientation case it lies on
\[
 u^sv^r=c_0,
 \qquad
 (1-u)^s(1-v)^r=c_1,
\]
with nonzero rational constants.  In either case these are two divisors of
bidegree \((s,r)\) on \(\mathbb P^1\times\mathbb P^1\), whose intersection
number is \(2rs\).

They have no common curve component.  Indeed, if one of \(u,v\) were
constant on a component, either binomial equation would make the other
constant as well, so the component would be a point.  On the normalization
of a genuine curve component, both \(u,v\) are therefore nonconstant, and
the two equations give two
linear relations among
\[
 \Delta_0=\operatorname{div}(u),\quad
 \Delta_1=\operatorname{div}(1-u),\quad
 E_0=\operatorname{div}(v),\quad
 E_1=\operatorname{div}(1-v).
\]
Lemma~\ref{lem:fourunit} forces \(v\) to be one of the six transforms of
\(u\).  In the basis \((\Delta_0,\Delta_1)\), their divisor pairs are
\[
\begin{array}{c|cc}
 v&E_0&E_1\\ \hline
 u&(1,0)&(0,1)\\
 1-u&(0,1)&(1,0)\\
 u^{-1}&(-1,0)&(-1,1)\\
 (1-u)^{-1}&(0,-1)&(1,-1)\\
 u/(u-1)&(1,-1)&(0,-1)\\
 (u-1)/u&(-1,1)&(-1,0)
\end{array}
\]
The same-orientation relations are \(s\Delta_i=rE_i\); the table permits
only \(v=u\) and \(r=s\), which would give the same package.  The
opposite-orientation relations are
\(s\Delta_i+rE_i=0\), and no row satisfies both.  Thus the intersection is
zero-dimensional.

The \(d\) conjugates of \(\alpha\) give \(d\) distinct intersection points.
Indeed equality of the corresponding \((u,v)\) recovers both
\(\alpha^r\) and \(\alpha^{\pm s}\), and \(\gcd(r,s)=1\) recovers
\(\alpha\).  Bezout's theorem now gives \((18.9)\).
\end{proof}

\subsection{The exact enumeration}

The proof pipeline deliberately avoids constructing any large polynomial.
Corollary~\ref{cor:unitsupportcap} first enumerates all primitive shapes
\(A<B\), \(N=A+B\leq D_1\), satisfying
\[
 \operatorname{rad}(A)\operatorname{rad}(B)\operatorname{rad}(N)
 \leq4\,398\,937.
\]
This enumeration is exhaustive without scanning every \(A\).  Order the
three radicals as \(\rho_1\leq\rho_2\leq\rho_3\).  Then
\[
 \rho_1\leq\lfloor4\,398\,937^{1/3}\rfloor=163,
 \qquad
 \rho_2\leq\lfloor4\,398\,937^{1/2}\rfloor=2097.
\]
The program generates every integer at most \(D_1\) with radical at most
these two bounds, assigns the two smallest radical roles in all three possible
ways, and reconstructs the remaining member from \(A+B=N\).  A canonical
role ordering removes duplicates.  A separate brute-force test checks this
logic on small boxes.

For every coprime scale pair \(1\leq r\leq s\leq139\), the pair program
then applies, in order:
\begin{enumerate}[label=(\arabic*)]
\item the support compatibility of Corollary~\ref{cor:supportcompatibility}
outside \(\{2,3\}\), allowing a missing prime only when it divides the
opposite scale;
\item the two support caps \((18.6)\), the range \(D_0\leq D\leq D_1\),
and disjointness of the two scaled additive triples;
\item the exact integer bounds \(d>D/140\), \(d^3<64D^2\),
Theorem~\ref{thm:schinzelbound}, and the proper-factor inequalities
\[
 sk\leq N-4,
 \qquad rk\leq M-4.
\]
Here global isolation of an irreducible primitive base makes each selected
primitive factor proper; Theorem~\ref{thm:rational} then rules out a
degree-one complementary factor.
\item the role-binomial bounds of Lemma~\ref{lem:rolebinomial};
\item the full-radical congruences
\[
 sk\equiv0,-2\pmod p\quad(p\mid ABN),
 \qquad
 rk\equiv0,-2\pmod p\quad(p\mid CEM),
\]
combined by an exact Chinese-remainder interval test.
\end{enumerate}
No irreducibility oracle, floating-point cutoff, finite-field factorization,
or modular gcd is used in this main pipeline.  The deterministic output is:
\[
\begin{array}{lr}
\text{radical-admissible primitive shapes}&237\,418\\
\text{scaled records}&2\,450\,291\\
\text{coprime scale pairs}&5\,952\\
\text{shape pairs after range, support cap, and disjointness}&67\,730\,966\\
\text{after the degree interval}&2\,124\,295\\
\text{after role-binomial bounds}&139\,640\\
\text{after full-radical CRT feasibility}&761\\
\text{same-shape rows among these}&761\\
\text{rows surviving Proposition~\ref{prop:sameshapebezout}}&0.
\end{array}
\]

\begin{theorem}[Computer-assisted elimination of the unit range]
\label{thm:upperelimination}
There is no normalized counterexample with
\[
 6\,816\,242\leq\max\{m,n,p,q\}\leq175\,394\,637.
\]
Consequently every normalized counterexample satisfies
\begin{equation}
 \boxed{\max\{m,n,p,q\}\leq6\,816\,241.}
 \tag{18.10}
\end{equation}
\end{theorem}

\begin{proof}
Every counterexample in the displayed range satisfies all of the
mathematical filters just listed, so the exhaustive exact program places it
among the \(761\) final rows.  In every row the two primitive shapes agree.
Proposition~\ref{prop:sameshapebezout} gives
\[
 d\leq2rs\leq2\cdot139\cdot138=38\,364,
\]
whereas Proposition~\ref{prop:uniformdegree} gives \(d\geq48\,688\).  Thus
all rows are impossible.  Combining this with
Theorem~\ref{thm:explicitbound} proves \((18.10)\).
\end{proof}

\begin{remark}[Certificate boundary]
The theorem is computer-assisted in the ordinary source-verification sense.
The bundle contains the complete shape list, the \(761\) pre-final rows, an
independent Python verifier for every recorded row, deterministic checksums,
and a brute-force small-box regression test for the shape enumerator.  The
coverage of the large pair scan is encoded transparently in the short C++
loop rather than in a proof-assistant certificate.  A future Lean-oriented
version should replace that last trust boundary by interval or branch
coverage certificates; no large polynomial arithmetic would be required.
\end{remark}

\subsection{Certified constants and reproduction commands}

The archived upper-range verifier is replayed from the current bundle by
\begin{lstlisting}[basicstyle=\ttfamily\footnotesize]
make q12-verify-upper
\end{lstlisting}
The historical Qab9 upper generator can also be rerun from
\path{archive/Qab9_upper_bundle} using its local Makefile.  That pipeline
requires a C++20 compiler with OpenMP and Python~3; its Python verification
uses only the standard library.  The exact rational certificate
\path{archive/Qab9_upper_bundle/code/certify_constants.py} prints
\begin{lstlisting}[basicstyle=\ttfamily\scriptsize]
plastic_power_coefficients=62608681,109870576,82938844
theta70_rational_lower=220094051941/625
twice_D_max=350789274
uniform_degree_bound=d>D/140
scale_max=139
same_shape_max=38364
unit_degree_integer_min=48688
support_endpoint_upper_floor=4398936
support_product_cap=4398937
\end{lstlisting}
The two analytic transition points in Theorem~\ref{thm:explicitbound} and
Corollary~\ref{cor:unitrange} are independently checked by the supplied Sage
and Arb script \path{archive/Qab9_upper_bundle/code/certify_thresholds.sage}.  All comparisons in
both scripts are interval or exact rational comparisons, not binary
floating-point guesses.

\subsection{Exact factorization through total sixty}

Exact factorization over \(\Q\) was rerun with SymPy 1.14 for every unordered
unequal pair with total \(a+b\leq60\).  There are
\[
 \sum_{N=3}^{60}\left\lfloor\frac{N-1}{2}\right\rfloor=870
\]
such pairs.  Every \(Q_{a,b}\) was irreducible, and no normalized reciprocal
factor orbit occurred in two different packages.  The calculation took exact
quotients and exact rational factorization; no numerical root matching was
used.

For every tested pair, its primitive reduction also has total at most \(60\)
and was found irreducible.  Theorem~\ref{thm:allliftsisolated} therefore
upgrades all \(870\) finite factorizations to global certificates: every one
of these packages is coprime to every distinct package in the entire infinite
family, including packages of arbitrarily large total.

The complete script is included below.
\begin{lstlisting}[language=Python,basicstyle=\ttfamily\scriptsize,breaklines=true,frame=single]
from math import gcd
import sympy as sp
from sympy import Poly, QQ

x = sp.symbols("x")
MAXN = 60
seen, reducible, collisions = {}, [], []
count = globally_certified = 0

def qpoly(a, b):
    g = gcd(a, b)
    num = Poly(a*x**(a+b) - (a+b)*x**a + b, x, domain=QQ)
    den = Poly((x**g - 1)**2, x, domain=QQ)
    return num.exquo(den)

def reciprocal_orbit_key(f):
    p = Poly(f, x, domain=QQ).monic()
    rev = Poly.from_list(list(reversed(p.all_coeffs())),
                         gens=x, domain=QQ).monic()
    return min(tuple(p.all_coeffs()), tuple(rev.all_coeffs()))

for N in range(3, MAXN + 1):
    for a in range(1, (N - 1)//2 + 1):
        b = N - a
        count += 1
        q = qpoly(a, b)
        _, factors = sp.factor_list(q.as_expr(), x, domain=QQ)
        is_irreducible = (len(factors) == 1 and factors[0][1] == 1
                and Poly(factors[0][0], x).degree() == q.degree())
        if not is_irreducible:
            reducible.append((a, b))
        g = gcd(a, b)
        base = qpoly(a//g, b//g)
        _, base_factors = sp.factor_list(base.as_expr(), x, domain=QQ)
        base_irreducible = (len(base_factors) == 1
                and base_factors[0][1] == 1
                and Poly(base_factors[0][0], x).degree() == base.degree())
        if base_irreducible:
            globally_certified += 1
        for f, exponent in factors:
            key = reciprocal_orbit_key(f)
            if key in seen and seen[key] != (a, b):
                collisions.append((seen[key], (a, b)))
            else:
                seen[key] = (a, b)

print(count, reducible, collisions, globally_certified)
# Output: 870 [] [] 870
\end{lstlisting}


\section{Lower-range extraction degrees and endpoint states}
\label{sec:lowerstates}

The audited upper computation gives
\[
 D:=\max\{m,n,p,q\}\leq6\,816\,241.
\]
This section repairs the exceptional fourth-power branch in the preceding
revision and records the lower-range identities in a form suitable for exact
enumeration.  Write the two normalized packages as
\[
 \{ra,rb\},\qquad \{sc,se\},
 \qquad \gcd(a,b)=\gcd(c,e)=1,
 \qquad \gcd(r,s)=1,
\]
and let \(\alpha\) be a common root of selected ordered orientations.  Put
\[
 E=\Q(\alpha),\quad K=\Q(\alpha^r),\quad L=\Q(\alpha^s),
\]
\[
 d=[E:\Q],\quad e_1=[K:\Q],\quad e_2=[L:\Q],
 \quad t=[E:K],\quad u=[E:L].
\]

\subsection{Extraction degrees and the reduced correspondence}

\begin{lemma}[Extraction-degree identities]\label{lem:extractiondegrees}
One has
\begin{equation}
 1\leq t\leq r,\qquad 1\leq u\leq s,
 \qquad d=te_1=ue_2,
 \qquad d\leq e_1e_2.
 \tag{19.1}
\end{equation}
Moreover \(E=KL\).
\end{lemma}

\begin{proof}
The element \(\alpha\) satisfies \(X^r-\alpha^r\) over \(K\), and similarly
on the second side.  The tower law gives the first identities.  An integral
B\'ezout combination of the coprime integers \(r,s\) recovers \(\alpha\)
from \(\alpha^r,\alpha^s\), so \(E=KL\).  Finally
\([KL:\Q]\leq[K:\Q][L:\Q]\).
\end{proof}

\begin{proposition}[General reduced correspondence]
\label{prop:reducedcorrespondence}
Suppose selected ordered primitive orientations \(Q_{a,b}(x^r)\) and
\(Q_{c,e}(x^s)\) share a noncyclotomic irreducible factor of degree \(d\).
Put
\[
 g_0=\gcd(rb,se),\quad B_0=\frac{rb}{g_0},\quad E_0=\frac{se}{g_0},
\]
\[
 g_1=\gcd(ra,sc),\quad A_0=\frac{ra}{g_1},\quad C_0=\frac{sc}{g_1}.
\]
Unless the normalized packages are identical,
\begin{equation}
 \boxed{d\leq E_0A_0+B_0C_0
 =\frac{rs(ae+bc)}{\gcd(rb,se)\gcd(ra,sc)}.}
 \tag{19.2}
\end{equation}
\end{proposition}

\begin{proof}
Set \(x=\alpha^r\), \(y=\alpha^s\), \(N=a+b\), \(M=c+e\), and
\[
 U_0=\frac aN x^b,\qquad 1-U_0=\frac bN x^{-a},
 \qquad
 V_0=\frac cM y^e,\qquad 1-V_0=\frac eM y^{-c}.
\]
The identity \(x^s=y^r\) gives two curves in
\(\mathbb P^1\times\mathbb P^1\):
\[
 U_0^{E_0}=\lambda_0V_0^{B_0},\qquad
 (1-U_0)^{C_0}=\lambda_1(1-V_0)^{A_0}.
\]
Their bidegrees are \((E_0,B_0)\) and \((C_0,A_0)\), hence their
intersection number is \(E_0A_0+B_0C_0\).  A common curve component would,
by Lemma~\ref{lem:fourunit}, make \(V_0\) one of the six M\"obius transforms
of \(U_0\).  The divisor table used in Proposition~\ref{prop:sameshapebezout}
then leaves only \(V_0=U_0\), \(se=rb\), and \(sc=ra\), which is the
identical-package case.

The value \(U_0\) determines both \(x^b\) and \(x^{-a}\), and therefore
\(x\), since \(\gcd(a,b)=1\); similarly \(V_0\) determines \(y\).  The
pair determines \(\alpha\) because \(\gcd(r,s)=1\).  Thus the \(d\)
conjugates give \(d\) distinct intersection points, and B\'ezout proves the
bound.
\end{proof}

For equal \emph{selected ordered} orientations, (19.2) gives \(d\leq2rs\).
When the unordered primitive shapes agree but the selected orientations are
opposite, Proposition~\ref{prop:sameshapebezout} still gives the stronger
bound \(d\leq2rs\); the two statements should not be conflated.

\subsection{Norm transfer, endpoint arithmetic, and packet occupancy}

Let
\[
 h_\alpha(X)=UX^d+\cdots+V\in\Z[X],\qquad U>0,
\]
be the primitive minimal polynomial of \(\alpha\), and let the primitive
minimal polynomials of \(\beta=\alpha^r\) and \(\gamma=\alpha^s\) have
leading and constant coefficients \((u_1,v_1)\) and \((u_2,v_2)\).

\begin{lemma}[Norm transfer]\label{lem:normtransfer}
One has
\begin{equation}
 u_1^t=U^r,\qquad |v_1|^t=|V|^r,
 \qquad
 u_2^u=U^s,\qquad |v_2|^u=|V|^s.
 \tag{19.3}
\end{equation}
\end{lemma}

\begin{proof}
Gauss's lemma applies because the relevant primitive factor divides a
primitive composition of \(Q_{a,b}\), whose leading and constant
coefficients divide the coprime integers \(a,b\).  Hence the leading and
constant coefficients of each primitive irreducible factor are coprime; in
particular \(V/U\) and \(v_1/u_1\) are reduced fractions.  The norm identity
\[
 N_{E/\Q}(\alpha^r)=N_{K/\Q}(\beta)^t
\]
then reads \((V/U)^r=(v_1/u_1)^t\), up to sign.  Comparing numerator and
denominator valuations gives the first pair of identities; the second is
identical.
\end{proof}

\begin{proposition}[Scaled endpoint divisibilities and slopes]
\label{prop:scaledendpoint}
For the first selected orientation,
\begin{equation}
 |V|^{ra}\mid b^d,\qquad U^{rb}\mid a^d;
 \tag{19.4}
\end{equation}
for the second,
\begin{equation}
 |V|^{sc}\mid e^d,\qquad U^{se}\mid c^d.
 \tag{19.5}
\end{equation}
If \(p\mid U\), then
\begin{equation}
 \frac{v_p(a)}{rb}=\frac{v_p(c)}{se};
 \tag{19.6}
\end{equation}
and if \(q\mid V\), then
\begin{equation}
 \frac{v_q(b)}{ra}=\frac{v_q(e)}{sc}.
 \tag{19.7}
\end{equation}
In particular \(U\mid a,c\) and \(|V|\mid b,e\).
\end{proposition}

\begin{proof}
Apply Proposition~\ref{prop:endpointarith} to the primitive factor of
\(Q_{a,b}\), raise its divisibilities to the extraction degree \(t\), and
use Lemma~\ref{lem:normtransfer} together with \(te_1=d\).  The slope
identities follow by choosing embeddings with negative, respectively
positive, valuation and reading the endpoint Newton slopes.  The final
assertion also follows directly from Gauss's lemma.
\end{proof}

\begin{proposition}[Exact endpoint packet occupancy]
\label{prop:packetoccupancy}
Suppose \(p^\kappa\Vert a\), \(p\mid U\), and put
\[
 \lambda=\frac r t v_p(U)\in\Z_{>0}.
\]
If the primitive factor of \(Q_{a,b}\) has degree \(e_1\), then its number
of roots in the negative-valuation endpoint block is
\begin{equation}
 w_p=\frac{b\lambda}{\kappa}\in\Z,
 \qquad 1\leq w_p\leq\min\{b,e_1\},
 \qquad e_1-w_p\equiv0\text{ or }-2\pmod p.
 \tag{19.8}
\end{equation}
There is a reciprocal statement for \(p^\kappa\Vert b\) and \(p\mid V\).
\end{proposition}

\begin{proof}
Every root in the nonunit endpoint block has valuation \(-\kappa/b\).
The factor leading coefficient has valuation
\(v_p(u_1)=rv_p(U)/t=\lambda\), while its constant coefficient is a
\(p\)-adic unit.  The product-of-roots formula therefore gives
\(w_p\kappa/b=\lambda\).  The remaining roots lie in the unit packets; the
center-
\(1\) packet contributes either zero or \(p-2\) roots and all other
clusters contribute multiples of \(p\), proving the congruence.
\end{proof}

\subsection{Sharper coefficient cutoffs}

\begin{proposition}[Two nonunit coefficients]\label{prop:twononunit}
If \(U>1\) and \(|V|>1\), then
\begin{equation}
 \boxed{D\leq16\,583.}
 \tag{19.9}
\end{equation}
\end{proposition}

\begin{proof}
Choose distinct primes \(p\mid U\), \(q\mid V\).  Equations
(19.6)--(19.7) and Proposition~\ref{prop:reducedcorrespondence} give
\[
 d\leq v_p(c)v_q(b)+v_p(a)v_q(e)
 \leq2\lfloor\log_2D\rfloor\lfloor\log_3D\rfloor.
\]
The plastic number \(\theta_0\) exceeds
\(1324717/10^6\), as follows by substituting this rational number into
\(X^3-X-1\).  The height estimate therefore gives the strict lower bound
\[
 d>
 \frac{D\log(1324717/10^6)}{2\log(2D)}.
\]
On each interval on which the two logarithmic floors are constant, the right
side is increasing.  Exact Arb interval comparisons at the finitely many
power-of-two and power-of-three boundaries show that the two bounds are
incompatible for every \(D\geq16\,584\).  The certificate is
\texttt{code/qab12/certify\_qab12\_constants.py}.
\end{proof}

\begin{corollary}[Coefficient trichotomy]\label{cor:coefftrichotomy}
For \(16\,584\leq D\leq6\,816\,241\), at most one of \(U,|V|\) is a
nonunit and
\[
 \max\{U,|V|\}\leq26.
\]
For \(D\geq50\,000\), the upper bound improves to \(12\).
\end{corollary}

\begin{proof}
The function \(8\log(2D)/D^{1/3}\) is decreasing in this range.  The Arb
certificate proves that its value at \(16\,584\) is below \(\log27\), and
its value at \(50\,000\) is below \(\log13\).  The absolute leading and
constant coefficients of a primitive polynomial are at most its Mahler
measure.  Proposition~\ref{prop:twononunit} excludes two nonunits.
\end{proof}

\subsection{The exceptional fourth-power branch}

\begin{lemma}[Capelli's exceptional branch]\label{lem:exceptionalfourth}
Let \(K=\Q(\beta)\), and suppose
\[
 \beta=-4\eta^4\qquad(\eta\in K).
\]
Then the exceptional branch cannot occur in a total-role packet or in a
\(2\)-adic unit endpoint packet.  If it occurs in the leading nonunit
endpoint block of \(Q_{a,b}\), where \(2^\kappa\Vert a\), then
\begin{equation}
 e_1\leq b,\qquad 2b\mid e_1\kappa.
 \tag{19.10}
\end{equation}
The reciprocal assertion holds for the constant endpoint.
\end{lemma}

\begin{proof}
The element \(-\beta=(2\eta^2)^2\) is a square in \(K\).  In every
\(2\)-adic unit packet used in the local power analysis, after an unramified
extension there is an odd-order root of unity \(\zeta\) and an integer
\(\ell>1\) such that
\[
 v_2(\beta-\zeta)=1/\ell<1.
\]
Since \(v_2(2\zeta)=1\), the ultrametric inequality gives
\[
 v_2(-\beta-\zeta)=v_2((\beta-\zeta)+2\zeta)=1/\ell.
\]
The Newton-edge power obstruction, applied to \(-\beta\), now contradicts
its being a square.  This excludes total packets and the unit part of an
endpoint packet.

Every conjugate of \(\beta\) is again of the form \(-4\eta_i^4\).  Hence a
surviving irreducible factor has all its roots in the nonunit endpoint block,
which has length \(b\); thus \(e_1\leq b\).  Finally
\[
 N_{K/\Q}(\beta)=(-4)^{e_1}N_{K/\Q}(\eta)^4
\]
has even \(2\)-adic valuation, while the endpoint slope gives
\(|v_2N_{K/\Q}(\beta)|=e_1\kappa/b\).  Therefore
\(2b\mid e_1\kappa\).
\end{proof}

\begin{proposition}[Localization of a Kummer defect]
\label{prop:defectlocalization}
Let \(p^A\Vert r\).  If \(v_p(t)<A\), then either
\(\beta\in K^p\), or \(p=2\), \(A\geq2\), and
\(\beta\in-4K^4\).  The following alternatives exhaust the branch:
\begin{enumerate}[label=(\alph*)]
\item if \(p\) is good for the primitive shape, then \(e_1\leq p-2\);
\item a total-role prime is impossible;
\item a locally \(p\)-adic unit endpoint factor is impossible;
\item in a nonunit endpoint block,
\[
 e_1\leq b,\qquad pb\mid e_1\kappa
\]
for a leading endpoint \(p^\kappa\Vert a\), with the reciprocal statement
at the constant endpoint.  Necessarily \(p\in\{2,3,5,7\}\).
\end{enumerate}
In particular every unit-coefficient Kummer-defect counterexample satisfies
\begin{equation}
 D<140\cdot139\cdot137=2\,666\,020,
 \qquad d\leq139\cdot137=19\,043.
 \tag{19.11}
\end{equation}
\end{proposition}

\begin{proof}
Capelli's criterion gives the first dichotomy.  In the ordinary internal
\(p\)-th-power branch, Theorem~\ref{thm:frobeniusfactor},
Corollary~\ref{cor:totalprimepower}, and Theorem~\ref{thm:endpointpower}
give (a)--(d).  Lemma~\ref{lem:exceptionalfourth} supplies exactly the
missing exceptional \(p=2\) case.  The endpoint condition implies
\(\kappa\geq p\); since \(p^\kappa\leq6\,816\,241\), only
\(2,3,5,7\) remain.  For a unit factor every deficient prime is therefore
good, so \(e_1\leq p-2\leq137\), and (19.11) follows from
\(d=te_1\leq re_1\), \(r\leq139\), and \(d>D/140\).
\end{proof}

\subsection{A strict endpoint-state object}

\begin{definition}[Admissible lower state]\label{def:endpointstate}
A package state records
\[
 (r,t,e_1,d,U,V;a,b,\varepsilon)
\]
together with every deficient scale prime, its valuations in \(r,t\), its
role, and its endpoint exponent \(\kappa\).  A pair of package states is
\emph{admissible} if it has common \((d,U,V)\), satisfies
\begin{gather*}
 d=te_1=ue_2,\qquad d\leq e_1e_2,
 \qquad 140d>D,
 \qquad d^3<64D^2,
 \qquad d\leq\min\{a+b,c+e\}-2,\tag{19.12}
\end{gather*}
all norm-transfer, endpoint, slope, packet-occupancy, radical-signature, and
correspondence conditions above, the coefficient trichotomy of
Corollary~\ref{cor:coefftrichotomy}, the alternatives of
Proposition~\ref{prop:defectlocalization}, coprime scales, and disjoint scaled
additive triples.
\end{definition}

In the current Qab12 bundle these fields are represented by typed CSV rows
under \path{data/qab12}.  The separate verifiers in \path{code/qab12}
recompute the global integer conditions from those rows, and the small-box
generators supply coverage regressions for the enumeration logic.

\begin{theorem}[Finite lower-state envelope]\label{thm:lowerstatereduction}
Every normalized counterexample with \(D\leq6\,816\,241\) determines an
admissible lower state.  The state set is finite and effectively enumerable.
\end{theorem}

\begin{proof}
The preceding lemmas and the exact local factor signatures give all listed
necessary conditions.  The explicit bounds on \(D,r,s,d,U,V\), together
with the endpoint divisibilities, leave only finitely many integer tuples.
\end{proof}

\section{Exact elimination of the complete unit-coefficient branch}
\label{sec:unitlowercomputation}

After the upper range is gone, a normalized counterexample in the lower box has either unit endpoint coefficients or at least one nonunit endpoint coefficient.  This section completes the unit branch by combining full-Kummer searches, modular gcd certificates, and modular factor-degree certificates.

This section records a new production computation.  It covers the whole
remaining box under the hypothesis
\[
 U=|V|=1,
\]
including both full-Kummer and Kummer-defect cases.  All branch decisions in
the C++ enumerators are integer decisions.  The one analytic post-filter uses
Arb balls and rounds in the conservative direction.

\begin{lemma}[Subset-degree irreducibility certificate]
\label{lem:subsetdegreecertificate}
Let \(F\in\Z[x]\) be primitive of degree \(n\), and let
\(p\nmid\operatorname{lc}(F)\).  If the reduction of \(F\) modulo \(p\)
has irreducible factor degrees \(d_{p,j}\), counted with multiplicity, then
the degree of every factor of \(F\) over \(\Q\) is a subset sum of the
\(d_{p,j}\).  Consequently, if the intersections of these subset-sum sets
for finitely many primes is \(\{0,n\}\), then \(F\) is irreducible.
\end{lemma}

\begin{proof}
A primitive rational factor may be chosen in \(\Z[x]\).  Since the total
leading coefficient is nonzero modulo \(p\), neither factor loses degree on
reduction.  Its reduction is a divisor of \(F\bmod p\), so its degree is the
sum of a submultiset of the modular irreducible factor degrees.
\end{proof}

\begin{theorem}[Unit-coefficient elimination]\label{thm:unitelimination}
There is no normalized counterexample whose primitive minimal polynomial has
leading and constant coefficients of absolute value one.
\end{theorem}

\begin{proof}[Computer-assisted proof]
The calculation has three exhaustive branches.  In each branch the search is
conservative: a shape or pair is retained whenever a local theorem is not
being used as a certified filter, and it is discarded only by an independently
recomputed arithmetic certificate.

\smallskip
\noindent\textbf{(I) Full Kummer degree and \(50\,000\leq D\leq6\,816\,241\).}
The one-package support estimate is bounded throughout this interval by
\[
 r\,\operatorname{rad}(ab(a+b))<1\,612\,000.
\]
The generator finds \(59\,063\) reducible-candidate primitive shapes.  The
scale-pair scan considers \(101\,752\,371\) indexed shape pairs and applies
only the exact radical, height, role-binomial, modular factor-degree, and
correspondence filters.  The conservative terminal list consists of seven
package pairs.  For each of the four selected orientation pairs attached to
each package pair, the certificate records a good auxiliary prime at which the
two collision trinomials have modular gcd of degree exactly two.  The
independent verifier reconstructs all twenty-eight modular computations.

\smallskip
\noindent\textbf{(II) Full Kummer degree and \(D<50\,000\).}
Here \(d\leq5428\).  For every primitive factor degree \(e\), the relation
\(\operatorname{rad}(ab(a+b))\mid e(e+2)\) permits direct smooth-number
enumeration.  It produces \(9491\) shape--degree records.  The exact pair scan
reduces \(7\,366\,972\) raw pairs to three terminal package pairs.  Their
twelve selected orientation pairs are certified by the same modular-gcd
criterion, and the verifier recomputes every claimed gcd degree.

\smallskip
\noindent\textbf{(III) A deficient scale prime.}
Proposition~\ref{prop:defectlocalization} gives
\(D<2\,666\,020\), \(d\leq19\,043\), and a deficient-side primitive degree
at most \(137\).  Enumerating extraction cores first and then the smooth
additive triples reduces the branch to \(21\,114\) integer state pairs,
representing \(9297\) package-orientation pairs but only \(46\) primitive
shapes.

Lemma~\ref{lem:subsetdegreecertificate} certifies every one of these
\(46\) shapes as irreducible.  The largest shape \((1,6655)\) is verified
with several auxiliary primes; the certificate file stores the exact modular
factor-degree lists and the separate verifier recomputes \(106\) modular
factorizations.  The global isolation theorem for irreducible primitive bases
then excludes every row in this branch.

The branch generators, output tables, certificates, and independent
recomputers are in \path{code/qab12} and \path{data/qab12}.  The bundled
checksums and branch-specific verifier outputs provide the row counts and
SHA-256 bindings used for replay.
\end{proof}

\begin{table}[ht]
\centering
\begin{tabular}{lrrr}
\hline
branch & principal input & preterminal rows & terminal objects\\
\hline
full Kummer, \(D\geq50\,000\) & \(101\,752\,371\) pairs & \(7\) pairs & \(28\) gcd certificates\\
full Kummer, \(D<50\,000\) & \(7\,366\,972\) pairs & \(3\) pairs & \(12\) gcd certificates\\
Kummer defect & state-pair scan & \(46\) shapes & \(46\) irreducibility certificates\\
\hline
\end{tabular}
\caption{The three unit-coefficient branches.  Counts are recomputed from the
bundled deterministic artifacts; timing data are kept outside the proof
manifest.}
\end{table}

\begin{corollary}[Remaining nonunit envelope]\label{cor:remainingnonunit}
Every normalized counterexample has \(U>1\) or \(|V|>1\).  After replacing
\(\alpha\) by \(\alpha^{-1}\) and reversing both selected orientations when
necessary, the one-nonunit branch may be written with \(U>1\), \(|V|=1\).
Moreover:
\begin{enumerate}[label=(\roman*)]
\item if both coefficients are nonunits, then \(D\leq16\,583\);
\item if \(D\geq16\,584\), then \(2\leq U\leq26\);
\item if \(D\geq50\,000\), then \(2\leq U\leq12\);
\item every deficient endpoint prime lies in \(\{2,3,5,7\}\) and satisfies
(19.8), (19.10), and the slope identities.
\end{enumerate}
\end{corollary}


\section{Exact nonunit branch artifacts}
\label{sec:nonunitcomputation}

The final branch keeps the endpoint coefficient data instead of collapsing to full Kummer degree.  The state records here are designed so that the verifiers can recompute norm transfer, endpoint occupancy, slope signatures, and terminal modular gcd certificates from explicit rows.

This section records the current nonunit computation.  It is deliberately
separated from the unit proof because the nonunit package states retain exact
extraction degrees and endpoint packet occupancies rather than full Kummer
multiplicity.

\begin{theorem}[Both-nonunit elimination]
\label{thm:twononunitelimination}
There is no normalized counterexample for which both the leading coefficient
and the absolute constant coefficient of the primitive minimal polynomial are
nonunits.
\end{theorem}

\begin{proof}[Computer-assisted proof]
Proposition~\ref{prop:twononunit} and the strengthened constants of
Corollary~\ref{cor:remainingnonunit} reduce this branch to
\(D\leq16\,583\), \(d\leq224\), and endpoint valuations at most \(14\).
The generator enumerates every one-package state satisfying the exact norm
transfer, endpoint divisibility, slope-signature, packet-occupancy,
deficient-prime, total-radical, degree, and size conditions.  It finds
\(1054\) package states.  The independent verifier recomputes every row and
checks that the common-factor keys
\[
 (U,V,d;\sigma_U,\sigma_V)
\]
are all distinct.  Hence two packages cannot share the same factor data; the
state-pair file is empty.  The verification output is
\path{data/qab12/two_nonunit_verification_output.txt}.
\end{proof}

\begin{theorem}[One-nonunit elimination above \(16\,583\)]
\label{thm:onenonunitlargeelimination}
There is no normalized counterexample with exactly one nonunit endpoint
coefficient and \(D\geq16\,584\).
\end{theorem}

\begin{proof}[Computer-assisted proof]
After reciprocal normalization write \(U>1\), \(|V|=1\).  The exact
coefficient-transition certificate gives \(2\leq U\leq26\).  For each such
\(U\), the bound
\[
 \log U<\frac{8\log(2D)}{D^{1/3}}
\]
gives the sharp coefficient-specific cap shown in the bundled table
\path{data/qab12/qab12_constants_output.txt}.  The largest cap is
\(6\,816\,241\) for \(U=2\), and the cap for \(U=26\) is \(16\,743\).
The same certificate supplies an integer \(C(U)\) such that \(D<C(U)d\)
throughout the corresponding coefficient range.

The all-shapes generator is conservative: it keeps primitive shapes even
when a known irreducibility theorem would already discard them.  Across
\(U=2,\ldots,26\) it produces \(608\,490\) package records, grouped into
\(514\,280\) common-factor keys.  The pair scan tests \(137\,530\) raw
same-key package pairs.  The successive exact filters leave
\[
\begin{array}{rcl}
10\,636 && \text{after coprime scales},\\
 7\,916 && \text{after range and extraction-degree checks},\\
 5\,430 && \text{after disjointness},\\
   483 && \text{after support compatibility},\\
    76 && \text{after the global degree bounds},\\
     0 && \text{after the local role-binomial bound}.
\end{array}
\]
Thus no terminal modular gcd computation is needed in this branch.  The
script \path{code/qab12/verify_one_nonunit.py} independently recomputes the
row conditions, confirms the per-\(U\) pair counts, and verifies that every
survivor file is header-only.
\end{proof}

\begin{corollary}[Residual envelope before the final small computation]
\label{cor:q12residualenvelope}
If a normalized counterexample remains after the computations in this
bundle, then after replacing \(\alpha\) by \(\alpha^{-1}\) and reversing both
selected orientations if necessary, its primitive minimal polynomial has
\[
 U>1,
 \qquad |V|=1,
 \qquad D\leq16\,583.
\]
Moreover \(U\mid a\), hence \(2\leq U\leq D\), every endpoint valuation is at
most \(14\), the safe degree bound
\[
 d^3<64D^2
\]
gives \(d\leq2601\), and every deficient endpoint prime that lies in the
nonunit endpoint block lies in \(\{2,3,5,7\}\).
\end{corollary}

\begin{theorem}[Residual one-nonunit elimination]
\label{thm:onenonunitresidualelimination}
There is no normalized counterexample with exactly one nonunit endpoint
coefficient and \(D\leq16\,583\).
\end{theorem}

\begin{proof}[Computer-assisted proof]
By Corollary~\ref{cor:q12residualenvelope}, after reciprocal normalization it
is enough to enumerate
\[
 D\leq16\,583,\qquad |V|=1,\qquad 2\leq U\leq D,\qquad d\leq2601.
\]
The residual generator uses five slope-signature slots, enough because an
integer at most \(16\,583\) has at most five distinct prime factors.  It
enumerates every one-package state satisfying the norm-transfer exponent
integrality, endpoint divisibility, packet occupancy, slope congruence,
deficient-prime, total-radical, extraction-degree, and size conditions.  It
is conservative and keeps rows that later degree filters can remove.  The
output is
\[
\begin{array}{rcl}
233\,278 && \text{package records},\\
212\,369 && \text{common-factor keys},\\
 29\,611 && \text{raw same-key state pairs},\\
  5\,680 && \text{after coprime scales and the residual range},\\
  5\,679 && \text{after extraction-degree checks},\\
  2\,345 && \text{after disjointness},\\
    428 && \text{after support compatibility},\\
    234 && \text{after the global degree bounds},\\
      4 && \text{after the local role-binomial bound},\\
      4 && \text{after the reduced correspondence bound}.
\end{array}
\]
The four terminal state pairs collapse to two package pairs:
\[
\begin{array}{ccl}
U=8,  & d=99:  &(3;440,1,441)\ \text{vs.}\ (4;3024,1,3025),\\
U=32, & d=123,205,285: &(1;6560,1,6561)\ \text{vs.}\ (2;1024,1,1025).
\end{array}
\]
For each of these two package pairs and each of the four selected
orientations, the certificate file
\path{data/qab12/one_nonunit_residual_modular_certificates.csv} records the
prime \(p=1009\) and modular gcd degree \(2\).  The verifier
\path{code/qab12/verify_one_nonunit_residual.py} recomputes every package
row, reruns the pair sieve, compares the terminal CSV files exactly, and
recomputes all eight modular gcds over \(\mathbf F_{1009}\).  Degree \(2\) is
accounted for by the forced double root at \(x=1\), so no noncyclotomic
common factor remains in any terminal orientation.
\end{proof}

\subsection*{Assurance status of the Qab12 bundle}

The bundled Qab12 artifacts now eliminate the audited upper range, the full
unit branch, the both-nonunit branch, the one-nonunit branch with
\(D\geq16\,584\), and the corrected residual one-nonunit branch with
\(D\leq16\,583\).  The light verifier replays the nonunit row/count checks,
the residual pair sieve, the residual modular gcd certificates, and the small
brute-force oracle for the residual generator.

\section*{Conclusion}

The mathematical reduction and the audited upper computation first give
\[
 D\leq6\,816\,241.
\]
The exceptional \(-4K^4\) branch is handled separately, so deficient scale
primes are localized without citing an ordinary internal-power theorem outside
its hypotheses.  The exact unit computation eliminates every unit-coefficient
common factor.  The new nonunit artifacts eliminate both endpoint coefficients
being nonunit and also eliminate the whole one-nonunit branch above
\(16\,583\).

The corrected residual one-nonunit computation then eliminates the remaining
range \(D\leq16\,583\), \(2\leq U\leq D\), \(d\leq2601\), using four terminal
state pairs and eight modular gcd certificates.

Consequently no normalized reciprocal-package counterexample remains in the
finite box reached by the mathematical reductions.  By the collision
dictionary and primitive normalization of Section~\ref{sec:dictionary}, this
proves Theorem~\ref{thm:main}.  The present bundle is therefore a reviewable
computer-assisted proof package for the reciprocal-package coprimality
statement.

\section*{Use of AI tools}

This proof was found through multiple rounds of prompting to ChatGPT 5.5 Pro.
The manuscript was edited and revised with the Codex CLI tool running
gpt-5.5 xhigh; that tool also finalized the computer verification part.  A
fuller account of the AI-assisted discovery, editing, and verification
process will be added in a later revision.

\begin{thebibliography}{99}

\bibitem{Borisov1998}
A.~Borisov,
\emph{On some polynomials allegedly related to the abc conjecture},
Acta Arith. \textbf{84} (1998), no.~2, 109--128.
\url{https://matwbn.icm.edu.pl/ksiazki/aa/aa84/aa8422.pdf}


\bibitem{AmorosoSombraZannier2017}
F.~Amoroso, M.~Sombra, and U.~Zannier,
\emph{Unlikely intersections and multiple roots of sparse polynomials},
Math. Z. \textbf{287} (2017), 1065--1081.
\url{https://arxiv.org/abs/1412.8059}

\bibitem{BeukersSchlickewei1996}
F.~Beukers and H.~P.~Schlickewei,
\emph{The equation \(x+y=1\) in finitely generated groups},
Acta Arith. \textbf{78} (1996), 189--199.
\url{https://eudml.org/doc/206941}

\bibitem{AndersonSaffVarga1979}
N.~Anderson, E.~B.~Saff, and R.~S.~Varga,
\emph{On the Enestr\"om--Kakeya theorem and its sharpness},
Linear Algebra Appl. \textbf{28} (1979), 5--16.

\bibitem{Cassels1997}
J.~W.~S.~Cassels,
\emph{An Introduction to the Geometry of Numbers},
Classics in Mathematics, Springer, 1997.

\bibitem{CoxLittleSchenck2011}
D.~A.~Cox, J.~B.~Little, and H.~K.~Schenck,
\emph{Toric Varieties},
Graduate Studies in Mathematics 124, American Mathematical Society, 2011.

\bibitem{Fulton1998}
W.~Fulton,
\emph{Intersection Theory}, second edition,
Ergebnisse der Mathematik und ihrer Grenzgebiete 2,
Springer, 1998.

\bibitem{SerreLocalFields}
J.-P.~Serre,
\emph{Local Fields},
Graduate Texts in Mathematics 67, Springer, 1979.

\bibitem{GuardiaMontesNart2008}
J.~Gu\`ardia, J.~Montes, and E.~Nart,
\emph{Newton polygons of higher order in algebraic number theory},
Trans. Amer. Math. Soc. \textbf{364} (2012), 361--416.
\url{https://arxiv.org/abs/0807.2620}

\bibitem{LangAlgebra}
S.~Lang,
\emph{Algebra}, third edition,
Graduate Texts in Mathematics 211, Springer, 2002.

\bibitem{Mahler1964}
K.~Mahler,
\emph{An inequality for the discriminant of a polynomial},
Michigan Math. J. \textbf{11} (1964), 257--262.

\bibitem{Smyth1971}
C.~J.~Smyth,
\emph{On the product of the conjugates outside the unit circle of an
algebraic integer},
Bull. London Math. Soc. \textbf{3} (1971), 169--175.

\end{thebibliography}

\end{document}
